% Single-file version of volume_obsolete/main.tex (2026-10-04): every
% \input is inlined below and the bibliography is the bibtex (alpha) output
% of ref.bib, inlined as \bibitem entries at the end.
% Build: pdflatex three times; no bibtex, no external dependency.
\documentclass[11pt]{article}
\usepackage{amsmath, amssymb, amsthm}
\usepackage{mathtools}

\usepackage{enumitem}
\usepackage[normalem]{ulem}

\usepackage{xcolor}
\usepackage{float}
\definecolor{ForestGreen}{rgb}{0.1333,0.5451,0.1333}
\usepackage[colorlinks=true,linkcolor=ForestGreen,citecolor=blue]{hyperref}

\usepackage[margin=1in]{geometry}
\usepackage{bbm}
\DeclareMathOperator*{\argmin}{argmin}
\usepackage[noabbrev,capitalize,nameinlink]{cleveref}
\crefname{equation}{}{}
\numberwithin{equation}{section}
\renewcommand{\theHequation}{\thesection.\arabic{equation}}
\usepackage{cite}
\usepackage{fancybox}
\newenvironment{fminipage}%
  {\begin{Sbox}\begin{minipage}}%
  {\end{minipage}\end{Sbox}\fbox{\TheSbox}}

\newenvironment{algbox}[0]{\vskip 0.2in
\noindent
\begin{fminipage}{6.3in}
}{
\end{fminipage}
\vskip 0.2in
}

\newtheorem{theorem}{Theorem}
\newtheorem{lemma}{Lemma}[section]

\usepackage{aliascnt}
\newaliascnt{corollary}{lemma}
\newtheorem{corollary}[corollary]{Corollary}
\aliascntresetthe{corollary}
\crefname{corollary}{Corollary}{Corollaries}
\newaliascnt{fact}{lemma}
\newtheorem{fact}[fact]{Fact}
\aliascntresetthe{fact}
\crefname{fact}{Fact}{Facts}

\theoremstyle{definition}
\newaliascnt{definition}{lemma}
\newtheorem{definition}[definition]{Definition}
\aliascntresetthe{definition}
\crefname{definition}{Definition}{Definitions}
\newaliascnt{proposition}{lemma}
\newtheorem{proposition}[proposition]{Proposition}
\aliascntresetthe{proposition}
\crefname{proposition}{Proposition}{Propositions}
\newaliascnt{conjecture}{lemma}
\newtheorem{conjecture}[conjecture]{Conjecture}
\aliascntresetthe{conjecture}
\crefname{conjecture}{Conjecture}{Conjectures}

\usepackage{etoolbox}
\newcommand{\thestmt}{}
\AtBeginEnvironment{theorem}{\renewcommand{\thestmt}{\thetheorem}}
\AtBeginEnvironment{lemma}{\renewcommand{\thestmt}{\thelemma}}
\AtBeginEnvironment{corollary}{\renewcommand{\thestmt}{\thecorollary}}
\AtBeginEnvironment{fact}{\renewcommand{\thestmt}{\thefact}}
\AtBeginEnvironment{definition}{\renewcommand{\thestmt}{\thedefinition}}
\AtBeginEnvironment{proposition}{\renewcommand{\thestmt}{\theproposition}}
\AtBeginEnvironment{conjecture}{\renewcommand{\thestmt}{\theconjecture}}
\setlist[enumerate,1]{ref=\thestmt(\arabic*)}
\crefname{enumi}{Item}{Items}
\Crefname{enumi}{Item}{Items}

\newcommand{\g}{\nabla}

\newcommand{\eps}{\varepsilon}
\newcommand{\R}{\mathbb{R}}
\newcommand{\Vol}{\mathrm{Vol}}
\renewcommand{\O}{\widetilde{O}}
\renewcommand{\l}{\langle}
\renewcommand{\r}{\rangle}

\newcommand{\E}[2]{\mathbb{E}_{#1}\left[#2\right]}
\renewcommand{\P}[2]{\mathbb{P}_{#1}\left[#2\right]}
\newcommand{\V}[2]{\mathsf{Var}_{#1}\left[#2\right]}
\newcommand{\C}{\mathsf{Cov}}
\newcommand{\Tr}[1]{\mathsf{Tr}\left(#1\right)}
\newcommand{\TV}{\mathrm{TV}}

\newcommand{\abs}[1]{\left|#1\right|}
\newcommand{\norm}[1]{\left\lVert#1\right\rVert}

\newcommand{\pe}{\preceq}
\newcommand{\se}{\succeq}

\newcommand{\poly}{\mathrm{poly}}
\newcommand{\op}{\mathrm{op}}

\newcommand{\phat}{\widehat{p}}

\renewcommand{\tilde}{\widetilde}
\newcommand{\Query}{\textsc{Query}}
\newcommand{\algname}[1]{\textnormal{\textsc{#1}}}
\newcommand{\ind}{\mathbbm{1}}

\newcommand{\todo}[1]{\text{\color{orange}\textbf{TODO: #1}}}
\newcommand{\as}[1]{\text{\color{green}\textbf{Alicia: #1}}}
\newcommand{\gap}[1]{\text{\color{red}\textbf{GAP: #1}}}
\newcommand{\tododevin}[1]{\text{\color{green!60!black}\textbf{Devin: #1}}}
\newcommand{\defeq}{\stackrel{\Delta}{=}}
\emergencystretch=2em

\begin{document}

\title{Approaching Cubic Oracle Calls for Isotropic Rounding}
\author{}
\date{October 4, 2026}


\clearpage\maketitle

\begin{abstract}
The main results are:
\begin{enumerate}
\item For arbitrary convex bodies, an improvement of the previous best query
  bound from $\tilde{O}(n^{3.5})$ to $\tilde{O}(n^{3.25})$,
  while also giving a simpler cooling path.
  This is obtained by directly prompting GPT 5.6 Sol
  and in Kimi Code with K3max (but not through prompt),
  and is in \cref{sec:asymmetric}.
  Its sharpness in \cref{sec:sharp} is obtained
  by directly prompting GPT 5.6 Sol.
\item For symmetric convex bodies, a matrix-regularized Gaussian cooling
  algorithm improves the rounding cost to $\tilde{O}(n^3)$, and hence gives
  volume estimation in $\tilde{O}(n^3\epsilon^{-2})$ membership queries
  (\cref{sec:symmetric}).
  This is obtained by prompting GPT 5.6 Sol to improve an algorithm from a
  previous submission with the above improvement in asymmetric settings.
  We have been unable to reproduce it by directly prompting / running Codex,
  including running a GPT 6 Astra agent swarm up to a budget of \$100.
\item The counter examples in \cref{sec:counter_coordinate_moment}
and \cref{sec:counter_ellipsoid_growth_asym} were obtained by directly
prompting GPT 5.6 Sol with the statements of those conjectures.
\end{enumerate}

\end{abstract}

\noindent\textbf{Note.}
This document contains intermediate results of a project that started
around February 2025
involving first Yang P. Liu and Colin Tang,
then September 2025 Alicia Stepin, and July 2026 Richard Peng.
These results were obtained using a combination of GPT 5.6 Sol,
FizzyProver (a heavy harness written by Junzhao Yang) running GPT 5.6 Sol,
and Kimi K3 / Kimi Code.

\newpage

\tableofcontents

\newpage

\section{Introduction}
\label{sec:intro}

Volume estimation of a convex set is a fundamental problem in computer science
with a long history
\cite{DFK91:journal,LS93:journal,KLS97:journal,LV06c:journal,LV06b:journal,LV07:journal,CV18:journal}.
There is a standard model for accessing convex bodies which most of these
previous works use.
For a convex body $K$, we assume that there is a routine
$\Query(x)$ which for a point
$x \in \R^n$ answers whether
$x \in K$.
We call $\Query$ a \emph{membership oracle}, and we measure the complexity of
our algorithm in the number of times we call this membership oracle.
All other running time overheads can be made to be polynomial in $n$, though we
do not carefully track it in this article.
Also, it is standard to assume that $B_n(0, r) \subseteq K \subseteq B_n(0, R)$
where
$B_n(x, r)$ is the
$n$-dimensional ball centered at
$x$ of radius
$r$.
Generally, the complexities of our algorithms (and previous algorithms) depend
logarithmically on $R/r$.
Some assumption like this is necessary to design nontrivial algorithms:
otherwise, $K$ could be a very small convex set inside a ball of radius
$R$, and it would take exponential time to even locate a single point inside
$K$.

The work of \cite{DFK91:journal} proved that the volume can be estimated to
within a constant factor using polynomially many membership oracle queries.
More recent work has focused on obtaining faster fine-grained runtimes for
volume estimation.
Generally, volume estimation algorithms can be broken down into two main phases,
with most works focusing on one phase or another.
The first phase is \emph{isotropic rounding}.
In this phase, the goal of the algorithm is to compute a linear transformation
$T$ on
$\R^n$ such that the covariance matrix of the body
$TK$ is approximately the identity matrix, i.e., to make the convex body nearly
isotropic.
The second phase is to design algorithms to estimate the volume of a
nearly-isotropic convex body.

The complexity of the second phase, i.e., estimating the volume of a
nearly-isotropic convex body, is reasonably well-understood by now.
In this case, the \emph{Gaussian cooling} algorithm of Cousins-Vempala
\cite{CV18:journal} gives a $(1+\eps)$-multiplicative approximation to the
volume in
$\O(n^3\eps^{-2})$ membership queries.
The complexity of the first isotropic rounding step has historically been
larger.
Before this work, the best general bound was due to \cite{JLLV26:journal}, using
$\O(n^{3.5} \log(nR/r))$ membership queries.
In the remainder of the introduction, we assume that $R/r \le \poly(n)$ for
simplicity, and use
$\O(\cdot)$ to suppress
$\poly(\log n)$ factors throughout the paper.

The main algorithmic results of this paper are two near-cubic rounding
algorithms.
The first works for arbitrary convex bodies and estimates the uncentered second
moment in $\O(n^{3.25})$ membership queries.
The second uses symmetry more strongly and obtains a cubic rounding bound for
symmetric convex bodies.

The next theorem concerns the general-body rounding.

\begin{theorem}
\label{thm:general-upper-bound}
The algorithm \algname{AsymmetricSecondMoment}
(\cref{fig:asymmetric-second-moment}), given membership query access to a convex
set $K \subseteq \R^n$ satisfying
$B_n(x_0, r) \subseteq K \subseteq B_n(x_0, R)$ and a failure parameter
$\beta \in (0,1)$, computes a positive semidefinite matrix
$\widehat M$ satisfying
\[
0.99\widehat M
\preceq \E{x \sim K}{xx^\top}
\preceq 1.01\widehat M.
\]
It succeeds with probability at least $1-\beta$ and makes
\[
O\left(n^{13/4}\log\left(2+\frac{nR}{r\beta}\right)^{O(1)}\right)
\]
membership query calls.
On the same success event, the algorithm also returns an invertible affine map
$\mathcal T$ such that the uniform law on
$\mathcal T(K)$ has covariance in
$[1/1.01,1.01]I$ and mean of norm at most
$1/250$; in particular
$B(0,0.99)\subseteq\mathcal T(K)$ and
$\E{x\sim\mathcal T(K)}{\|x\|^2}\leq1.02n$.
\end{theorem}

The next theorem concerns the cubic symmetric rounding.

\begin{theorem}
\label{thm:symmetric-upper-bound}
The algorithm \algname{Symmetric-Covariance} (\cref{fig:symmetric-covariance}),
given membership query access to a \emph{symmetric} convex set
$K \subseteq \R^n$ satisfying
$B_n(0, r) \subseteq K \subseteq B_n(0, R)$ and a failure parameter
$\beta \in (0,1)$, computes a matrix
$\Sigma$ satisfying
\[
0.99\Sigma
\preceq \E{x \sim K}{xx^\top}
\preceq 1.01\Sigma.
\]
It succeeds with probability at least $1-\beta$ and makes
\[
O\left(n^3\log\left(2+\frac{nR}{r\beta}\right)^{O(1)}\right)
\]
membership query calls.
\end{theorem}

The next corollary concerns the volume estimation.

\begin{corollary}
\label{cor:volume}
There is an algorithm that, given membership query access to a convex set
$K \subseteq \R^n$ satisfying
$B_n(x_0, r) \subseteq K \subseteq B_n(x_0, R)$ and an error parameter
$\eps \in (0,1)$, outputs with probability at least
$3/4$ a number within a factor
$1\pm\eps$ of the volume of
$K$ after making
\[
O\left(n^{13/4}\log(2+nR/r)^{O(1)} + n^3\eps^{-2}\log(n/\eps)^{O(1)}\right)
\]
membership queries.
If $K$ is symmetric and the two balls are centered at the origin, the query
bound improves to
\[
O\left(n^3\log(2+nR/r)^{O(1)} + n^3\eps^{-2}\log(n/\eps)^{O(1)}\right).
\]
\end{corollary}

The corollary is proved at the end of \cref{sec:asymmetric} from the two
theorems, the inner-ball bound of \cref{fact:kls-inner-ball}, and the Gaussian
cooling volume algorithm of \cite{CV18:journal} recalled in
\cref{fact:cv18-volume}.

The Gaussian cooling perspective gives a simpler route to rounding than the
previous ellipsoid-growth framework and improves the general membership-query
bound from $\O(n^{3.5})$ to
$\O(n^{13/4})$.
The algorithms are discussed in \cref{sec:overview} and presented in
\cref{sec:asymmetric,sec:symmetric}.

\subsection{Volume Estimation of Convex Bodies}
\label{subsec:polytopes}

In this section we discuss volume estimation algorithms for polytopes and how
various techniques (including ours) fit within the context of prior work.
The main primitive for volume estimation is \emph{sampling} from a convex body
$K$ using membership oracle queries.
Sampling from a convex body is done with Markov chain methods, for example the
\emph{Ball-Walk} or \emph{hit-and-run} algorithm.
The mixing time of these methods depends on the conductance of the polytope $K$,
where conductance refers to minimum surface area to small volume ratio of a
partition of
$K$ into two equal-size pieces.
So, in discussion of these kinds of volume algorithms, it was important to
understand the conductance of convex bodies in high dimensions.
The famous KLS conjecture asks whether this is always at least $\Omega(1)$, if
$K$ is \emph{isotropic}, i.e., the uniform distribution has covariance equal to
$I$.
Exciting recent progress has led to a bound of $\Omega(\log^{-1/4} n)$
\cite{L26:arxiv} (see \cref{subsec:highdim} for more discussion on these lines
of work), which in turn implies that the Ball-Walk mixes in
$\widetilde{\O}(n^2)$ steps for isotropic bodies.
Volume estimation requires sampling several points, and it is known that the
volume of an isotropic (even well-rounded) convex body can be estimated in
$\O(n^3)$ membership queries, using the Gaussian cooling algorithm of
Cousins-Vempala \cite{CV18:journal}, as mentioned above.
However, the question remains on how to transform any convex body to isotropic
position efficiently, and that is the focus of this paper.

There are several related lines of work on sampling general log-concave
distributions in different settings, such as sampling and volume estimation in
polytopes (that is, convex bodies explicitly described by their $m$ faces)
\cite{KN12:journal,N16:journal,NR17:journal,SV16:journal,CDWY17,CDWY18:journal,LLV20,LV22:journal,LV18,KLSV22,KLSV23,GKV24,MV23,GKMSZ24,JC25,GC11:journal}.
Several different algorithms are known in this setting, often motivated by
barrier functions from interior point methods.

\subsection{High-Dimensional Geometry}
\label{subsec:highdim}

The KLS conjecture of Kannan-Lovasz-Simonovits \cite{KLS95:journal} asks about
the minimal conductance of a cut of a convex body $K$.
The conjecture claims that if $K$ is isotropic, then the conductance is
$\Omega(1)$.
This is an important conjecture, because the KLS constant not only governs the
mixing time of Markov chains for sampling from $K$, but captures how well
functions distributed according to
$K$ concentrate.
For example, the thin-shell bound (concentration of the norm)
\[
\E{x \sim K}{(\|x\|_2 - n^{1/2})^2} \le O(1)
\]
would follow from the KLS conjecture.
Our sampling bounds use proved estimates toward the KLS conjecture.
The cooling analysis uses Brascamp--Lieb and Paouris inequalities in the
general case, and an even-sector Poincar\'e inequality in the symmetric case.

Recent progress towards the KLS conjecture started when Eldan introduced a new
method for proving lower bounds on the conductance called \emph{stochastic
localization}.
This initiated a chain of work
\cite{E13:journal,LV24:journal,C21:journal,KL22:journal,JLV22:arxiv,K23:journal,KL25:arxiv,L26:arxiv}
culminating in the current $\Omega(\log^{-1/4} n)$ KLS lower bound.
Related breakthroughs resolved the thin-shell and slicing conjectures
\cite{K06:journal,G24:arxiv,KL25:journal,F10:journal,GM11:journal,EL14:book}.
These analytic inputs are developed in
\cref{sec:overview,sec:asymmetric,sec:symmetric}.

\section{Preliminaries}
\label{sec:prelim}
\label{sec:preliminaries}

\paragraph{General notation.} The notation of $\Tilde{O}(\cdot)$ throughout indicates hidden dependence on the error parameter and logarithmic factors. We use $\R_{\ge0}$ to denote the nonnegative real numbers. For a twice-differentiable function $f: \R^n \to \R$ we let $\g f$ and $\g^2 f$ denote its gradient and Hessian respectively. We call a function $f:\R^n\to\R$ \emph{locally Lipschitz} if, for every $x\in\R^n$, there is a neighborhood $N_x$ of $x$ and a constant $L_x<\infty$ such that
\[
 |f(y)-f(z)|\leq L_x\|y-z\|_2
 \qquad\text{for all }y,z\in N_x.
\]
We first restate some notions and facts fundamental to high dimensional
geometry.

For positive semidefinite (PSD) matrices $A, B \in \R^{n \times n}$ we say that
$A$ and
$B$ are
$\alpha$-spectral-approximations (or
$\alpha$-approximations where it does not create confusion) for
$\alpha \ge 1$ if
$\alpha^{-1}A \pe B \pe \alpha A$ where
$\pe$ denotes the Loewner ordering of PSD matrices.
For a PSD matrix $A$ we let
$\|A\|_{\op} = \max_{\|v\|_2 = 1} v^\top Av$ denote its operator norm.

Furthermore, we abbreviate two sided multiplicative approximations using
$\alpha$.
For non-negative scalars $a, b \in \R_{\geq 0}$, we use
$a \approx_{\alpha} b$ to denote
$\frac{1}{\alpha} a \leq b \leq \alpha a$, and similarly for matrices
$A, B \in \R^{n \times n}$, we use
$A \approx_{\alpha} B$ to denote
$\frac{1}{\alpha} A \preceq B \preceq \alpha A$.

The Frobenius norm of a matrix $A$ is denoted as
$\|A\|_F \coloneq \Tr{A^\top A}^{1/2}$.
The operator norm of matrix $A$ is denoted by
$\|A\|_{\op}$ and defined to be
\[
\norm{A}_{\op}
\coloneq
\sup_{u \in \R^n : \norm{u}_2=1}
\norm{Au}_2.
\]

When the domain of integration is omitted, integrals are taken over $\R^n$.

\subsection{Probability}

A probability density function, or distribution, on $\R^{n}$ is a measure
$p: \R^{n} \rightarrow \R_{\geq 0}$ such that
$\int_{x} p(x) \, dx = 1$.
We use \emph{with high probability} (abbreviated w.h.p.) to mean that an event
holds with probability $1 - n^{-C}$ for any constant
$C > 0$ (possibly by adjusting constants in the definition of the event).

We measure distances between two probability density functions using total
variation distance.
The next definition concerns the total variation distance.

\begin{definition}
For two probability density functions $p, q$ on
$\R^n$ we define their total-variation distance as
\[
\TV\left(p, q\right)
=
\sup_A \abs{\int_A \left(p\left(x\right) - q\left(x\right)\right)\,dx},
\]
where $A \subseteq \R^n$ ranges over all measurable subsets of
$\R^n$.
Alternatively, we can write
\[
\TV\left(p, q\right)
= \frac{1}{2} \int_{\R^n} \abs{p\left(x\right) - q\left(x\right)}\,dx
= \frac{1}{2}\norm{p - q}_1.
\]
\end{definition}

Given a probability density function $p$ on
$\R^{n}$, and a function
$f: \R^{n} \rightarrow \R$, we use
$\E{x \sim p}{f(x)}$,
$\V{p}{f(x)}$ to denote expectation and variance of
$f(x)$.
Precisely, $\E{x \sim p}{f(x)} = \int_{x \in \R^{n}} f(x) p(x) dx$ and
\[
\V{x \sim p}{f\left(x\right)}
= \int_{x \in \R^{n}} \left(f\left(x\right) - \E{x \sim p}{f\left(x\right)}\right)^2 p\left(x\right)\,dx
= \E{x \sim p}{f\left(x\right)^2} - \E{x \sim p}{f\left(x\right)}^2.
\]
Sometimes to avoid notational clutter we will just write $\E{p}{f(x)}$ instead
of
$\E{x \sim p}{f(x)}$ when the variable being drawn from
$p$ is clear from context.

Now we define the mean and covariance matrix of a probability density function
on $\R^n$

\begin{definition}
\label{def:cov}
Given a probability density function $p$ on
$\R^n$, we denote its mean
$m_p$ as
\[
\R^n \ni m_p
\coloneq
\E{x \sim p}{x}
\]
and covariance matrix $\C_p$ as
\[
\R^{n \times n} \ni \C_p
\coloneq \E{x \sim p}{\left(x-m_p\right)\left(x-m_p\right)^\top}
= \E{x \sim p}{xx^\top} - m_pm_p^\top.
\]
\end{definition}

Matrix concentration bounds state that a sum of random matrices is close to the
sum of their expectations, assuming the norm of each matrix in their sum is not
too large.
Below we use the following version of Matrix Bernstein.

The next fact concerns the matrix Bernstein inequality, symmetric case of
Theorem 1.6.2 of \cite{T15:journal}
\footnote{\href{https://arxiv.org/pdf/1501.01571v1}{arXiv:1501.01571v1}}.

\begin{fact}
\label{prelim:mat_bern}
Let $Y_1, \dots, Y_m \in \R^{n \times n}$ be independent mean 0 random Hermitian
matrices such that for each
$i \in [m]$,
$\|Y_i\|_{\op} \leq R$ with probability
$1$.
Then
\[
\P{Y_1 \ldots Y_m}{\norm{\sum_{i = 1}^{m} Y_i}_{\op} \geq t}
\leq
2n \cdot \exp\left(
-\frac{t^2/2}{\norm{ \sum_{i}\E{Y_i}{Y_i^2}}_{\op}  + Rt/3}
\right)
\]
\end{fact}

\subsection{Convex geometry}

Let $B_n(x, \delta)$ denote the
$n$-dimensional Euclidean ball with radius
$\delta$ centered at
$x$.
Alternatively $B_n(x, \delta) = \{z \in \R^n : \|x-z\| \leq \delta\}$.

We say that a probability distribution (or density function) $p$ on
$\R^n$ is \emph{isotropic} if
$m_p = \vec{0}$ and
$\C_p = I_n$ where
$I_n \in \R^{n \times n}$ is the
$n \times n$ identity matrix.
We say that $p$ is \emph{near-isotropic} up to a factor of
$\alpha$ if
$m_p = \vec{0}$ and
$\C_p \approx_{\alpha} I_n$.

Let $K$ be a convex body in
$\R^n$.
Its associated indicator function $1_K: \R^n \to \R$ is given by
\[
1_K(x)
\coloneq
\begin{cases}
1 & \text{if $x \in K$}\\
0 & \text{otherwise}
\end{cases},
\]
and its volume is given by $\Vol(K) \coloneq \int_{x} 1_K(x) dx$.
We say that $K$ is isotropic if its associated probability density function
$p_K \coloneq \frac{1}{\Vol(K)} 1_K$ is isotropic.

\subsection{Log-Concave Distributions}\label{subsec:lc_dist}

In this section we collect some standard facts on log-concave distributions.
We first define a log-concave function.
The next definition concerns the log-concave.

\begin{definition}
A function $f: \R^n \rightarrow \R_{\ge0}$ is \emph{log-concave} if for all
$x,y \in \R^n$ and
$\alpha \in [0,1]$ it satisfies
\[
f(\alpha x + (1-\alpha)y) \geq f(x)^\alpha f(y)^{1-\alpha}
\]
\end{definition}
Throughout we refer to a log-concave probability density function as a
\emph{log-concave distribution} for brevity.
We now define a strongly log-concave density.

The next definition concerns $M$-strongly log-concave.

\begin{definition}
For a PSD matrix $M$, we call a log-concave function
$f(x)$
\emph{$M$-strongly log-concave} if
$f$ is twice-differentiable and satisfies
$\nabla^2(-\log(f(x))\se M$ for all
$x \in \R^n$.
For a positive real number $t$, we say that
$f$ is
$t$-strongly log-concave if
$\nabla^2(-\log(f(x)) \se tI$ for all
$x \in \R^n$.
\end{definition}

We note that log-concave distributions have sub-exponential tails.
The next lemma states that log-concave distributions have sub-exponential tails.

\begin{lemma}\label{prelim:lem:lg_subtails}
There exists a universal constant $c>0$ such that for every isotropic
log-concave distribution
$\mu$, for every unit vector
$v \in \R^n$, and for every
$t \ge 0$, the following holds:
    \[
\P{z \sim \mu}{|\langle v, z \rangle|\geq t} \leq \exp(-ct).
\]
\end{lemma}

We require the following bound on the values achieved by log-concave
distributions due to Fradelizi (we cite the version in \cite[Theorem
3.3]{P12:journal}).
\begin{lemma}
\label{lemma:frad}
Let $\mu(x) \propto e^{-\phi(x)}$ be log-concave with mean
$0$.
Let $x^* \coloneq \argmin_{x \in \R^n} \phi(x)$.
Then $\mu(x^*) \le e^n \mu(0)$.
Equivalently, $\phi(0) \le n + \phi(x^*)$.
\end{lemma}

We require the Brascamp-Lieb inequality.
This will be applied to bound the covariance of $M$-strongly log-concave
distributions.
The next lemma concerns the Brascamp-Lieb inequality.

\begin{lemma}\label{lem:brascamp_lieb}
For any log-concave distribution $p: \R^{n} \rightarrow \R_{\ge0}$ of the form
$p(x) = \exp(-\phi(x))$, and any differentiable function
$f: \R^{n} \rightarrow \R$ it holds that
\[
\V{x \sim p}{f(x)} \le \E{x \sim p}{\l \g f(x), \g^2 \phi(x)^{-1} \g f(x) \r}.
\]
\end{lemma}

The next corollary concerns the Brascamp-Lieb inequality for covariances.

\begin{corollary}
\label{cor:brascamp_lieb_cov}
Let $p(x)\propto e^{-\phi(x)}$ be a probability density on
$\mathbb{R}^n$ and suppose
$p$ is
$M$-strongly log-concave for some positive definite matrix
$M$.
Then
\[
\C_p \pe M^{-1}.
\]
\end{corollary}
\begin{proof}

We fix $v \in \R^n$ and set
$f(x)=\langle x,v\rangle$.
Then $\nabla f(x) = v$, so by the Brascamp-Lieb inequality (Lemma
\ref{lem:brascamp_lieb})
\[
\V{x \sim p}{f(x)}
\leq \E{x \sim p}{\langle \nabla f(x),(\nabla^2\phi(x))^{-1}\nabla f(x)\rangle}
= \E{x\sim p}{v^\top(\nabla^2\phi(x))^{-1} v}
\leq v^\top M^{-1} v
\]
where the final inequality comes from the definition of $M$-strongly
log-concave.
Since \(v^\top \C_p v=\V{x\sim p}{\langle x,v\rangle}\), this yields the claimed
spectral bound $\C_p \pe M^{-1}$.
\end{proof}

We require the standard fact that log-concavity is preserved under marginals
\cite{P73:journal}.
The next lemma concerns the Prekopa-Leindler inequality.

\begin{lemma}\label{lem:prekopa_leindler}
Let $p$ be a log-concave probability density on
$\R^n$ and let
$v \in \R^n$.
Then the distribution $\l x, v\r$ on
$\R$ (where
$x \sim p$) is a log-concave probability measure.
\end{lemma}
For $v\ne0$, this measure has density
\[
\phat(t)=\frac1{\|v\|_2}
\int_{\{x\in\R^n:\langle x,v\rangle=t\}}p(x)\,d\mathcal H^{n-1}(x),
\]
where $\mathcal H^{n-1}$ is the induced Euclidean surface measure.
For $v=0$, it is the point mass at zero.

We require the Paouris Inequality which allows us to bound the tail behavior of
the Euclidean norm of samples drawn from an isotropic log-concave distribution
\cite{P06:journal}.

The next lemma concerns the Paouris inequality.

\begin{lemma}\label{lem:prelim:paouris}
There is an absolute constant $C > 0$ such that for any isotropic log-concave
distribution
$p$ on
$\R^n$ and every
$t \geq 1$ it holds that
\[
\P{x \sim p}{\norm{x}_2 \geq Ct n^{1/2}}
\leq
\exp\left(-tn^{1/2}\right).
\]
\end{lemma}

We need the following inequality to bound moments of log-concave variables
\cite[Theorem~5.22]{LV07:journal}.

The next lemma concerns the reverse H{\"{o}}lder.

\begin{lemma}
\label{lem:reverse_holder}
Let $p$ be a log-concave probability density on $\R^n$.
For every positive integer $k$,
\[
\bigl(\E{p}{\|x\|_2^k}\bigr)^{1/k}
\leq 2k\,\E{p}{\|x\|_2}
\leq 2k\bigl(\E{p}{\|x\|_2^2}\bigr)^{1/2}.
\]
\end{lemma}

We use the following Gaussian-source form of Caffarelli's contraction theorem
\cite[Theorem~11]{C00:journal}.

\begin{theorem}
\label{thm:caffarelli-contraction}
Let $V:\R^n\to\R$ be twice continuously differentiable with
$\nabla^2V(x)\succeq2I$ for every $x\in\R^n$, and define
\[
d\gamma(x)=\pi^{-n/2}e^{-\|x\|_2^2}\,dx,
\qquad
d\mu(x)=Z^{-1}e^{-V(x)}\,dx,
\qquad
Z=\int_{\R^n}e^{-V(x)}\,dx.
\]
The quadratic-cost Brenier map $T=\nabla\varphi$ satisfying $T_\#\gamma=\mu$
is unique up to a $\gamma$-null set and has a $1$-Lipschitz representative.
In particular,
\[
\lVert DT(x)\rVert_{\mathrm{op}}\leq 1
\qquad\text{for $\gamma$-almost every $x$.}
\]
\end{theorem}

\subsection{Conductance of Log-Concave Distributions}\label{subsec:ball_walk}

We start by defining the conductance (KLS constant) of a log-concave
distribution.

\begin{definition}
\label{def:kls}
Let $p$ be a log-concave distribution on
$\R^n$.
We define its KLS constant $\psi_p$ as
\[
\frac{1}{\psi_p} \coloneq \inf_{S \subseteq \R^n} \frac{\int_{\partial S} p(x) \, dx}{\min\left\{\int_S p(x) \, dx, \int_{\R^n \setminus S} p(x) \, dx \right\}}.
\]
\end{definition}
Here, $\partial S$ denotes the \emph{boundary} of the set
$S$.
The KLS conjecture \cite{KLS95:journal} asks whether
$\psi_p = O(\|\C_p\|_{\op}^{1/2})$ for all log-concave distributions
$p$.
While this has not been proven, there has been significant recent work on
tightening the bound.
The current best bound is due to Letwin \cite{L26:arxiv}.

\begin{theorem}
\label{thm:klsbound}
If $p$ is a log-concave distribution, then
$\psi_p \leq O(\|\C_p\|_{\op}^{1/2} \log^{1/4} n)$.
\end{theorem}

The next inequality converts conductance bounds into variance bounds for smooth
test functions.
The next lemma concerns the Poincare inequality.

\begin{lemma}
\label{lemma:poincare}
Let $p$ be a log-concave distribution on
$\R^n$ and let
$f: \R^n \to \R$ be a differentiable function.
Then
\[
\V{x \sim p}{f\left(x\right)}
\le
O\left(\psi_p^2\right) \cdot \E{x \sim p}{\norm{\g f\left(x\right)}_2^2}.
\]
\end{lemma}

Two algorithmic black boxes from \cite{KV25} drive the full-density rounding
primitive of \cref{sec:grounded-core}: the near-cubic isotropic rounding of a
well-rounded log-concave law, and the estimation of the mean and covariance of
a well-rounded law with bounded covariance from one run of a Markov chain.
We state both with the exact hypotheses under which they are proved in
\cite{KV25}; the reductions of our inputs to these hypotheses are carried out
in \cref{sec:grounded-core}.
The constants in the following definition are those of
\cite[Assumption~4.4]{KV25}.

The next definition concerns well-rounded log-concave distributions.

\begin{definition}
\label{def:rounded}
Let $\pi\propto e^{-W}$ be a log-concave probability law on
$\R^n$, let
$z\in\R^n$, and let
$\rho,H,B>0$.
We say that $\pi$ is \emph{$(\rho,H,B)$-well-rounded around
$z$} if
\[
 B(z,\rho)\subseteq\{x\in\R^n: W(x)\leq\inf W+Hn\}
 \qquad\text{and}\qquad
 \E{X\sim\pi}{\|X-z\|^2}\leq Bn .
\]
We say that $\pi$ is \emph{grounded around
$z$ with respect to
$W$} if
$\rho=1$ and
$H=10$ can be taken, and in addition
$\min W\geq0$ and
$W\leq10n$ on the support of
$\pi$, so that the support of
$\pi$ equals the level set
$\{W\leq\inf W+10n\}$.
\end{definition}
Note that well-roundedness is a relaxation of near-isotropy.

The following adapts the well-rounded isotropizer of
\cite[Algorithm~4 and Lemma~4.11]{KV25} under their Assumption~4.4.
We record its output through an affine map and give the implementation
and proof below.

\begin{theorem}
\label{prelim:thm:wellr_iso}
There is a universal constant $c_{\rm iso}>0$ and, for every
$C\geq1$, there are constants
$A_C,a_C>0$ such that the following holds.
Let $\pi\propto e^{-W}$ be a log-concave law on
$\R^n$ that is grounded around a known point
$z$ with respect to
$W$ (\cref{def:rounded}), with
$\E{X\sim\pi}{\|X-z\|^2}\leq C^2n$ for the known constant
$C$, and let
$W$ be given by an evaluation oracle.
There is a randomized algorithm which, on an event of probability at least
$1-c_{\rm iso}n^{-1/2}$, returns an invertible affine map
$F:\R^n\to\R^n$ with
\[
 C_{F_\#\pi}\approx_{1.01}I_n
 \qquad\text{and}\qquad
 \|m_{F_\#\pi}\|\leq0.01,
\]
and the expected number of evaluation queries it makes on this event is at
most
$A_Cn^3\log^{a_C}(2n)$.
\end{theorem}
Writing $F(x)=A(x-\bar x)$ with
$A$ invertible and
$\Sigma=(A^\top A)^{-1}$, the conclusion reads
$\Sigma\approx_{1.01}C_\pi$ and
$\|\Sigma^{-1/2}(\bar x-m_\pi)\|\leq0.01$, because
$C_{F_\#\pi}=AC_\pi A^\top$ and
$\|\Sigma^{-1/2}v\|=\|Av\|$ for every
$v$.

\begin{proof}[Implementation and proof of \cref{prelim:thm:wellr_iso}]
We specify the rounding primitive using the sampling and stationary
covariance-estimation results of \cite{KV25,KZ24:arxiv}.
Translate the original input by its known point $z$, so its potential is
$W_0(x)=W(x+z)$ and $\E{\pi_0}{\|X\|^2}\leq C^2n$.
All affine maps below are recorded and composed with this translation.
Set
\[
 L=\max\{4,\log(2n)\},\qquad q=L^{-1},\qquad
 d_0=\frac{nq^2}{100},\qquad r_*=\frac{q^2\sqrt n}{100}.
\]
Start with $\pi_0$ and $r_0=1$.
At stage $i$, write $m_i,C_i$ for the true mean and covariance of the current
law $\pi_i$.
While $r_i<r_*$, estimate its covariance by a positive semidefinite empirical
matrix $S_i$ satisfying
\begin{equation}
\label{eq:primitive-covariance-accuracy}
 -\tfrac1{10}C_i-d_0I\preceq S_i-C_i
 \preceq\tfrac1{10}C_i+d_0I.
\end{equation}
Let $P_i$ be the spectral projection of $S_i$ onto eigenvalues at most $n$,
put $M_i=I+P_i$, and update
\[
 \pi_{i+1}=(M_i)_\#\pi_i,\qquad
 r_{i+1}=2(1-q)r_i.
\]
The estimates, their probability of success, and their query cost are
justified below. Their required accuracy differs from a constant additive
multiple of $n$ only by a logarithmic factor.

We first give the deterministic geometric argument on the event that
\eqref{eq:primitive-covariance-accuracy} holds at every completed stage.
The covariance obeys $C_{i+1}=M_iC_iM_i$, whence
$\operatorname{tr}C_{i+1}\leq4\operatorname{tr}C_i$.
The lower estimate in \eqref{eq:primitive-covariance-accuracy} gives
$P_iC_iP_i\preceq2nP_i$. Consequently,
\[
 \|C_{i+1}\|_{\op}
 =\|C_i^{1/2}(I+3P_i)C_i^{1/2}\|_{\op}
 \leq\|C_i\|_{\op}+6n.
\]
There are $O(L)$ stages, because $2(1-q)\geq3/2$ and $r_*\leq\sqrt n$.
Since $(1-q)^{-O(L)}=O(1)$, these observations imply
\begin{equation}
\label{eq:primitive-covariance-growth}
 \operatorname{tr}C_i=O_C(r_i^2n),
 \qquad
 \|C_i\|_{\op}=O_C(nL).
\end{equation}

We next prove that the support of $\pi_i$ contains a ball $B(c_i,r_i)$
for some center $c_i$; the algorithm need not know this center.
The assertion holds initially.
The support also contains
$m_i+e^{-1}C_i^{1/2}B(0,1)$ by the log-concave support-ball bound
\cite[Lemma~A.10]{KV25}.
Suppress the index $i$ and put $M=I+P$.
After applying $M$, the support contains the two ellipsoids
\[
 A=M B(c_i,r_i),\qquad
 B=M\bigl(m_i+e^{-1}C_i^{1/2}B(0,1)\bigr),
\]
and hence their weighted Minkowski sum $(1-q/2)A+(q/2)B$.
Center this sum at
$c'=(1-q/2)Mc_i+(q/2)Mm_i$.
For a unit vector $u$, write $a=\|(I-P)u\|$.
The centered support function of $A$ is
$r_i\|Mu\|=r_i\sqrt{4-3a^2}\geq r_i(2-\tfrac32a^2)$.
Since $S_i\succeq n(I-P)$, the upper estimate in
\eqref{eq:primitive-covariance-accuracy} gives
\[
 u^\top MC_iMu\geq\frac{na^2-4d_0}{1.1}.
\]
If $a^2\leq8d_0/n$, the contribution of $A$ alone is at least
$2r_i-qr_i-0.12q^2r_i\geq2(1-q)r_i$.
Otherwise the centered support function of $B$ is at least
$\sqrt n\,a/(e\sqrt{2.2})$.
Its weight $q/2$ compensates the term $\tfrac32r_i a^2$, because
\[
 r_i<r_*=\frac{q^2\sqrt n}{100}
 \leq\frac{q\sqrt n}{3e\sqrt{2.2}}.
\]
The support function of the sum is therefore at least
$2(1-q)r_i$ in every unit direction.
Thus the sum contains $B(c',2(1-q)r_i)$, proving the induction.
This argument uses support functions of the full ellipsoids and does not
require their principal sections to be large.

Here is the sampling implementation.
Before each covariance-estimation stage, including the final one, independently
apply \cite[Theorem~3.16]{KV25} to the original input $W_0$, with accuracy
$\log2$.
Its hypotheses hold with known center $0$ and $R^2=C^2n$.
The construction uses evaluations of $W_0$, including $W_0(0)$, and does
not require the value of $\inf W_0$.
It generates a fresh ideal sample from a law bounded by twice the original
exponential lift.
Push this sample forward by the current recorded linear map, leaving the
lift coordinate $t$ unchanged.
Conditioned on the preceding stages, the resulting ideal starting law is
bounded by twice the lift of $\pi_i$.

For sampling at stage $i$, rescale only the $x$-coordinates by $r_i^{-1}$.
The ground set then has unit inradius.
By \cite[Lemma~2.5]{KV25}, the lifted covariance has operator norm at most
\[
 2\bigl(\|C_i\|_{\op}/r_i^2+160\bigr).
\]
The exact proximal kernel is reversible.
Its $\chi^2$ contraction, the Poincare bound, and
\cite[Lemma~2.14]{KV25} provide a thinned chain with spectral gap at least
$0.99$, using
\[
 \widetilde O\!\left(n^2\left(\|C_i\|_{\op}/r_i^2+1\right)\right)
\]
queries per retained sample, with logarithmic dependence on its abort budget.
The mixing lengths and sample counts use deterministic upper bounds from
\eqref{eq:primitive-covariance-growth}, computed from the known $C,n,r_i$
and stage index; they do not require the unknown covariance $C_i$.
Undoing the known coordinate rescaling preserves this spectral gap.
The resulting stationary law is the unscaled exponential lift, whose
covariance has trace at most $\operatorname{tr}C_i+160$ and operator norm
at most $2(\|C_i\|_{\op}+160)$.
Apply \cite[Theorem~8]{KZ24:arxiv} to this lifted chain with multiplicative
accuracy $1/10$ and additive accuracy $d_0$.
Using \eqref{eq:primitive-covariance-growth}, its sample requirement is
$k_i=\widetilde O_C(r_i^2)$.
Compressing the empirical covariance to the $x$-coordinates gives
\eqref{eq:primitive-covariance-accuracy} with stationary-start failure
probability $O(1/n)$.
Since $r_i^2\leq n$ at every updating stage, the total sampling cost of
that stage is $\widetilde O_C(n^3)$.

At termination, $r_i\geq\max\{1,r_*\}$.
The per-sample bound above is now $\widetilde O_C(n^2)$.
Use $\widetilde O(n)$ retained samples for the final empirical mean
$\bar m$ and covariance $S$.
If $S$ is not positive definite, declare failure.
For their analysis, whiten the $x$-coordinates by $C_i^{-1/2}$ and leave
$t$ unchanged; this is an invertible change of variables of the stationary
chain and preserves its spectral gap.
The resulting lifted covariance has operator norm at most $322$.
The stationary covariance and mean estimates
\cite[Corollaries~10 and~33]{KZ24:arxiv}, with sufficiently small fixed
accuracy, therefore give
\[
 (1-\varepsilon)C_i\preceq S\preceq(1+\varepsilon)C_i,
 \qquad
 \|C_i^{-1/2}(\bar m-m_i)\|\leq\varepsilon,
 \qquad \varepsilon=10^{-3},
\]
except with probability $O(1/n)$.
Whitening here is used only in the analysis: the algorithm computes the
ordinary empirical statistics in its current coordinates.
The final map $x\mapsto S^{-1/2}(x-\bar m)$ has covariance between
$(1+\varepsilon)^{-1}I$ and $(1-\varepsilon)^{-1}I$, and its mean has norm
at most $\varepsilon/\sqrt{1-\varepsilon}<0.01$.
Compose it with all recorded maps to obtain the required $F$.

Finally, we account for failure and queries without conditioning a sampler
on successful rejection steps.
Fix a deterministic bound $J=\lceil4L\rceil$ on the number of
covariance-estimation stages.
Assign every warm-start run and every sampling run an abort budget
$n^{-2}/(2J)$.
Ideal processes use rejection sampling without a cutoff.
Couple each implemented run to its ideal run using the same proposals until
a cutoff is exceeded.
The per-step estimates in \cite[Lemmas~2.14 and~3.10]{KV25}, summed along
the ideal processes on accurate preceding stages, bound the probability of
an abort before any earlier statistical failure by $n^{-2}$.
For any event $E$ of a stage's ideal chain, its conditional probability,
given the preceding stages, is at most twice the corresponding
stationary-start probability: if $\nu_i\leq2\widehat\pi_i$, then
\[
 \mathbb P_{\nu_i}(E)
 =\int\mathbb P_x(E)\,\nu_i(dx)
 \leq2\int\mathbb P_x(E)\,\widehat\pi_i(dx).
\]
Thus each first statistical failure, conditional on accurate preceding
estimates, has probability $O(1/n)$.
A union bound gives overall failure probability
$O(J/n)+n^{-2}=O(n^{-1/2})$.
Each fresh start costs $\widetilde O_C(n^3)$, as does each updating stage
and the final stage.
The same per-step estimates bound the expected capped query cost along
accurate prefixes; summing over $J=O(L)$ stages gives the stated expected
$\widetilde O_C(n^3)$ query bound on the success event.
As usual, fixed small dimensions are absorbed into the universal constants;
the later application explicitly pads to a sufficiently large dimension.
\end{proof}

The second black box is the final step of \cite[Algorithm~4]{KV25}, run on a
law whose covariance is already known to be bounded: a warm start from
\cite[Theorem~3.16]{KV25}, one run of the proximal sampler of
\cite[Section~2]{KV25}, and the dependent-sample mean and covariance estimator
of \cite[Corollaries~10 and~33]{KZ24:arxiv}, as recorded in
\cite[Lemma~4.6]{KV25}.
We state the resulting guarantee and then record how it is assembled from the
cited statements.

\begin{theorem}
\label{thm:moment-estimation}
There is a universal constant $c_{\rm est}>0$ and, for every
$c_+\geq1$ and every
$\xi\in(0,1)$, there are constants
$A,a>0$ such that the following holds.
Let $\pi\propto e^{-W}$ be a log-concave law on
$\R^n$ with
$\|C_\pi\|_{\op}\leq c_+$, let
$W$ be given by an evaluation oracle, and let a point
$z$ and a radius
$R>0$ be known with
\[
 B(z,1)\subseteq\{W\leq\inf W+10n\}
 \qquad\text{and}\qquad
 \E{X\sim\pi}{\|X-z\|^2}\leq R^2 .
\]
For every $\eta\in(0,1)$ there is a randomized algorithm which, on an event of
probability at least
$1-2\eta-c_{\rm est}/n$, returns
$(\widehat m,\widehat C)$ with
\[
 (1-\xi)C_\pi\preceq\widehat C\preceq(1+\xi)C_\pi
 \qquad\text{and}\qquad
 \|\widehat m-m_\pi\|\leq\xi,
\]
and the expected number of evaluation queries it makes on this event is at
most
$A\,(n^2R^2+n^3)\log^{a}(2n/\eta)$.
\end{theorem}

The theorem is extracted from \cite{KV25} as follows.
Let $\widehat\pi(x,t)\propto e^{-nt}$ on
$\mathcal K=\{(x,t)\in\R^n\times\R: W(x)\leq nt\}$ be the exponential
reduction (exp-red) of \cite[Section~2.1]{KV25}: it is log-concave, its
$x$-marginal is
$\pi$, and its covariance
$\widehat\Sigma$ satisfies
$\|\widehat\Sigma\|_{\op}\leq2(\|C_\pi\|_{\op}+160)\leq2c_++320$ by
\cite[Lemma~2.5]{KV25}.
The two displayed hypotheses are those of \cite[Problems~2.1 and~3.1]{KV25}
with
$x_0=z$.
First, use \cite[Theorem~3.16]{KV25} with accuracy $\log2$ and abort budget
$\eta$ to generate a warm start. Its ideal process uses rejection sampling
without a cutoff and has output law $\nu\leq2\widehat\pi$; the implemented
algorithm agrees with that output except on an abort event of probability
at most $\eta$. The expected query bound is
$\O(n^2(R^2\vee n)\log^8(1/\eta))$.

Let $P$ be the exact proximal-sampler kernel for $\widehat\pi$.
It is reversible, and \cite[Lemma~2.9]{KV25}, applied in dimension $n+1$, gives
\[
 \chi^2(\lambda P\Vert\widehat\pi)
 \leq(1+h/C_{\rm PI}(\widehat\pi))^{-2}
       \chi^2(\lambda\Vert\widehat\pi).
\]
The covariance bound above and \cref{thm:klsbound,lemma:poincare} imply
\[
 C_{\rm PI}(\widehat\pi)
 =O((c_++1)\log^{1/2}(n+2)).
\]
Choose $N$ so that $P^N$ has spectral gap at least $0.99$, using this known
upper bound in place of the unknown Poincar\'e constant. Set
\[
 \xi'=\frac{\xi}{\sqrt{2c_++320}},\qquad
 k\asymp(\xi')^{-2}(n+1)
 \log^2\!\frac{n+2}{\xi'}\,\log^2(n+2).
\]
Choose the step size and rejection cutoff by \cite[Lemma~2.14]{KV25} for
$kN$ steps, warmness $M=2$, and total abort budget $\eta$.
These choices can be made with
$N=O((c_++1)n^2\log^{O(1)}((n+2)/\eta))$, where constants may depend on
$\xi$ and $c_+$.

For the ideal chain started from $\nu$, every marginal is $2$-warm because
$\nu P^j\leq2\widehat\pi P^j=2\widehat\pi$.
The per-step estimates in that lemma, summed along the ideal chain, bound
the probability that any rejection cutoff is exceeded by $\eta$ and the
expected implemented query count by $\O(kN\log^2(kN/\eta))$.
Couple the implemented and ideal chains using the same proposals until the
first abort. This argument does not condition the chain on successful
rejection steps.

Let $Z_1,\ldots,Z_k$ be every $N$-th ideal iterate.
Under a stationary start, \cite[Lemma~4.6]{KV25} gives empirical mean
$\overline m$ and covariance $\overline\Sigma$ satisfying
\[
 \begin{gathered}
 -\xi'\widehat\Sigma\preceq\overline\Sigma-\widehat\Sigma
 \preceq\xi'\widehat\Sigma,\\
 \|\overline m-m_{\widehat\pi}\|^2
 \leq(\xi')^2\|\widehat\Sigma\|_{\op}\leq\xi^2,
 \end{gathered}
\]
except with probability $O(1/n)$.
The same conclusion holds from $\nu$ with at most twice this failure
probability: for every path event $A$,
\[
 \mathbb P_\nu(A)=\int\mathbb P_z(A)\,\nu(dz)
 \leq2\int\mathbb P_z(A)\,\widehat\pi(dz)
 =2\mathbb P_{\widehat\pi}(A).
\]
Taking the $x$-blocks of $\overline m$ and $\overline\Sigma$ gives
$\widehat m$ and $\widehat C$ with the claimed bounds, since $\xi'\leq\xi$.
Outside the two abort events, these are exactly the implemented outputs.
A union bound gives total failure probability at most $2\eta+c_{\rm est}/n$.
Adding the warm-start and chain costs gives the stated query bound.

\subsection{Additional conventions}

The algorithmic parts of the paper distinguish centered covariance from the
uncentered second moment, and the technical upper bounds are stated through an
evaluation oracle for a convex potential before being specialized back to
membership oracles.
We record the relevant notation here so that the main body can remain focused on
the main membership-oracle convex-body second moment estimation.

We will also use the following layer-cake consequence of Tonelli's theorem.
The next lemma concerns the layer-cake identity.

\begin{lemma}
\label{lem:layer-cake}
Let $Y,Y'$ be random variables on a common probability space satisfying
$Y,Y'\ge a$ almost surely, let
$s,t>-a$, and let
$k,\ell>0$.
Then, with the expectation and the probability taken over the joint law of
$(Y,Y')$,
\begin{align*}
\E{}{(Y+s)^{-k}(Y'+t)^{-\ell}}
&=\int_a^\infty\!\!\int_a^\infty
 \mathbb P(Y\le x,Y'\le y)
 k(x+s)^{-k-1}\ell(y+t)^{-\ell-1}\,dx\,dy.
\end{align*}
The corresponding one-variable identity is
\[
(Y+s)^{-k}=\int_a^\infty
\ind_{\{Y\le x\}}k(x+s)^{-k-1}\,dx.
\]
\end{lemma}

\begin{proof}
The one-variable identity follows by integrating the derivative of
$u\mapsto(u+s)^{-k}$ from
$Y$ to
$\infty$.
Multiplying the two identities and applying Tonelli's theorem to the resulting
nonnegative integrand gives the two-variable formula.
\end{proof}

\subsection{Matrices, moments, and affine changes}
\label{sec:prelim-moments}

Vectors are columns, $\|\cdot\|$ is the Euclidean norm, and we define a ball of
radius
$r$ centered at
$x$ as
$B(x,r)=\{y:\|y-x\|\leq r\}$; we abbreviate
$B_2^n=B(0,1)$.
For symmetric matrices, $A\preceq B$ means that
$B-A$ is positive semidefinite.
We write $A^{-\mathsf T}=(A^{-1})^{\mathsf T}$ for an invertible matrix
$A$.

The next definition concerns the moment matrices and relative Loewner
approximation.

\begin{definition}
\label{def:moment-matrices}
Let $\pi$ be a probability law on
$\mathbb R^n$ with finite second moment, and let
$X\sim\pi$.
Its mean, uncentered second moment, and centered covariance are
\[
m_\pi:=\E{X\sim\pi}{X},
\qquad
M_\pi:=\E{X\sim\pi}{XX^{\mathsf T}},
\qquad
C_\pi:=\E{X\sim\pi}{(X-m_\pi)(X-m_\pi)^{\mathsf T}}.
\]
Thus
\begin{equation}
\label{eq:raw-centered}
 M_\pi=C_\pi+m_\pi m_\pi^{\mathsf T}.
\end{equation}
For positive definite $A,B$ and
$\kappa\geq1$, write
\[
A\approx_\kappa B
\quad
\Longleftrightarrow
\quad
\kappa^{-1}B
\preceq A
\preceq\kappa B.
\]
When an accuracy is given additively, we state the oriented inequalities
explicitly; for example, $(1-\varepsilon)B\preceq A\preceq(1+\varepsilon)B$.
For a convex body $K$, let
$\pi_K=\operatorname{Unif}(K)$ and abbreviate
$m_K:=m_{\pi_K}$,
$C_K:=C_{\pi_K}$, and
$M_K:=M_{\pi_K}$.
\end{definition}

If a law is even, then its mean is zero and its uncentered second moment is its
covariance.
In particular, this holds for $\pi_K$ when
$K=-K$.
This identity explains why the symmetric result may speak of covariance, whereas
the asymmetric results must retain the phrase \emph{uncentered
second moment}.

The next definition concerns the affine pushforwards.

\begin{definition}
\label{def:affine-pushforward}
Let $\mathcal T(x)=B_{\rm aff}(x-b)$, where
$B_{\rm aff}$ is known and invertible, and write
$\mathcal T_\#\pi$ for the law of
$\mathcal T(X)$ when
$X\sim\pi$.
Then
\begin{align}
 m_{\mathcal T_\#\pi}
   &=B_{\rm aff}(m_\pi-b), \label{eq:affine-mean}\\
 C_{\mathcal T_\#\pi}
   &=B_{\rm aff}C_\pi B_{\rm aff}^{\mathsf T}, \label{eq:affine-cov}\\
 M_{\mathcal T_\#\pi}
   &=B_{\rm aff}\!\left(C_\pi+
      (m_\pi-b)(m_\pi-b)^{\mathsf T}\right)B_{\rm aff}^{\mathsf T}.
      \label{eq:affine-raw}
\end{align}
In particular, Loewner inequalities between covariance matrices are preserved by
affine congruence.
Raw second moments obey the same congruence under a linear map; under a
translation one must use Equation~\eqref{eq:affine-raw}, rather than treating a
raw moment as a covariance.
\end{definition}

\subsection{Oracle models and adaptive randomized calls}
\label{sec:prelim-oracles}

This subsection separates the binary membership model used for convex bodies
from the exact-real evaluation model used for general log-concave laws.
It also records the conditional-success and deterministic query-limit
conventions needed when the later algorithms choose their oracle calls
adaptively.

The next definition concerns the membership and evaluation oracles.

\begin{definition}
\label{def:oracle-models}
A membership oracle for a convex body $K$ returns the bit
$\ind_{\{x\in K\}}$ on input
$x\in\mathbb R^n$.
An exact-real evaluation oracle for an extended-valued convex potential
$V:\mathbb R^n\to\mathbb R\cup\{+\infty\}$ returns
$V(x)$ on input
$x$.
The associated log-concave law has density
\[
\pi(dx)=Z_V^{-1}e^{-V(x)}\,dx,
\qquad
Z_V=\int_{\mathbb R^n}e^{-V(x)}\,dx.
\]
Neither $Z_V$ nor the irrelevant additive normalization of
$V$ is assumed as known.
For a convex body $K$, its indicator potential
\[
 I_K(x)=
 \begin{cases}
 0,&x\in K,\\
 +\infty,&x\notin K
 \end{cases}
\]
satisfies $e^{-I_K}=\ind_{K}$; hence evaluation access to
$I_K$ is exactly membership access to
$K$.
Computation between oracle calls is not charged unless explicitly stated.
\end{definition}

General log-concave densities require their potential values, while the
uniform-body specialization $V=I_K$ uses only membership.
This is a model distinction, not an impossibility result.

The next definition concerns the conditional guarantees and query limits.

\begin{definition}
\label{def:adaptive-oracle-guarantee}
Suppose a randomized routine is called on inputs chosen from earlier outputs.
A failure guarantee $p$ \emph{under adaptive conditioning} means that,
conditional on every realized pre-call history for which the routine's displayed
input promises hold, fresh internal randomness makes the call succeed with
probability at least
$1-p$.
A deterministic query limit is a bound that holds for every valid input, every
history, and every realization of the internal randomness, including executions
that return a failure symbol.
\end{definition}

This convention allows failure budgets to be assigned chronologically and summed
by conditional union bounds.
It also keeps correctness events separate from the all-outcomes query limits
used in the algorithmic theorems.

The next definition concerns the grounded evaluation-oracle input.

\begin{definition}
\label{def:grounded-input}
Fix a universal constant $h_0>0$.
A grounded input in dimension $n\geq1$ consists of a proper lower-semicontinuous
convex function
$V:\mathbb R^n\to\mathbb R\cup\{+\infty\}$, exact-real evaluation access to
$V$, numbers
$r,R>0$, and a known point
$x_0\in\mathbb R^n$ such that
\[
0
<\int_{\mathbb R^n}e^{-V(x)}\,dx
<\infty,
\qquad
B(x_0,r)\subseteq\{x:V(x)\leq\inf V+h_0n\},
\]
and, for $X\sim\pi\propto e^{-V}$,
\[
 \E{X\sim\pi}{\|X-x_0\|^2}\leq R^2.
\]
Constants hidden by asymptotic notation may depend on $h_0$, but not on
$n,r,R$, or a caller-supplied failure probability.
\end{definition}

Here and below, a sublevel set of $V$ is a set of the form
$\{x:V(x)\leq a\}$.
The promise concerns one such set, not necessarily the entire support.
It gives a known ball in a low sublevel set and a global second-moment bound;
the algorithm need not know $\inf V$.

\subsection{Finite-order change of measure}
\label{sec:prelim-renyi}

This subsection fixes the change-of-measure quantity used to pass between
successive laws on the cooling path.
Distinguishing the logarithmic divergence from its exponentiated form will keep
the warmness assumptions and one-step estimates in the overview and upper-bound
proof unambiguous.

The next definition concerns the R\'enyi divergence and finite-order warmness.

\begin{definition}
\label{def:renyi-divergence}
Let $\nu$ and
$\pi$ be probability laws and let
$q>1$.
If $\nu\ll\pi$, define
\[
D_q(\nu\Vert\pi) :=\frac{1}{q-1}\log \int\left(\frac{d\nu}{d\pi}\right)^q\,d\pi,
\qquad
\mathcal R_q(\nu\Vert\pi):=\exp(D_q(\nu\Vert\pi));
\]
otherwise both quantities are $+\infty$.
In particular,
\[
 \mathcal R_2(\nu\Vert\pi)
 =\int\left(\frac{d\nu}{d\pi}\right)^2d\pi.
\]
We use $D_q$ for the logarithmic divergence and
$\mathcal R_q$ for its exponentiated finite-order warmness.
\end{definition}

The order-two orientation matters: in the cooling argument, the new target is
compared to the law from which the rounding call inherits its state.
Writing both arguments in \cref{def:renyi-divergence} prevents this direction
from being suppressed.

\subsection{Standard log-concave facts}
\label{sec:prelim-log-concave}

The following proposition records exactly the background consequences used
later.
Its constants are universal and may change from one occurrence to the next.

The next proposition concerns the log-concave toolkit.

\begin{proposition}
\label{prop:log-concave-toolkit}
Let $\nu$ be a full-dimensional log-concave probability law on
$\mathbb R^n$ with potential
$W$, mean
$m$, and covariance
$C$, and let
$Y\sim\nu$.
\begin{enumerate}[label=\textnormal{(\roman*)}]
\item \label{itm:log-concave-toolkit:1} The barycenter inequality gives
\[
 W(m)\leq\inf W+n.
\]
\item \label{itm:log-concave-toolkit:2} If $m=0$ and
  $\frac12I\preceq C\preceq2I$, then there is a universal
  $c_P>0$ such that, for every
  $u\geq1$,
\[
 \Pr\!\left(\|Y\|\geq c_P \cdot un^{1/2}\right)
 \leq \exp(-un^{1/2}/c_P).
\]
\item \label{itm:log-concave-toolkit:3} There is a universal $C_4>0$ such that
  every affine functional
  $\ell$ satisfies
\[
 \E{Y\sim\nu}{|\ell(Y)|^4}
 \leq C_4\bigl(\E{Y\sim\nu}{|\ell(Y)|^2}\bigr)^2.
\]
\item \label{itm:log-concave-toolkit:4} There are universal $\rho_L,H_L>0$ such
  that, if
  $m=0$ and
  $\frac12I\preceq C\preceq2I$, then
\[
 B(0,\rho_L)\subseteq\{W\leq\inf W+H_Ln\}.
\]
Under the same covariance bounds, the conclusion with radius $\rho_L/2$ still
holds when the centering hypothesis is relaxed to
$\|m\|\leq\rho_L/2$.
\item \label{itm:log-concave-toolkit:5} Every affine functional $\ell$ obeys
\[
 \bigl(\E{Y\sim\nu}{|\ell(Y)|^2}\bigr)^{1/2}
 \leq4\E{Y\sim\nu}{|\ell(Y)|}.
\]
Moreover, if $m=0$ and
$s^{-1}I\preceq C\preceq sI$ for
$s\geq1$, then
\[
 B\!\left(0,\frac{1}{es^{1/2}}\right)
 \subseteq\operatorname{supp}\nu.
\]
\item \label{itm:log-concave-toolkit:6} If $W$ is twice continuously
  differentiable and strictly convex on its support, then every locally Lipschitz
  $g$ in the associated Sobolev space satisfies the Brascamp--Lieb variance
  inequality
\[
 \V{Y\sim\nu}{g}
 \leq\E{Y\sim\nu}{
 \nabla g(Y)^{\mathsf T}
 \bigl(\nabla^2W(Y)\bigr)^{-1}\nabla g(Y)}.
\]
Consequently, $\nabla^2W\succeq\alpha I$ implies
$C\preceq\alpha^{-1}I$ and
$\V{Y\sim\nu}{g}\leq \alpha^{-1}\E{Y\sim\nu}{\|\nabla g\|^2}$.
\end{enumerate}
\end{proposition}

The barycenter estimate in part (i) is due to Fradelizi \cite{F97:journal}.
Part (ii) is Paouris' norm concentration \cite{P06:journal}.
The one-dimensional moment and level-set estimates of Lov\'asz and Vempala
\cite{LV07:journal} imply the precise reverse-H\"older, support, and level-set
consequences in parts (iii)--(v); their short reductions are retained with the
full-density-rounder proof in \cref{sec:grounded-core}.
Part (vi) is the classical Brascamp--Lieb inequality \cite{BL76:journal}.
The later cooling proof establishes the nonsmooth extended-valued form it needs,
rather than silently applying part (vi) beyond its hypotheses.

For orientation, parts (iv)--(v) are used to pass from a bounded restriction to
the full law in \cref{sec:grounded-core}; parts (i), (ii), and (vi) enter the
cooling estimates in \cref{sec:cooling-estimates}; and the fourth-moment and
reverse-H\"older bounds in parts (iii) and (v) are used in the moment-transfer
and endpoint arguments in
\cref{sec:cooling-estimates,sec:endpoint-reconstruction}.

The next fact is Theorem~4.1 of \cite{KLS95:journal}, the inner ball of an
isotropic convex body, in the scaled form used for the volume corollary.

\begin{fact}
\label{fact:kls-inner-ball}
Let $K\subseteq\R^n$ be a convex body with centroid
$m$ and covariance
$C$.
If $C\succeq\lambda I$ for some
$\lambda>0$, then
$B(m,\lambda^{1/2})\subseteq K$.
\end{fact}

\begin{proof}
The body $C^{-1/2}(K-m)$ has centroid
$0$ and covariance
$I$, so Theorem~4.1 of \cite{KLS95:journal} gives
$B(0,((n+2)/n)^{1/2})\subseteq C^{-1/2}(K-m)$.
Hence $K-m\supseteq C^{1/2}B(0,1)\supseteq B(0,\lambda^{1/2})$, because
$\|C^{-1/2}u\|\leq\lambda^{-1/2}\|u\|$ for every
$u$.
\end{proof}

The next fact is Theorem~1.1 of \cite{CV18:journal}, in the form stated there
for a general second moment.

\begin{fact}
\label{fact:cv18-volume}
There is an algorithm that, for any $\eps,p\in(0,1)$ and any convex body
$K\subseteq\R^n$ containing the unit ball and satisfying
$\E{x\sim K}{\|x\|^2}\leq R^2$, with probability at least
$1-p$ outputs a number within a factor
$1\pm\eps$ of
$\Vol(K)$ after
\[
O\left(\frac{\max\{R^2n^2,n^3\}}{\eps^2}
\log^2n\,\log^2\frac1\eps\,\log^2\frac n\eps\,\log\frac1p\right)
\]
membership queries.
\end{fact}

\section{Overview}
\label{sec:overview}

We begin with the general-body scalar Gaussian-cooling algorithm.
Previous general rounding algorithms followed an ellipsoid-growth framework, and
the best known membership-query bound had remained at $\O(n^{3.5})$.
Our scalar cooling path gives a more direct analysis and improves the query
bound to $\O(n^{13/4})$.
After describing this general-body algorithm, we explain the symmetric
matrix-cooling algorithm, which uses symmetry to reach cubic complexity.
We then explain how the counterexamples to natural conjectural shortcuts fit
with this algorithmic picture.

\subsection{General-Body Scalar Cooling Algorithm}

The scalar cooling path has an inherent limitation: the single temperature $t$
relaxes the Gaussian factor by the same amount in every direction.
If only a few directions are close to becoming poorly conditioned, the scalar
update must slow down all directions anyway.
The general-body algorithm still gets a speedup by combining this scalar cooling
path with a stronger near-cubic rounding primitive for already well-rounded
log-concave distributions.
This lets the algorithm decrease the temperature by a factor of
$1-\Theta(n^{-1/4})$ while preserving enough finite-order warmness between
consecutive targets.

For an asymmetric body, the target matrix is naturally the uncentered second
moment
\[
M_K
=\E{x\sim K}{xx^\top}
=C_K+m_Km_K^\top.
\]
At each temperature the algorithm estimates both the centered covariance and the
mean, then uses the covariance estimate to define the next affine coordinates.
The key comparison between consecutive temperatures is a finite-order
change-of-measure estimate.
If $s=(1-\delta)t$ and
$\mu_u(x)\propto e^{-u\|x\|_2^2}1_K(x)$, then
\[
D_2(\mu_s\Vert\mu_t) =\Lambda(t)+\Lambda(2s-t)-2\Lambda(s),
\qquad
\Lambda(u)=\log\int_K e^{-u\|x\|_2^2}\,dx.
\]
Using the variance bound $\Lambda''(u)=\V{X\sim\mu_u}{\|X\|_2^2}\le O(n/u^2)$
gives
$D_2(\mu_s\Vert\mu_t)=O(n\delta^2)$.
Choosing $\delta=\Theta(n^{-1/4})$ makes this divergence
$O(n^{1/2})$.
The full-density primitive can tolerate this finite-order warmness because,
after truncating to the Paouris event for the old rounded law, the complement
contributes negligibly and the second moment of the new target remains $O(n)$ in
the old coordinates.
This is the source of the $n^{1/4}$ factor in the
$O(n^{1/4}\log(nR/r))$ cooling-stage bound.

Thus the algorithm makes $O(n^{1/4}\log(nR/r))$ calls to the near-cubic
full-density rounding primitive.
Each call returns an approximate mean and covariance in the current coordinates;
the next coordinates are obtained by subtracting the empirical mean and applying
the inverse square root of the empirical covariance.
At the endpoint the Gaussian tilt is weak enough that the cooled law is
comparable to the uniform law on $K$, and one final call estimates the desired
covariance and mean.
Combining them as $\widehat M=\widehat C+\widehat m\widehat m^\top$ gives
\cref{thm:general-upper-bound}.

Compared with the previous $\O(n^{3.5})$ rounding framework, this is both faster
and structurally simpler: the proof follows an explicit Gaussian-cooling path
rather than tracking intersections with a growing ellipsoid.
The symmetric cubic algorithm in the next section goes further by replacing $tI$
with a matrix precision, so the algorithm can continue cooling the
well-conditioned directions while separately regularizing the troublesome ones.

For symmetric bodies this scalar algorithm is still valid, with the
simplification that all means are zero.
The next section describes a genuinely symmetric improvement: because
multiplying the density by any even quadratic weight preserves its center, the
algorithm can choose an anisotropic matrix precision and remove the extra
$n^{1/4}$ factor.

\subsection{Cubic Symmetric Algorithm}
\label{subsec:overview-cubic-symmetric}

For a symmetric body, the uncentered second moment and covariance coincide
throughout the cooling path.
More importantly, symmetry lets us replace the scalar temperature by a
matrix-valued quadratic penalty without having to track a moving mean.
For a positive semidefinite matrix $H$, define
\[
\mu_H(x)\propto e^{-x^\top Hx}1_K(x),
\qquad
C_H=\E{x\sim\mu_H}{xx^\top}.
\]
The matrix $H$ specifies the strength of the quadratic penalty in each
direction; it is also called the precision matrix of the Gaussian factor.
We start from $H_0=r^{-2}I$ and drive
$H$ toward
$0$.
At every successful iteration the algorithm maintains an estimate
$W\approx C_H$.
In the normalized coordinates $z=W^{-1/2}x$, the current covariance is
\[
 W^{-1/2}C_HW^{-1/2}\approx I.
\]

The update is the main new feature.
Set
\[
H^+=(1-\delta)H+
\eta_{\rm ridge}\frac{\Tr{W\,H}}{n}W^{-1}.
\]
If $B=W^{1/2}HW^{1/2}$, then the quadratic-penalty matrix in the normalized
coordinates changes to
\[
 (1-\delta)B+
 \eta_{\rm ridge}\frac{\Tr{B}}{n}I.
\]
Thus every eigenvalue is decreased by the scalar cooling step and receives the
same additive ridge.
The ridge has the largest relative effect on directions in which the current
quadratic penalty is smallest.
The interpolation estimates in the proof show that the next transformed law
still has $O(n)$ second moment and a fixed-radius ball in an
$O(n)$-height sublevel set.
These are the hypotheses of the near-cubic covariance estimator used at each
step.

The analysis tracks
\[
D_{\rm sym}(H)=\log\frac{\Vol(K)}{\int_K e^{-x^\top Hx}\,dx},
\qquad
q_{\rm sym}(H)=\Tr{C_H\,H}.
\]
If $X$ is uniform on
$K$, then
\[
 e^{-D_{\rm sym}(H)}=\E{X\sim K}{e^{-X^\top H X}}.
\]
Thus $D_{\rm sym}$ measures the reduction in the normalizing integral caused by
the quadratic weight, while
$q_{\rm sym}(H)=\E{X\sim\mu_H}{X^\top HX}$ is the expected quadratic penalty
under
$\mu_H$.
With $L=\Theta(1+\log(2+nR/r))$,
$\delta=L^{-2}$, and
$\eta_{\rm ridge}=L^{-3}$, the proof gives
\[
D_{\rm sym}(H^+)\leq D_{\rm sym}(H)-\frac{\delta}{2}q_{\rm sym}(H),
\qquad
q_{\rm sym}(H)
\leq D_{\rm sym}(H)
\leq O(L)q_{\rm sym}(H).
\]
Hence $D_{\rm sym}$ contracts by a factor
$1-\Theta(L^{-3})$ at each successful update, and
$O(L^4)$ updates suffice.
Once $D_{\rm sym}$ is a small constant, a terminal comparison shows that the
uniform measure on
$K$ is well-rounded in the current coordinates, so one final near-cubic
covariance call estimates
$C_K$.

This gives the cubic symmetric algorithm of \cref{thm:symmetric-upper-bound}.
Sampling-based membership-oracle volume algorithms already spend $\O(n^2)$
membership queries per well-rounded sample and use
$O(n)$ samples for stable cooling ratios.
Thus we do not expect a sampling-based approach to beat cubic query complexity
except by changing the model or avoiding the bottleneck invoked by sampling
procedures.

\subsection{Counterexamples to Previously Stated Conjectures}

We also include several refutations of previously conjectured directions for
further improving these algorithms.
Specifically, we refute the possibility of constant-size scalar Gaussian-cooling
steps from the algorithm in this paper, and the possibility of large
ellipsoid-growth steps from an already rounded intersection referenced in
\cite{JLLV26:journal}.

For the scalar-cooling obstruction, an astute reader may observe that if
\[
 \left\|A_t^{-1/2}\frac{dA_t}{dt}A_t^{-1/2}\right\|_{\op}\le \O(1/t),
\]
then one could simply decrease $t\to(1-1/\O(1))t$ per iteration, and obtain a
cubic scalar-cooling algorithm even for asymmetric bodies.
For a compact convex body $K\subset\R^n$ and positive definite
$Q$, write
\[
  \mu_{K,Q}(\mathrm dx)
  =\frac{1}{Z_{K,Q}}\mathbf 1_K(x)
    \exp\!\left(-\frac12 x^{\mathsf T}Q^{-1}x\right)\mathrm dx,
\]
and say that $\mu_{K,Q}$ is isotropic if it has mean zero and covariance
$I_n$.
For $0\leq\delta\leq1$, define
\[
  R_n(\delta)
  =\sup\Bigl\{
    \operatorname{Tr}\operatorname{Cov}(\mu_{K,(1+\delta)Q}):
    K\subset\R^n\text{ a convex body},\ Q\succ0,\
    \mu_{K,Q}\text{ isotropic}
  \Bigr\}.
\]

The next theorem concerns the sharp trace growth under anisotropic Gaussian
warming.

\begin{theorem}
\label{thm:sharp-trace-lower}
There is a universal constant $c>0$ such that, for every
$n\geq1$ and
$0\leq\delta\leq1$,
\[
  R_n(\delta)\geq c\bigl(n+n^2\delta^4\bigr).
\]
For this lower bound, $K$ may be taken to be a convex polytope.
\end{theorem}

The proof is given in \cref{sec:sharp}.
Together with \cref{thm:coordinate-moment-counterexample}, this rules out the
shortcut above.
In particular, \cref{thm:sharp-trace-lower} constructs restricted Gaussians for
which the trace after a small anisotropic warming is $\Omega(n+n^2\delta^4)$.
Thus a constant warming step can make the next target have trace $\Omega(n^2)$
in the old coordinates, far outside the well-rounded regime needed by the
near-cubic primitive.
The scale $\delta=\Theta(n^{-1/4})$ is exactly where this obstruction stops
being superlinear in
$n$.
Similarly, \cref{thm:coordinate-moment-counterexample} gives a symmetric
log-concave counterexample to the coordinate-moment simplification behind a
constant-size derivative bound, so this shortcut fails even before one has to
handle moving means.

For the ellipsoid-growth obstruction, the framework of \cite{JLLV26:journal}
suggests another possible simplification: if $K\cap E$ is already approximately
isotropic, then modestly enlarging
$E$ might keep
$K\cap(1+\delta)E$ well-rounded.
Such a statement would allow the older ball/ellipsoid-growth method to take much
larger growth steps.
The following theorem shows that this route fails as well.

The next theorem concerns the counterexample to Conjecture~3.5 of
\cite{JLLV26:journal}.

\begin{theorem}
\label{thm:overview-counterexample-conj35}
There are absolute constants $c,C>0$ such that, for every sufficiently large
$n$, one can construct a convex body
$K\subset\R^n$ and an ellipsoid
$E\subset\R^n$ such that
$K\cap E$ has mean zero and covariance
$I_n$, but for every
$0<\delta\le1$ satisfying
$(n-1)\delta\ge C$,
\[
\Tr{\operatorname{Cov}(K\cap(1+\delta)E)}
\ge c(n-1)^2\delta^2.
\]
\end{theorem}

\Cref{thm:overview-counterexample-conj35} is proved in
\cref{sec:counter_ellipsoid_growth_asym}; the quantitative form used there is
\cref{app:counterexample:theorem:quantitative}.
This limitation still leaves open the possibility of a future asymmetric volume
algorithm that avoids full second-moment estimation.

We now state the two conjectures refuted in \cref{sec:counter_coordinate_moment,sec:counter_ellipsoid_growth_asym}.
Recall the notation from \cref{sec:prelim}: a convex body $K$ has
uniform law $\pi_K$, and its mean, covariance, and uncentered second
moment are $m_K$,
$C_K$, and
$M_K$ respectively.
When $K=-K$, we have
$m_K=0$ and
$C_K=M_K$.

We first give a counterexample to a coordinate-moment conjecture that would
make the scalar Gaussian-cooling derivative uniformly bounded.
If true, the resulting bound would permit constant-size temperature decreases in
the scalar cooling method of \cref{sec:asymmetric}, removing its $n^{1/4}$ stage
overhead and yielding a cubic asymmetric rounding algorithm up to logarithmic
factors.

The next conjecture concerns the coordinate-moment conjecture.

\begin{conjecture}\label{conj:first}
\label{conj:gaussian-cooling-coordinate-moment}
Let $f:\R^n\to[0,+\infty)$ be an isotropic log-concave function.
Then
\[
  \left|\E{x\sim f}{x_1^2\bigl(\norm{x}^2-n\bigr)}\right|=O(1),
\]
where $x_1=\langle x,e_1\rangle$ and the implicit constant is absolute.
\end{conjecture}

The second obstruction is the ellipsoid-growth statement that would let the
previous method of \cite{JLLV26:journal} take substantially larger growth steps.
If true, it would allow that method to replace its $1+\Theta(n^{-1/2})$ updates
by constant-factor updates, reducing the number of stages from
$\widetilde O(n^{1/2}\log(R/r))$ to
$\widetilde O(\log(R/r))$ and giving a cubic asymmetric rounding algorithm.

The next conjecture concerns the ellipsoid growth, Conjecture 3.5 from
\cite{JLLV26:journal}.

\begin{conjecture}
\label{conj:jllv-ellipsoid-growth}
Suppose that for a convex body $K$ and an ellipsoid
$E$, the set
$K\cap E$ is
$2$-isotropic.
Then $K\cap 2E$ is well-rounded.
\end{conjecture}

For this body-level conjecture, we say that a convex body $L$ is
$2$-isotropic if
\[
\frac12 I_n
\preceq M_L
\preceq 2I_n,
\]
and well-rounded if $\Tr{C_L}=\widetilde O(n)$ (the
$\widetilde O$ hides polylogarithmic factors) and
$B(0,1)\subseteq L$.
This is the body-level specialization of the well-roundedness terminology used
in the main text.

The next theorem concerns the failure of the coordinate-moment conjecture.

\begin{theorem}
\label{thm:coordinate-moment-counterexample}
There is a universal constant $c>0$ such that, for every sufficiently large
$n$, there is an even isotropic log-concave density
$f$ on
$\R^n$ satisfying
\[
  \left|\E{x\sim f}{x_1^2\bigl(\norm{x}^2-n\bigr)}\right|
  \geq cn^{1/2}.
\]
\end{theorem}

The conjecture counterexamples indicate that the scalar Gaussian-cooling route
is unlikely to be pushed to constant-size steps using
the moment simplifications discussed in \cref{sec:step-size-remarks},
and the previous ellipsoid-growth route cannot rely on Conjecture~3.5 of \cite{JLLV26:journal}.
We therefore view cubic complexity as the natural limit for these second-moment
methods in the membership-oracle model, while a different idea might still
achieve cubic asymmetric volume estimation.

% ===== begin 4_asymmetric_kimi.tex =====
\section{General-Body Algorithm: Scalar Cooling}
\label{sec:asymmetric}

In this section we prove the $\O(n^{13/4})$ rounding algorithm for a general
convex body, specifically~\cref{thm:general-upper-bound}.

The proof replaces the ellipsoid-growth chain of previous rounding algorithms by
an explicit Gaussian-cooling path.
Rather than intersecting the body with a growing ball or ellipsoid and tracking
how the covariance changes under that geometric operation, we keep the body
fixed and change the density: with
\[
\Phi_t(x)=V(x)+t\|x-x_0\|_2^2,
\qquad
\pi_t\propto e^{-\Phi_t},
\]
we cool the temperature $t$ from
$r^{-2}$ down to
$\Theta(R^{-2})$ by multiplicative steps
$t\leftarrow(1-\delta)t$ with
$\delta=\Theta(n^{-1/4})$.
At each temperature a near-cubic full-density rounding primitive,
\algname{FullDensityRound}, estimates the mean and covariance of the current
cooled law; the estimate defines the affine coordinates in which the next,
slightly warmer, law is again well-rounded.
The step size is dictated by a finite-order change-of-measure estimate: the
order-two R\'enyi divergence between consecutive cooled laws is
$O(n\delta^2)=O(n^{1/2})$, exactly the scale at which the Paouris tail bound
transfers the
$O(n)$ second-moment guarantee from one temperature to the next.
After $O(n^{1/4}\log(2+R/r))$ stages the Gaussian tilt is negligible, one final
call estimates the mean and covariance of the target law, and the uncentered
second moment is reconstructed as
$\widehat M=\widehat C+\widehat m\widehat m^\top$.

The section is organized as follows.
\cref{sec:grounded-core} constructs the full-density rounding interface
(\cref{lem:fdwr}) from the near-cubic well-rounded rounding black box recorded
in \cref{prelim:thm:wellr_iso} and the moment-estimation black box of
\cref{thm:moment-estimation}; the short reductions of \cref{prop:log-concave-toolkit}
parts (iii)--(v) promised in \cref{sec:prelim} are also given there.
\cref{sec:cooling-estimates} proves the cooling estimates: the nonsmooth
extended Brascamp--Lieb inequality (\cref{lem:bl}), the quadratic energy and
variance bound (\cref{lem:energy}), the one-step R\'enyi cost, and the
well-roundedness induction at the next temperature.
\cref{sec:endpoint-reconstruction} compares the cooled moments to the target
moments and reconstructs the raw second moment.
\cref{sec:general-algorithm} gives the algorithm
(\cref{fig:asymmetric-second-moment}), the well-roundedness induction, and the
query accounting, and proves \cref{thm:general-upper-bound}.
\cref{sec:step-size-remarks} records why the scalar route cannot be pushed to
constant-size steps.

\subsection{The Full-Density Rounding Interface}
\label{sec:grounded-core}

\paragraph{The two black boxes.}
This subsection rests on the two imported statements of \cref{sec:prelim}:
the isotropic rounding of a grounded law (\cref{prelim:thm:wellr_iso},
\cite[Lemma~4.11]{KV25}) and the moment estimation of a well-rounded law with
bounded covariance (\cref{thm:moment-estimation}, assembled from
\cite[Theorem~3.16, Lemma~2.14, Lemma~2.9, and Lemma~4.6]{KV25}).
Both are stated with the exact constants of \cite[Assumption~4.4]{KV25}: unit
inradius, level-set height
$10n$, and (for the rounding) support equal to that level set.
An input of \algname{FullDensityRound} is only
$(\rho,H,B)$-well-rounded (\cref{def:rounded}) for universal
$\rho,H,B$, so the proof of \cref{lem:fdwr} first brings it to these exact
constants by an explicit reduction (\cref{lem:restriction}): a rescaling, a
lift by independent Gaussian coordinates to a dimension
$N=O(n)$ in which the level-set height
$Hn$ is below
$N$, and a restriction to a level set of mass at least
$1-10^{-5}$.
Constants hidden by asymptotic notation below may depend on $(\rho,H,B)$, but
not on
$n$ or on a caller-supplied failure probability.
Throughout, $n_0\geq4$ is a universal integer fixed in the proof of
\cref{lem:fdwr}.

\paragraph{Reductions of the toolkit parts (iii)--(v).}
We discharge the promise made in \cref{sec:prelim} that the one-dimensional
moment and level-set estimates of \cite{LV07:journal} imply parts (iii)--(v) of
\cref{prop:log-concave-toolkit}.
\begin{proof}[Proof of \cref{prop:log-concave-toolkit} parts
\textnormal{(iii)--(v)}] Part (iii).
Let $\ell$ be an affine functional and write
$L=|\ell(Y)|$.
If $\ell$ is constant, the moment claims are immediate; otherwise its law
is a one-dimensional log-concave density.
Applying \cref{lem:reverse_holder} with $k=4$ and with
$k=2$ to the one-dimensional log-concave law of
$\ell(Y)$ (\cref{lem:prekopa_leindler}) gives
\[
\E{}{L^4}
\leq 8^4\bigl(\E{}{L}\bigr)^4
\leq 8^4\bigl(\E{}{L^2}\bigr)^2,
\]
which is part (iii) with $C_4=8^4$.

Part (v), first claim.
The same bound with $k=2$ reads
$(\E{}{L^2})^{1/2}\leq 4\E{}{L}$.

Part (v), second claim, and the support ball.
For an isotropic log-concave law, each one-dimensional marginal has density
at most $1$, and each halfspace containing the mean has probability at
least $1/e$ \cite[Lemmas~5.5(a) and~5.12]{LV07:journal}.
Consequently, every supporting hyperplane of its support is at distance
at least $1/e$ from its mean.
Applying the inverse whitening map gives the radius
$1/(es^{1/2})$ when $s^{-1}I\preceq C\preceq sI$;
see also \cite[Lemma~A.10]{KV25}.

Part (iv).
The isotropic density lower bound of
\cite[Theorem~5.14(a)]{LV07:journal} gives universal
$\rho_0,H_0>0$ such that every isotropic
log-concave density
$e^{-W}$ obeys
$e^{-W(x)}\geq e^{-H_0n}e^{-W(0)}$ on
$B(0,2\rho_0)$.
Suppose first that $m=0$ and
$\frac12I\preceq C\preceq2I$.
Pushing to isotropic coordinates preserves the conclusion up to a factor $2$, so
$W(x)\leq W(0)+H_0n$ on
$B(0,2^{1/2}\,\rho_0)$; the barycenter inequality (part (i)) gives
$W(0)\leq\inf W+n$, and hence
$B(0,\rho_L)\subseteq\{W\leq\inf W+H_Ln\}$ with
$\rho_L=2^{1/2}\,\rho_0$ and
$H_L=H_0+1$.
If instead $\|m\|\leq\rho_L/2$, apply the centered statement to
$\widetilde W(x)=W(x+m)$, whose law is centered and still has covariance in
$[\frac12I,2I]$: then
$B(m,\rho_L)\subseteq\{W\leq\inf W+H_Ln\}$, and
$B(0,\rho_L/2)\subseteq B(m,\rho_L)$.
\end{proof}

The rest of this subsection constructs the rounding interface.
The first step is a level-set mass estimate, which will let us truncate a
log-concave law to a sublevel set of its potential while barely changing it.

The next lemma concerns the level-set mass.

\begin{lemma}
\label{lem:level-mass}
Let $W:\R^n\to\R\cup\{+\infty\}$ be proper, lower-semicontinuous, and convex
with
$0<\int e^{-W}<\infty$, let
$a=\inf W$, and let
$\mu\propto e^{-W}$ be the associated probability law.
Then for every $\bar H\geq2$ and every
$s\geq2$,
\[
 \P{X\sim\mu}{W(X)\geq a+(\bar H+s)n}\leq\frac{4}{\bar H}\,e^{-sn/2}.
\]
\end{lemma}

\begin{proof}
Replacing $W$ by
$W(\cdot-x^*)-a$ for a minimizer
$x^*$ (which exists: integrability forces
$W$ to tend to
$+\infty$ along every ray, so its sublevel sets are bounded), assume
$a=0$ and
$W(0)=0$.
For $t\geq0$ let
$v(t)=\Vol(\{W\leq t\})$.
For $0<s<t$ and
$x\in\{W\leq t\}$, convexity gives
$W(\frac stx)\leq\frac stW(x)+(1-\frac st)W(0) \leq s$, that is,
$\frac st\{W\leq t\}\subseteq\{W\leq s\}$, and therefore
\[
 v(s)\geq\Bigl(\frac st\Bigr)^{\!n}v(t)
 \qquad(0<s<t).
\]
Layer cake gives $Z=\int e^{-W}\,dx=\int_0^\infty e^{-t}v(t)\,dt$, hence with
$T_0=\bar Hn$,
\[
Z
\geq v(T_0)\,T_0^{-n}\!\!\int_0^{T_0}\!\!e^{-t}t^n\,dt
\geq v(T_0)\,T_0^{-n}e^{-T_0}\frac{T_0^{\,n+1}}{n+1},
\]
while for $T>T_0$, using
$(1+u/T)^n\leq e^{nu/T}$,
\[
\int_T^\infty e^{-t}v(t)\,dt
\leq v(T_0)\,T_0^{-n}e^{-T}T^n\!\int_0^\infty\!\!e^{-u+nu/T}du
= v(T_0)\,T_0^{-n}e^{-T}T^n\Bigl(1-\frac nT\Bigr)^{-1}.
\]
Taking $T=(\bar H+s)n$ with
$s\geq2$ and dividing, and using
$(1+s/\bar H)^n\leq e^{sn/\bar H}$,
$\bar H\geq2$, and
$(1-n/T)^{-1}\leq2$,
\[
\mu\{W\geq T\}
\leq \frac{(n+1)\,e^{\bar Hn}}{\bar H n\,(1-n/T)} \Bigl(1+\frac{s}{\bar H}\Bigr)^{\!n}e^{-(\bar H+s)n}
\leq \frac{2(n+1)}{\bar H n}\,e^{-sn(1-1/\bar H)}
\leq \frac{4}{\bar H}\,e^{-sn/2},
\]
since $s\geq2$.
\end{proof}

The next lemma concerns the lifted grounded restriction.

\begin{lemma}
\label{lem:restriction}
Let $\mu\propto e^{-U}$ on
$\R^n$ be
$(\rho,H,B)$-well-rounded around
$z$, let
$k=\lceil Hn\rceil+n_0$ and
$N=n+k$, and define on
$\R^N=\R^n\times\R^k$
\[
 \bar U(x,y)=U(z+\rho x)+\tfrac12\|y\|^2,
 \qquad
 \bar\mu\propto e^{-\bar U},
\]
and
\[
 K=\{(x,y):\bar U(x,y)\leq\bar U(0)+9N\},
 \qquad
 \nu=\bar\mu(\,\cdot\mid K).
\]
Then $\bar\mu=\mu_\rho\otimes\gamma_k$, where
$\mu_\rho$ is the law of
$\rho^{-1}(Y-z)$ for
$Y\sim\mu$ and
$\gamma_k$ is the standard Gaussian law on
$\R^k$, and the following hold with
$\bar B=B/\rho^2+1$.
\begin{enumerate}
\item \label{itm:restriction:1} $B(0,1)\subseteq\{\bar U\leq\inf\bar U+N\}\subseteq K$ and
  $\E{X\sim\bar\mu}{\|X\|^2}\leq\bar BN$.
\item \label{itm:restriction:2} $\bar\mu(K)\geq q_0:=1-10^{-5}$ and
  $d\nu/d\bar\mu\leq q_0^{-1}$ on
  $K$.
\item \label{itm:restriction:3} $\inf_K\bar U=\inf\bar U$ and
  $\sup_K\bar U-\inf\bar U\leq10N$.
\item \label{itm:restriction:4} With $U'=\bar U+I_K-\bar U(0)+N$, the law
  $\nu\propto e^{-U'}$ is grounded around
  $0$ with respect to
  $U'$ (\cref{def:rounded}), and
  $\E{X\sim\nu}{\|X\|^2}\leq2\bar BN$.
\item \label{itm:restriction:5} One evaluation of
  $U'$ costs one evaluation of
  $U$ after
  $U(z)$ has been computed once.
\end{enumerate}
\end{lemma}

\begin{proof}
The product form is the change of variables $x=\rho^{-1}(y_1-z)$ on the first
factor, the Gaussian factor being
$e^{-\|y\|^2/2}$.
Write $\bar a=\inf\bar U=\inf U$, the infimum of the Gaussian part being
$0$.
For $\|(x,y)\|\leq1$ we have
$z+\rho x\in B(z,\rho)$ and
$\frac12\|y\|^2\leq\frac12$, so
$\bar U(x,y)\leq\inf U+Hn+\frac12\leq\bar a+N$, because
$N\geq Hn+n_0\geq Hn+1$; this is the first inclusion of
Item~\ref{itm:restriction:1}.
Since $\bar a\leq\bar U(0)$, the set
$\{\bar U\leq\bar a+N\}$ lies in
$\{\bar U\leq\bar U(0)+9N\}=K$.
The second moment is
$\E{\bar\mu}{\|X\|^2}=\E{\mu}{\|\rho^{-1}(Y-z)\|^2}+k\leq Bn/\rho^2+k \leq\bar BN$.

For Item~\ref{itm:restriction:2}, $K\supseteq\{\bar U\leq\bar a+9N\}$, so
\cref{lem:level-mass} in dimension
$N$ with
$\bar H=2$ and
$s=7$ gives
\[
 \bar\mu(K^c)
 \leq\P{X\sim\bar\mu}{\bar U(X)\geq\bar a+9N}
 \leq2e^{-7N/2}
 \leq2e^{-14}
 \leq10^{-5},
\]
using $N\geq n_0\geq4$; and
$d\nu/d\bar\mu=\ind_K/\bar\mu(K)\leq q_0^{-1}$.

For Item~\ref{itm:restriction:3}, $K$ contains
$\{\bar U\leq\bar a+9N\}$, which contains points of
$\bar U$-value arbitrarily close to
$\bar a$, so
$\inf_K\bar U=\bar a$; and
$\sup_K\bar U\leq\bar U(0)+9N\leq\bar a+10N$ because
$0\in B(0,1)\subseteq\{\bar U\leq\bar a+N\}$.

For Item~\ref{itm:restriction:4}, $\nu\propto e^{-U'}$ since
$U'$ differs from
$\bar U+I_K$ by a constant.
We have $\min U'=\inf_K\bar U-\bar U(0)+N=\bar a-\bar U(0)+N\geq0$ by
Item~\ref{itm:restriction:3}, and on
$K$,
$U'\leq\bar U(0)+9N-\bar U(0)+N=10N$.
Hence $\operatorname{supp}\nu=K\subseteq\{U'\leq\min U'+10N\}$, while
$\{U'\leq\min U'+10N\}\subseteq\{U'<\infty\}=K$; so the support of
$\nu$ equals the level set
$\{U'\leq\inf U'+10N\}$, which contains
$B(0,1)$ by Item~\ref{itm:restriction:1}.
Finally
$\E{\nu}{\|X\|^2}=\bar\mu(K)^{-1}\E{\bar\mu}{\|X\|^2\ind_K} \leq\bar BN/q_0\leq2\bar BN$.

For Item~\ref{itm:restriction:5}, membership in $K$ is decided by comparing
$\bar U(x,y)$ with
$\bar U(0)+9N=U(z)+9N$, and
$\bar U(x,y)$ is one evaluation of
$U$ at
$z+\rho x$ plus explicit arithmetic.
\end{proof}

The next lemma transfers rounding information from a large-mass restriction back
to the full law; it is the content of Steps~2 and~3 of \cite[Section~4.3]{KV25},
recorded here in the form we use.

The next lemma concerns the restricted-to-full transfer.

\begin{lemma}
\label{lem:transfer}
Let $\lambda$ be a log-concave probability law on
$\R^n$, let
$S\subseteq\R^n$ be measurable with
$\lambda(S)\geq1-\eps$ for some
$\eps\in(0,10^{-4}]$, and let
$\nu=\lambda(\,\cdot\mid S)$.
Then for every unit vector $v$ and every center
$c\in\R^n$,
\[
 \E{\lambda}{|v^\top(X-c)|}
 \leq1.05\,\E{\nu}{|v^\top(X-c)|},
\]
and, writing the variances as $\sigma_\lambda^2=\V{\lambda}{v^\top X}$ and
$\sigma_\nu^2=\V{\nu}{v^\top X}$,
\[
(1-\eps)\,\sigma_\nu^2
\leq\sigma_\lambda^2
\leq18.1\,\sigma_\nu^2 ,
\qquad
\bigl(v^\top(m_\lambda-m_\nu)\bigr)^2 \leq\frac{\eps}{1-2\eps}\,\sigma_\lambda^2 .
\]
In particular, if an affine map $F$ makes
$F_\#\nu$ satisfy
$C_{F_\#\nu}\approx_{1.01}I$ and
$\|m_{F_\#\nu}\|\leq0.01$, then
$0.99\,I\preceq C_{F_\#\lambda}\preceq18.5\,I$ and
$\|m_{F_\#\lambda}\|\leq0.06$.
\end{lemma}

\begin{proof}
Set $f(x)=|v^\top(x-c)|$ and note the
$\chi^2$-warmness of a conditional law: since
$d\nu=\ind_S\,d\lambda/\lambda(S)$,
\[
\chi^2(\nu\Vert\lambda)
=\E{\lambda}{\Bigl(\frac{d\nu}{d\lambda}-1\Bigr)^{\!2}}
=\frac{1}{\lambda(S)}-1
\leq\frac{\eps}{1-\eps}.
\]
Cauchy--Schwarz and \cref{prop:log-concave-toolkit}(v) applied to the affine
functional $v^\top(x-c)$ under the log-concave law
$\lambda$ give
\[
\bigl|\E{\lambda}{f}-\E{\nu}{f}\bigr|
\leq\E{\lambda}{f^2}^{1/2}(\chi^2(\nu\Vert\lambda))^{1/2}
\leq4\,\E{\lambda}{f}(\eps/(1-\eps))^{1/2},
\]
so, since $4(\eps/(1-\eps))^{1/2}\leq0.05$ for
$\eps\leq10^{-4}$,
\[
\E{\lambda}{f}
\leq\bigl(1-4(\eps/(1-\eps))^{1/2}\bigr)^{-1}\E{\nu}{f}
\leq1.05\,\E{\nu}{f}.
\]
Squaring after another application of \cref{prop:log-concave-toolkit}(v),
\[
\E{\lambda}{f^2}
\leq16\bigl(\E{\lambda}{f}\bigr)^2
\leq16(1.05)^2\bigl(\E{\nu}{f}\bigr)^2
\leq18\,\E{\nu}{f^2},
\]
since $16(1.05)^2=17.64\leq18$.
For the centered bounds, let $g(x)=v^\top(x-m_\nu)$ and
$\Delta=(v^\top(m_\lambda-m_\nu))^2$.
The same Cauchy--Schwarz step with $g$ in place of
$f$ gives
$\Delta=|\E{\lambda}{g}-\E{\nu}{g}|^2 \leq\E{\lambda}{g^2}\,\chi^2(\nu\Vert\lambda) \leq(\sigma_\lambda^2+\Delta)\,\eps/(1-\eps)$,
hence
$\Delta\leq\frac{\eps}{1-2\eps}\sigma_\lambda^2$.
The lower bound follows from
$\E{\lambda}{(X-m_\lambda)(X-m_\lambda)^\top} \succeq\lambda(S)\,\E{\nu}{(X-m_\lambda)(X-m_\lambda)^\top} \succeq(1-\eps)C_\nu$.
For the upper bound, take $c=m_\lambda$ above:
$\sigma_\lambda^2=\E{\lambda}{f^2}\leq18\,\E{\nu}{f^2}$ with
$\E{\nu}{f^2}=\sigma_\nu^2+\Delta \leq\sigma_\nu^2+\frac{\eps}{1-2\eps}\sigma_\lambda^2$;
rearranging,
$\sigma_\lambda^2 \leq\frac{18}{1-18\eps/(1-2\eps)}\,\sigma_\nu^2 \leq18.1\,\sigma_\nu^2$,
since
$18\eps/(1-2\eps)\leq1.9\cdot10^{-3}$.

For the final claim, in $F$-coordinates
$\sigma_\nu^2\in[1/1.01,1.01]$ for every unit
$v$. Hence
\[
\sigma_\lambda^2\in[0.9999/1.01,\,18.1\cdot1.01]
\subseteq[0.99,18.5].
\]
Moreover,
\[
\|m_{F_\#\lambda}-m_{F_\#\nu}\|^2
=\sup_v\bigl(v^\top(m_{F_\#\lambda}-m_{F_\#\nu})\bigr)^2
\leq\frac{\eps}{1-2\eps}\sup_v\sigma_\lambda^2
\leq\frac{10^{-4}}{1-2\cdot10^{-4}}\cdot18.5
\leq0.002,
\]
so $\|m_{F_\#\lambda}\|\leq\|m_{F_\#\nu}\|+0.002^{1/2} \leq0.01+0.05\leq0.06$.
\end{proof}

The next lemma concerns the robust selection.

\begin{lemma}
\label{lem:selection}
Fix $\xi\in(0,1/50)$ and
$p\in(0,1/10)$, and let
$L=\lceil512\log(2/p)\rceil$.
Fix $m\in\R^n$ and $C\succ0$.
Let $L$ independent random outputs each be either a failure symbol or a pair
$(\widehat m_i,\widehat C_i)$ with $\widehat C_i\succ0$.
Suppose each output is \emph{$\xi$-accurate} for $(m,C)$ with probability
at least $17/32$, meaning
$(1-\xi)C\preceq\widehat C_i\preceq(1+\xi)C$ and
$\|C^{-1/2}(\widehat m_i-m)\|\leq\xi$.
Failure symbols have no partners. For nonfailure outputs, declare $i\sim j$ when
\[
(1-3\xi)\widehat C_i
\preceq\widehat C_j
\preceq(1+3\xi)\widehat C_i,
\qquad
\|\widehat C_i^{-1/2}(\widehat m_i-\widehat m_j)\|\leq3\xi .
\]
Select a nonfailure index $i^*$ with the most partners, returning a failure
symbol if none exists. With probability at least $1-p/2$, the selected pair
is $5\xi$-accurate.
\end{lemma}

\begin{proof}
Accurate pairs are mutually compatible: for accurate $i,j$,
$\widehat C_i\preceq(1+\xi)C\preceq\frac{1+\xi}{1-\xi}\widehat C_j \preceq(1+3\xi)\widehat C_j$
since
$\xi\leq1/50$, and symmetrically; and
\[
\|\widehat C_i^{-1/2}(\widehat m_i-\widehat m_j)\|
\leq\|\widehat C_i^{-1/2}C^{1/2}\|\, \bigl(\|C^{-1/2}(\widehat m_i-m)\| +\|C^{-1/2}(m-\widehat m_j)\|\bigr)
\leq2(1-\xi)^{-1/2}\,\xi
\leq3\xi ,
\]
using
$\|\widehat C_i^{-1/2}C^{1/2}\|^2 =\|\widehat C_i^{-1/2}C\widehat C_i^{-1/2}\|_{\op}\leq(1-\xi)^{-1}$.
By Hoeffding's inequality, the number of accurate indices exceeds $L/2$
with probability at least
$1-e^{-L/512}\geq1-p/2$.
On that event, every accurate index has more than $L/2$ partners,
so $i^*$ also has more than $L/2$ partners.
Since fewer than $L/2$ indices are inaccurate,
$i^*$ has an accurate partner $j$.
Compatibility then gives
\[
\widehat C_{i^*}
\preceq\frac1{1-3\xi}\widehat C_j
\preceq\frac{1+\xi}{1-3\xi}C
\preceq(1+5\xi)C,
\qquad
\widehat C_{i^*}
\succeq\frac1{1+3\xi}\widehat C_j
\succeq\frac{1-\xi}{1+3\xi}C
\succeq(1-5\xi)C,
\]
and
\[
\|C^{-1/2}(\widehat m_{i^*}-m)\|
\leq\|C^{-1/2}\widehat C_{i^*}^{1/2}\|\, \|\widehat C_{i^*}^{-1/2}(\widehat m_{i^*}-\widehat m_j)\| +\|C^{-1/2}(\widehat m_j-m)\|
\leq(1+5\xi)^{1/2}\cdot3\xi+\xi
\leq5\xi .\qedhere
\]
\end{proof}

\begin{figure}[H]
\begin{algbox}
\textbf{\uline{\algname{FullDensityRound}}}$(U,z,\rho,H,B,\eps,p)$

\textbf{Input:} evaluation oracle for a proper lower-semicontinuous convex
potential $U$; a known point
$z\in\R^n$; universal constants
$\rho,H,B>0$; a fixed accuracy
$\eps\in(0,1/10)$; a failure tolerance
$p\in(0,1/10)$.

\textbf{Output:} a failure symbol $\bot$, or a pair
$(\widehat m,\widehat C)$ with
$\widehat C\succ0$.

\textbf{Assumptions:} $\mu\propto e^{-U}$ is a probability law with
$B(z,\rho)\subseteq\{U\leq\inf U+Hn\}$ and
$\E{Y\sim\mu}{\|Y-z\|^2}\leq Bn$.

\begin{enumerate}[leftmargin=*,itemsep=2pt]
\item Set $L=\lceil512\log(2/p)\rceil$,
  $\xi=\eps/5$,
  $k=\lceil Hn\rceil+n_0$,
  $N=n+k$, and compute
  $U(z)$.
  Implement the grounded potential $U'$ on
  $\R^N$ of \cref{lem:restriction}; each evaluation costs one evaluation of
  $U$.
  Let $Q_{\rm iso}(N)=A_CN^3\log^{a_C}(2N)$ be the expected-query bound of
  \cref{prelim:thm:wellr_iso} for
  $C^2=2\bar B=2B/\rho^2+2$, and let
  $Q_{\rm est}(N)=168AN^3\log^a(128N)$ be that of \cref{thm:moment-estimation}
  for
  $c_+=167$, accuracy
  $\xi$,
  $R^2=167N$, and
  $\eta=1/64$.
\item For each $i\in[L]$, with fresh randomness:
  \begin{enumerate}[itemsep=1pt]
\item Run the rounding algorithm of \cref{prelim:thm:wellr_iso} on
  $U'$ with the known point
  $0$ and the known constant
  $C$, capped at
  $16Q_{\rm iso}(N)$ evaluations.
  If it fails or times out, mark trial $i$ failed
  and continue with the next trial.
  Otherwise it returns an invertible affine map $F_i(x)=A_i(x-\bar x_i)$ on
  $\R^N$.
\item Run the estimator of \cref{thm:moment-estimation} on the potential
  $W_i(u)=\bar U(F_i^{-1}(u/3))$ on
  $\R^N$ with the known point
  $0$,
  $R^2=167N$, accuracy
  $\xi$, and
  $\eta=1/64$, capped at
  $16Q_{\rm est}(N)$ evaluations.
  If it fails or times out, mark trial $i$ failed
  and continue with the next trial.
  Otherwise it returns $(\widehat m_i',\widehat C_i')$; set
  $\widehat m_i=\widehat m_i'/3$ and
  $\widehat C_i=\widehat C_i'/9$, pull back to the lifted coordinates by
  $\widetilde m_i=F_i^{-1}\widehat m_i$ and
  $\widetilde C_i=A_i^{-1}\widehat C_iA_i^{-\top}$, and let
  $\widetilde m_i^x\in\R^n$ and
  $\widetilde C_i^x\in\R^{n\times n}$ be their blocks on the first
  $n$ coordinates.
  If $\widetilde C_i^x$ is not positive definite, mark trial $i$ failed.
  \end{enumerate}
\item Among the successful trials, select an index $i^*$ whose pair
  $(\widetilde m_i^x,\widetilde C_i^x)$ has the most
  $\sim$-partners in the sense of \cref{lem:selection}; if no trial succeeded,
  return
  $\bot$.
\item Return $\widehat m=z+\rho\,\widetilde m_{i^*}^x$ and
  $\widehat C=\rho^2\widetilde C_{i^*}^x$.
\end{enumerate}
\end{algbox}
\caption{The full-density rounding primitive used by the scalar
cooling algorithm.}
\label{fig:full-density-round}
\end{figure}

The next lemma concerns the full-density well-rounded rounding.

\begin{lemma}
\label{lem:fdwr}
Fix universal constants $\rho,H,B>0$ and a fixed accuracy
$\eps\in(0,1/10)$.
The algorithm \algname{FullDensityRound} (\cref{fig:full-density-round}), given
exact-real evaluation access to a proper lower-semicontinuous convex $U$ whose
law
$\mu\propto e^{-U}$ is a probability, a known point
$z$ with
\[
 B(z,\rho)\subseteq\{U\leq\inf U+Hn\},
 \qquad
 \E{Y\sim\mu}{\|Y-z\|^2}\leq Bn ,
\]
and a failure tolerance $p\in(0,1/10)$, returns either a failure symbol or a
pair
$(\widehat m,\widehat C)$ with
$\widehat C\succ0$.
Conditional on every realized pre-call history for which these input promises
hold (\cref{def:adaptive-oracle-guarantee}), with probability at least $1-p$ it
returns a pair satisfying
\[
(1-\eps)C_\mu
\preceq\widehat C
\preceq(1+\eps)C_\mu,
\qquad
\|C_\mu^{-1/2}(\widehat m-m_\mu)\|\leq\eps .
\]
The number of evaluation queries is at most $O(n^3\log(n/p)^{O(1)})$ on every
execution.
\end{lemma}

\begin{proof}
We analyze \cref{fig:full-density-round}.
Fix the universal integer $n_0\geq4$ so large that
$c_{\rm iso}n_0^{-1/2}\leq1/16$ and
$c_{\rm est}/n_0\leq1/16$, with
$c_{\rm iso}$ and
$c_{\rm est}$ from \cref{prelim:thm:wellr_iso,thm:moment-estimation}.
Then $N=n+\lceil Hn\rceil+n_0$ satisfies
$n_0\leq N\leq(H+1)n+n_0+1=O(n)$.
Step~1 is \cref{lem:restriction}: the lifted law
$\bar\mu=\mu_\rho\otimes\gamma_k$ on
$\R^N$ has the restriction
$\nu=\bar\mu(\,\cdot\mid K)$ with
$\bar\mu(K)\geq1-10^{-5}$, and
$\nu\propto e^{-U'}$ is grounded around
$0$ with
$\E{\nu}{\|X\|^2}\leq2\bar BN=C^2N$ for the universal constant
$C^2=2B/\rho^2+2$.
Fix a trial $i$ and condition on the pre-call history and on the other trials;
the hypotheses on
$\mu$ are deterministic properties of the input, so the two black boxes apply
with their stated guarantees over the fresh randomness of trial
$i$.

Step 2(a).
\cref{prelim:thm:wellr_iso} applies to $\nu$ with the known point
$0$ and the known constant
$C$, and returns, on an event
$E_1$ of probability at least
$1-c_{\rm iso}N^{-1/2}\geq1-1/16$, an invertible affine map
$F_i(x)=A_i(x-\bar x_i)$ with
$C_{F_{i\#}\nu}\approx_{1.01}I_N$ and
$\|m_{F_{i\#}\nu}\|\leq0.01$.
On $E_1$ the expected number of queries is at most
$Q_{\rm iso}(N)$, so by Markov's inequality the run capped at
$16Q_{\rm iso}(N)$ queries times out on
$E_1$ with probability at most
$1/16$.
Hence Step~2(a) succeeds with probability at least $1-1/8$.

Step 2(b).
Apply \cref{lem:transfer} with $\lambda=\bar\mu$,
$S=K$, and
$\eps=10^{-5}$: the law
$\pi_i=F_{i\#}\bar\mu$ satisfies
$0.99I\preceq C_{\pi_i}\preceq18.5I$ and
$\|m_{\pi_i}\|\leq0.06$.
Its potential is $\bar U\circ F_i^{-1}$ up to an additive constant, and
$\operatorname{supp}F_{i\#}\nu=F_i(K)$.
Since $F_{i\#}\nu$ is
$1.01$-isotropic with mean of norm at most
$0.01$, \cref{prop:log-concave-toolkit}(v) applied to its centering gives
\[
 B(0,0.35)
 \subseteq B\!\left(m_{F_{i\#}\nu},\frac{1}{e\,1.01^{1/2}}\right)
 \subseteq F_i(K)
 \subseteq\{\bar U\circ F_i^{-1}\leq\inf(\bar U\circ F_i^{-1})+10N\},
\]
the last inclusion by Item~\ref{itm:restriction:3}.
The potential $W_i(u)=\bar U(F_i^{-1}(u/3))$ is, up to an additive constant,
that of the law
$\pi_i'=(3\,\cdot)_\#\pi_i$, for which
\[
 B(0,1)\subseteq B(0,1.05)\subseteq\{W_i\leq\inf W_i+10N\},
 \qquad
 8.9I\preceq C_{\pi_i'}=9C_{\pi_i}\preceq166.5I,
\]
and
\[
 \E{X\sim\pi_i'}{\|X\|^2}
 =\Tr C_{\pi_i'}+\|m_{\pi_i'}\|^2
 \leq166.5N+0.18^2
 \leq167N .
\]
Thus \cref{thm:moment-estimation} applies to $\pi_i'$ with the known point
$0$,
$R^2=167N$,
$c_+=167$, accuracy
$\xi$, and
$\eta=1/64$: on an event
$E_2$ of probability at least
$1-1/32-c_{\rm est}/N\geq1-3/32$ it returns
$(\widehat m_i',\widehat C_i')$ with
$(1-\xi)C_{\pi_i'}\preceq\widehat C_i'\preceq(1+\xi)C_{\pi_i'}$ and
$\|\widehat m_i'-m_{\pi_i'}\|\leq\xi$, and on
$E_2$ the run capped at
$16Q_{\rm est}(N)$ queries times out with probability at most
$1/16$ by Markov's inequality.
Each evaluation of $W_i$ is one evaluation of
$U$ by Item~\ref{itm:restriction:5}, the Gaussian part being explicit.
On $E_2$ without time-out, since
$C_{\pi_i'}\succeq8.9I$,
\[
 \|C_{\pi_i'}^{-1/2}(\widehat m_i'-m_{\pi_i'})\|
 \leq8.9^{-1/2}\|\widehat m_i'-m_{\pi_i'}\|
 \leq\xi ,
\]
so $(\widehat m_i',\widehat C_i')$ is
$\xi$-accurate for
$(m_{\pi_i'},C_{\pi_i'})$ in the sense of \cref{lem:selection}, and dividing
by
$3$ and by
$9$ gives a
$\xi$-accurate pair
$(\widehat m_i,\widehat C_i)$ for
$(m_{\pi_i},C_{\pi_i})$, both accuracy notions being invariant under
scaling.
Altogether trial $i$ produces a
$\xi$-accurate pair for
$(m_{\pi_i},C_{\pi_i})$ with probability at least
\[
 1-\frac1{16}-\frac1{16}-\frac{3}{32}-\frac1{16}
 =\frac{23}{32}
 \geq\frac{17}{32}
\]
over its fresh randomness, conditional on the pre-call history and on the
other trials.

Step 3.
Since $C_{\pi_i}=A_iC_{\bar\mu}A_i^\top$ and
$m_{\pi_i}=F_i(m_{\bar\mu})$, congruence invariance (compare
\cref{def:affine-pushforward}) makes the pulled-back pair
$(\widetilde m_i,\widetilde C_i)$ of a successful trial
$\xi$-accurate for
$(m_{\bar\mu},C_{\bar\mu})$: indeed
$\widetilde C_i=A_i^{-1}\widehat C_iA_i^{-\top}$ shares the Loewner window of
$\widehat C_i$ relative to
$C_{\pi_i}$, and
$\|C_{\bar\mu}^{-1/2}(\widetilde m_i-m_{\bar\mu})\| =\|C_{\pi_i}^{-1/2}(\widehat m_i-m_{\pi_i})\|$.
As $\bar\mu=\mu_\rho\otimes\gamma_k$, we have
$m_{\bar\mu}=(m_{\mu_\rho},0)$ and
$C_{\bar\mu}=\operatorname{diag}(C_{\mu_\rho},I_k)$.
Compressing $(1-\xi)C_{\bar\mu}\preceq\widetilde C_i\preceq(1+\xi)C_{\bar\mu}$
to the first
$n$ coordinates gives
$(1-\xi)C_{\mu_\rho}\preceq\widetilde C_i^x\preceq(1+\xi)C_{\mu_\rho}$, and
since $C_{\bar\mu}^{-1/2}$ is block diagonal,
\[
 \|C_{\mu_\rho}^{-1/2}(\widetilde m_i^x-m_{\mu_\rho})\|
 \leq\|C_{\bar\mu}^{-1/2}(\widetilde m_i-m_{\bar\mu})\|
 \leq\xi ,
\]
so $(\widetilde m_i^x,\widetilde C_i^x)$ is
$\xi$-accurate for
$(m_{\mu_\rho},C_{\mu_\rho})$.
Every accurate covariance block is positive definite, so the additional
check rejects no accurate trial. Treat failed trials as failure symbols
in \cref{lem:selection}.
The trials use fresh randomness, so conditional on the pre-call history their
outcomes are independent, each
$\xi$-accurate with probability at least
$17/32$.
\cref{lem:selection} with $L=\lceil512\log(2/p)\rceil$ then gives, except with
probability
$p/2\leq p$, a selected pair that is
$5\xi=\eps$-accurate for
$(m_{\mu_\rho},C_{\mu_\rho})$.

Step 4.
Since $C_{\mu_\rho}=\rho^{-2}C_\mu$ and
$m_\mu=z+\rho\,m_{\mu_\rho}$, the returned pair
$(z+\rho\,\widetilde m_{i^*}^x,\,\rho^2\widetilde C_{i^*}^x)$ is
$\eps$-accurate for
$(m_\mu,C_\mu)$ on the selection event, again by congruence invariance.

Query cap.
Each trial costs at most $16Q_{\rm iso}(N)+16Q_{\rm est}(N)=O(n^3\log^{O(1)}n)$
evaluations deterministically, since
$N=O(n)$ and the constants
$A_C,a_C,A,a$ depend only on
$(\rho,H,B,\eps)$; with
$L=\lceil512\log(2/p)\rceil$ trials the total is
$O(n^3\log(n/p)^{O(1)})$ on every execution.
Each query is one evaluation of $U$ plus arithmetic.
\end{proof}

\subsection{Cooling Estimates}
\label{sec:cooling-estimates}

For the rest of this section, fix a grounded input $(V,x_0,r,R)$ in the sense of
\cref{def:grounded-input}:
$V$ is proper, lower-semicontinuous, and convex with
$0<\int e^{-V}<\infty$,
$B(x_0,r)\subseteq\{V\leq\inf V+h_0n\}$, and
$\E{X\sim\pi}{\|X-x_0\|^2}\leq R^2$ for
$\pi\propto e^{-V}$.
For $t>0$ set
\[
\Phi_t(x)=V(x)+t\|x-x_0\|^2,
\qquad
Z_t=\int e^{-\Phi_t(x)}\,dx,
\qquad
\pi_t=Z_t^{-1}e^{-\Phi_t},
\]
and write $m_t=m_{\pi_t}$,
$C_t=C_{\pi_t}$, and
$a_t=\inf\Phi_t$.
Note $Z_t\leq e^{-\inf V}\int e^{-t\|x-x_0\|^2}dx<\infty$ and
$Z_t\geq\Vol(B(x_0,r))e^{-\inf V-h_0n-tr^2}>0$, so
$\pi_t$ is a well-defined log-concave probability law.
Evaluating $\Phi_t\circ F^{-1}$ for a known invertible affine
$F$ costs one evaluation of
$V$ plus arithmetic.

The first estimate is the nonsmooth form of Brascamp--Lieb promised in
\cref{sec:prelim}; it is used with $G=I_K$ in \cref{sec:symmetric}.

The next lemma concerns the extended Brascamp--Lieb consequence.

\begin{lemma}
\label{lem:bl}
Let $G:\R^n\to\R\cup\{+\infty\}$ be proper, lower-semicontinuous, and convex,
let
$t>0$, and suppose
$\nu_t\propto e^{-G-t\|x\|^2}$ is a probability law.
Then
\[
C_{\nu_t}\preceq\frac1{2t}I,
\qquad
\V{\nu_t}{g}\leq\frac1{2t}\,\E{\nu_t}{\|\nabla g\|^2}
\]
for every locally Lipschitz $g$ with
$\|\nabla g(x)\|\leq c_g(1+\|x\|)$ almost everywhere.
\end{lemma}

\begin{proof}
Suppose first that $G$ is twice continuously differentiable.
Then $W(x)=G(x)+t\|x\|^2$ satisfies
$\nabla^2W\succeq2tI$, so \cref{prop:log-concave-toolkit}(vi) applies and gives
both claims:
$C_{\nu_t}\preceq(2t)^{-1}I$ and
$\V{\nu_t}{g}\leq(2t)^{-1}\E{\nu_t}{\|\nabla g\|^2}$.

For the general case, fix an affine minorant $\ell$ of
$G$ (which exists by properness) and approximate in two steps.
First, the Moreau envelope $G_\eta(x)=\inf_y\{G(y)+\frac1{2\eta}\|x-y\|^2\}$ is
convex, continuously differentiable with Lipschitz gradient, increases to
$G$ pointwise as
$\eta\downarrow0$, and satisfies
$G_\eta\geq\ell-C_\eta$ for some
$C_\eta<\infty$.
Second, the Gaussian mollification $G_{\eta,\sigma}=G_\eta*\rho_\sigma$ is
smooth and convex, and by Jensen
$G_{\eta,\sigma}(x)=\E{}{G_\eta(x+\sigma\xi)}\geq G_\eta(x)$.
The smooth case applies to $G_{\eta,\sigma}$, giving both inequalities for
$\nu_{\eta,\sigma}\propto e^{-G_{\eta,\sigma}-t\|x\|^2}$.
The densities are dominated uniformly in $(\eta,\sigma)$: with
$\tilde\ell(x)=\ell(x)-C_1$,
\[
e^{-G_{\eta,\sigma}(x)-t\|x\|^2}
\leq e^{-G_\eta(x)-t\|x\|^2}
\leq e^{-\tilde\ell(x)-t\|x\|^2},
\]
whose right side is integrable together with $\|x\|^2$ times itself.
As $\sigma\downarrow0$,
$G_{\eta,\sigma}\to G_\eta$ locally uniformly, so dominated convergence
identifies
$\nu_{\eta,\sigma}\to\nu_\eta$ and convergence of the covariance and of both
sides of the variance inequality (using the growth bound on
$\|\nabla g\|$); then
$\eta\downarrow0$ gives
$G_\eta\uparrow G$ and the same passage yields both claims for
$\nu_t$.
\end{proof}

The next lemma concerns the quadratic energy and variance.

\begin{lemma}
\label{lem:energy}
There is a constant $K_E=K_E(h_0)$ such that, for every
$t\in(0,r^{-2}]$,
\[
t\,\E{\pi_t}{\|X-x_0\|^2}\leq K_En,
\qquad
\V{\pi_t}{\|X-x_0\|^2}\leq\frac{2K_En}{t^2}.
\]
\end{lemma}

\begin{proof}
Covariance is unaffected by translation, so applying \cref{lem:bl} with
$G(x)=V(x+x_0)$ gives
$\Tr{C_t}\leq\frac{n}{2t}$.
For the mean, the barycenter inequality (\cref{prop:log-concave-toolkit}(i))
gives $\Phi_t(m_t)\leq a_t+n$, and any
$y\in B(x_0,r)$ with
$V(y)\leq\inf V+h_0n$ gives
$a_t\leq\Phi_t(y)\leq\inf V+h_0n+tr^2$; hence
\[
t\|m_t-x_0\|^2
\leq\Phi_t(m_t)-\inf V
\leq a_t-\inf V+n
\leq(h_0+1)n+tr^2
\leq(h_0+2)n ,
\]
using $tr^2\leq1$.
Summing,
$t\,\E{\pi_t}{\|X-x_0\|^2} =t\Tr{C_t}+t\|m_t-x_0\|^2 \leq\frac n2+(h_0+2)n=:K_En$.
The variance claim is the same application of \cref{lem:bl} to
$g(x)=\|x-x_0\|^2$ with
$\|\nabla g(x)\|=2\|x-x_0\|$:
$\V{\pi_t}{g}\leq\frac1{2t}\E{\pi_t}{4\|X-x_0\|^2} \leq\frac{2K_En}{t^2}$.
\end{proof}

The next lemma concerns the one-step R\'enyi cost.

\begin{lemma}
\label{lem:renyi}
For $0<\delta\leq1/4$,
$s=(1-\delta)t$, and
$t\leq r^{-2}$,
\[
D_2(\pi_s\Vert\pi_t)\leq K_Dn\delta^2,
\qquad
K_D=8K_E .
\]
\end{lemma}

\begin{proof}
Since $2\Phi_s-\Phi_t=\Phi_{2s-t}$ with
$2s-t=(1-2\delta)t>0$,
\[
\mathcal R_2(\pi_s\Vert\pi_t)
=\int\frac{\pi_s^2}{\pi_t}
=\frac{Z_tZ_{2s-t}}{Z_s^2},
\qquad
\text{so with }\Lambda(u)=\log Z_u,
\quad
D_2(\pi_s\Vert\pi_t)=\Lambda(t)+\Lambda(2s-t)-2\Lambda(s),
\]
a centered second difference.
Differentiating twice, $\Lambda''(u)=\V{\pi_u}{\|X-x_0\|^2}\leq2K_En/u^2$ by
\cref{lem:energy}.
Taylor's theorem with $h=\delta t$, noting
$s-h=(1-2\delta)t\geq t/2$, gives
\[
D_2(\pi_s\Vert\pi_t)
=\int_0^h(h-w)\bigl[\Lambda''(s+w)+\Lambda''(s-w)\bigr]\,dw
\leq h^2\cdot\frac{2K_En}{(t/2)^2}
=8K_En\delta^2.\qedhere
\]
\end{proof}

The next lemma concerns the R\'enyi-to-trace transfer.

\begin{lemma}
\label{lem:renyi-trace}
There are universal constants $\gamma\in(0,1/4]$ and
$K_T=K_T(h_0)$ such that, with
$\delta=\gamma n^{-1/4}$ and
$s=(1-\delta)t$, the following holds for every
$t\leq r^{-2}$.
If $F$ is an invertible affine map for which
$Z=F(X)$,
$X\sim\pi_t$, satisfies
\[
\tfrac12I
\preceq C_{F_\#\pi_t}
\preceq2I,
\qquad
\|\E{F_\#\pi_t}{Z}\|\leq\rho_L/2,
\]
then $\E{\pi_s}{\|F(X)\|^2}\leq K_Tn$.
\end{lemma}

\begin{proof}
Write $L=\frac{d\pi_s}{d\pi_t}(X)$ and split
$\E{\pi_s}{\|Z\|^2}=\E{\pi_t}{\|Z\|^2L}$ on the Paouris event
$A=\{\|Z\|\leq uc_Pn^{1/2}\}$, with
$c_P$ from \cref{prop:log-concave-toolkit}(ii) and
$u\geq2$ universal, chosen below.
On $A$ the integrand is at most
$u^2c_P^2nL$, contributing at most
$u^2c_P^2n$.
Off $A$, H\"older with exponents
$(4,2,4)$ gives
\[
 \E{\pi_t}{\|Z\|^2L\ind_{A^c}}
 \leq\bigl(\E{\pi_t}{\|Z\|^8}\bigr)^{1/4}
     \bigl(\E{\pi_t}{L^2}\bigr)^{1/2}
     \P{\pi_t}{A^c}^{1/4}.
\]
The three factors are:
$(\E{\pi_t}{\|Z\|^8})^{1/4} \leq16^2\,\E{\pi_t}{\|Z\|^2}\leq768n$ by
\cref{lem:reverse_holder} with
$k=8$ and
$\E{\pi_t}{\|Z\|^2}=\Tr{C_{F_\#\pi_t}}+\|m_{F_\#\pi_t}\|^2 \leq2n+\rho_L^2/4\leq3n$;
$(\E{\pi_t}{L^2})^{1/2} =\mathcal R_2(\pi_s\Vert\pi_t)^{1/2} \leq e^{K_D\gamma^2\sqrt n/2}$
by \cref{lem:renyi}; and
$\P{\pi_t}{A^c}^{1/4}\leq e^{-u\sqrt n/(8c_P)}$ by
\cref{prop:log-concave-toolkit}(ii) applied to
$Z-m_{F_\#\pi_t}$, with the mean shift
$\rho_L/2$ absorbed by the choice
$u\geq2$.
Choosing $u=\max(2,8c_PK_D\gamma^2)$ makes the combined exponent at most
$-K_D\gamma^2n^{1/2}/4$ in either case, so
\[
\E{\pi_t}{\|Z\|^2L\ind_{A^c}}
\leq768n\cdot e^{-K_D\gamma^2\sqrt n/4}
\leq768n.
\]
In total,
\[
\E{\pi_s}{\|Z\|^2}\leq (u^2c_P^2+768)n=K_Tn,
\qquad K_T:=u^2c_P^2+768.
\]
\end{proof}

The next lemma concerns the moving minimum.

\begin{lemma}
\label{lem:moving-min}
For $s=(1-\delta)t$ and
$t\leq r^{-2}$,
\[
0
\leq a_t-a_s
\leq\frac{\delta}{1-\delta}(h_0+1)n .
\]
\end{lemma}

\begin{proof}
Since $\Phi_t-\Phi_s=(t-s)\|x-x_0\|^2\geq0$, we have
$a_t\geq a_s$.
Let $y$ minimize
$\Phi_s$.
As in \cref{lem:energy}, $a_s\leq\inf V+h_0n+sr^2$, so
$s\|y-x_0\|^2\leq\Phi_s(y)-\inf V\leq(h_0+1)n$.
Then
\[
a_t
\leq\Phi_t(y)
=a_s+(t-s)\|y-x_0\|^2
\leq a_s+\frac{t-s}{s}(h_0+1)n
=a_s+\frac{\delta}{1-\delta}(h_0+1)n .\qedhere
\]
\end{proof}

We now fix the universal constants used by the cooling induction:
\[
\rho_*=\rho_L/2,
\qquad
H_*=H_L+h_0+1,
\qquad
B_*=\max\bigl(K_T,\,K_E/\rho_*^2\bigr),
\qquad
\eps_{\rm cool}=\min\bigl(1/500,\,\rho_L/3\bigr),
\]
where $\rho_L,H_L$ are from \cref{prop:log-concave-toolkit}(iv) and
$K_E,K_T$ from \cref{lem:energy,lem:renyi-trace}.

The next lemma concerns the well-roundedness at the next temperature.

\begin{lemma}
\label{lem:next-temp}
Let $s=(1-\delta)t$ with
$\delta=\gamma n^{-1/4}$ and
$t\leq r^{-2}$.
Suppose a call to \algname{FullDensityRound} on $\pi_t$ (in affine coordinates
in which
$\pi_t$ is
$(\rho_*,H_*,B_*)$-well-rounded around
$0$) succeeded and returned
$(\widehat m,\widehat C)$, and set
$F(x)=\widehat C^{-1/2}(x-\widehat m)$.
Then in these coordinates $\pi_s$ is
$(\rho_*,H_*,B_*)$-well-rounded around
$0$:
\[
 B(0,\rho_*)\subseteq\{\Phi_s\circ F^{-1}\leq a_s+H_*n\},
 \qquad
 \E{\pi_s}{\|F(X)\|^2}\leq B_*n .
\]
\end{lemma}

\begin{proof}
On the success event,
$(1-\eps_{\rm cool})C_t\preceq\widehat C\preceq(1+\eps_{\rm cool})C_t$ and
$\|C_t^{-1/2}(\widehat m-m_t)\|\leq\eps_{\rm cool}$.
Hence
\[
 C_{F_\#\pi_t}
 =\widehat C^{-1/2}C_t\widehat C^{-1/2}
 \in\bigl[(1+\eps_{\rm cool})^{-1},(1-\eps_{\rm cool})^{-1}\bigr]I
 \subseteq\bigl[\tfrac12I,2I\bigr],
\]
and
$\|m_{F_\#\pi_t}\| =\|\widehat C^{-1/2}(m_t-\widehat m)\| \leq(1-\eps_{\rm cool})^{-1/2}\,\eps_{\rm cool} \leq\rho_L/2$.
\cref{prop:log-concave-toolkit}(iv) (the relaxed-centering form) now gives
$B(0,\rho_L/2)\subseteq\{\Phi_t\circ F^{-1}\leq a_t+H_Ln\}$.
For $z$ in that ball, since
$\Phi_s\leq\Phi_t$ pointwise and by \cref{lem:moving-min},
\[
\Phi_s(F^{-1}z)
\leq a_t+H_Ln
\leq a_s+H_Ln+\frac{\delta}{1-\delta}(h_0+1)n
\leq a_s+(H_L+h_0+1)n
=a_s+H_*n ,
\]
using $\frac{\delta}{1-\delta}\leq1$.
The second-moment bound is \cref{lem:renyi-trace}:
$\E{\pi_s}{\|F(X)\|^2}\leq K_Tn\leq B_*n$.
\end{proof}

The next lemma concerns the initialization.

\begin{lemma}
\label{lem:init}
For $t_0=r^{-2}$ and
$F_{\rm base}(x)=(x-x_0)/(r\rho_*)$, the law
$\pi_{t_0}$ in
$F_{\rm base}$-coordinates is
$(\rho_*,H_*,B_*)$-well-rounded around
$0$.
\end{lemma}

\begin{proof}
By \cref{lem:energy} with $t_0r^2=1$,
\[
\E{\pi_{t_0}}{\|F_{\rm base}(X)\|^2}
=\frac{\E{\pi_{t_0}}{\|X-x_0\|^2}}{r^2\rho_*^2}
\leq\frac{K_En}{t_0r^2\rho_*^2}
=\frac{K_En}{\rho_*^2}
\leq B_*n .
\]
For $\|z\|\leq\rho_*$
($\leq1$ without loss of generality, since the constants may be shrunk), the
point
$F_{\rm base}^{-1}z=x_0+r\rho_*z$ lies in
$B(x_0,r)$, so
$V(F_{\rm base}^{-1}z)\leq\inf V+h_0n$ and
\[
\Phi_{t_0}(F_{\rm base}^{-1}z)
\leq\inf V+h_0n+t_0r^2\rho_*^2\|z\|^2
\leq\inf V+(h_0+1)n
\leq a_{t_0}+H_*n ,
\]
using $a_{t_0}\geq\inf V$ and
$H_*\geq h_0+1$.
\end{proof}

\subsection{Endpoint Comparison and Raw Reconstruction}
\label{sec:endpoint-reconstruction}

At the end of the cooling path the Gaussian tilt is weak enough that the cooled
moments and the target moments are comparable, uniformly in the centering.
Write
\[
M(b)=\E{\pi}{(X-b)(X-b)^\top},
\qquad
M_t(b)=\E{\pi_t}{(X-b)(X-b)^\top}.
\]

The next lemma concerns the Reverse-H\"older endpoint comparison.

\begin{lemma}
\label{lem:endpoint}
There is a universal $K_H$ such that, whenever
$\eta=tR^2\leq1/2$,
\[
(1-K_H\eta)\,M(b)
\preceq M_t(b)
\preceq(1+2\eta)\,M(b)
\qquad
\text{for every }b\in\R^n .
\]
\end{lemma}

\begin{proof}
The likelihood ratio is $\frac{d\pi_t}{d\pi}(x)=\frac{1}{c_t}e^{-t\|x-x_0\|^2}$
with
$c_t=Z_t/Z_V=\E{\pi}{e^{-t\|X-x_0\|^2}}$.
Jensen and the grounded moment bound give
$c_t\geq e^{-t\E{\pi}{\|X-x_0\|^2}}\geq e^{-\eta}$, while
$c_t\leq1$.
The upper bound is immediate:
$M_t(b)=\frac1{c_t}\E{\pi}{(X-b)(X-b)^\top e^{-t\|X-x_0\|^2}} \preceq c_t^{-1}M(b)\preceq e^{\eta}M(b)\preceq(1+2\eta)M(b)$
for
$\eta\leq1/2$.
For the lower bound, since $c_t\leq1$ and
$1-e^{-u}\leq u$, for every unit vector
$v$,
\[
v^\top\bigl(M(b)-M_t(b)\bigr)v
\leq\E{\pi}{\langle X-b,v\rangle^2 \bigl(1-e^{-t\|X-x_0\|^2}\bigr)}
\leq t\,\E{\pi}{\langle X-b,v\rangle^2\|X-x_0\|^2}.
\]
Cauchy--Schwarz, \cref{prop:log-concave-toolkit}(iii) applied to the affine
functional $\langle X-b,v\rangle$, and \cref{lem:reverse_holder} with
$k=4$ give
\[
t\bigl(\E{\pi}{\langle X-b,v\rangle^4}\bigr)^{1/2} \bigl(\E{\pi}{\|X-x_0\|^4}\bigr)^{1/2}
\leq tC_4^{1/2}\,\bigl(v^\top M(b)v\bigr)\cdot64R^2
=64C_4^{1/2}\,\eta\,v^\top M(b)v ,
\]
so $M_t(b)\succeq(1-K_H\eta)M(b)$ with
$K_H=64C_4^{1/2}$.
\end{proof}

The next lemma states that exact target is well-rounded.

\begin{lemma}
\label{lem:exact-wr}
Set $\eta_0=\min(1/2,1/(4K_H))$.
Suppose $tR^2\leq\eta_0$ and that, in some invertible affine coordinates
$F$, the cooled law
$\pi_t$ is
$(\rho_*,H_*,B_*)$-well-rounded around
$0$.
Then the target law $\pi$ in the same coordinates is
$(\rho_*,H_*',2B_*)$-well-rounded around
$0$ with
$H_*'=H_*+h_0+1$:
\[
 B(0,\rho_*)\subseteq\{V\circ F^{-1}\leq\inf(V\circ F^{-1})+H_*'n\},
 \qquad
 \E{\pi}{\|F(X)\|^2}\leq2B_*n .
\]
\end{lemma}

\begin{proof}
Write $F(x)=A(x-b)$.
By \cref{lem:endpoint} with $K_H\eta_0\leq1/4$,
$M(b)\preceq(1-K_H\eta_0)^{-1}M_t(b) \preceq\frac43M_t(b)$, so
\[
\E{\pi}{\|F(X)\|^2}
=\Tr{A^\top A\,M(b)}
\leq\tfrac43\Tr{A^\top A\,M_t(b)}
=\tfrac43\,\E{\pi_t}{\|F(X)\|^2}
\leq2B_*n .
\]
For the ball, since $\Phi_t=V+t\|x-x_0\|^2\geq V$, on
$B(0,\rho_*)$ we have
$V\circ F^{-1}\leq\Phi_t\circ F^{-1}\leq a_t+H_*n$, and
$a_t\leq\inf V+h_0n+tr^2\leq\inf V+(h_0+1)n$ because
$tr^2\leq t_0r^2=1$; hence
$V\circ F^{-1}\leq\inf V+(H_*+h_0+1)n$ on
$B(0,\rho_*)$.
\end{proof}

The next lemma concerns the raw second-moment reconstruction.

\begin{lemma}
\label{lem:raw}
Let $C\succ0$,
$m\in\R^n$, and fix any center
$c\in\R^n$; write
$M(c)=C+(m-c)(m-c)^\top$.
If $(1-\eps)C\preceq\widehat C\preceq(1+\eps)C$ and
$\|C^{-1/2}(\widehat m-m)\|\leq\zeta$, then
$\widehat M(c)=\widehat C+(\widehat m-c)(\widehat m-c)^\top$ is positive
semidefinite and
\[
-(\eps+\zeta+\zeta^2)\,M(c)
\preceq\widehat M(c)-M(c)
\preceq(\eps+\zeta+\zeta^2)\,M(c).
\]
For $\eps=\zeta=1/500$ this gives
$0.99\widehat M(c)\preceq M(c)\preceq1.01\widehat M(c)$.
\end{lemma}

\begin{proof}
Let $d=\widehat m-m$.
Then $\widehat M(c)-M(c)=(\widehat C-C)+dd^\top+(m-c)d^\top+d(m-c)^\top$.
The first term obeys $-\eps C\preceq\widehat C-C\preceq\eps C\preceq\eps M(c)$;
the second
$0\preceq dd^\top\preceq\zeta^2C\preceq\zeta^2M(c)$ from
$\|C^{-1/2}d\|\leq\zeta$; and for the cross term, every unit vector
$v$ obeys
\[
\begin{aligned}
\bigl|v^\top\bigl((m-c)d^\top+d(m-c)^\top\bigr)v\bigr|
&=2|\langle v,m-c\rangle||\langle v,d\rangle|\\
&\leq2(v^\top(m-c)(m-c)^\top v)^{1/2}(\zeta^2v^\top Cv)^{1/2}\\
&\leq\zeta\,v^\top M(c)v,
\end{aligned}
\]
by $2ab\leq a^2+b^2$.
Summing the three gives the Loewner claim.
With $\alpha=\eps+\zeta+\zeta^2=0.004004<\frac1{249}$,
the bounds
$(1-\alpha)M(c)\preceq\widehat M(c)\preceq(1+\alpha)M(c)$ imply
\[
(1+\alpha)^{-1}\widehat M(c)
\preceq M(c)
\preceq(1-\alpha)^{-1}\widehat M(c).
\]
Thus $0.99\widehat M(c)\preceq M(c)\preceq1.01\widehat M(c)$.
\end{proof}

\subsection{Algorithm, Induction, and Complexity}
\label{sec:general-algorithm}

\begin{figure}[H]
\begin{algbox}
\textbf{\uline{\algname{AsymmetricSecondMoment}}}$(V,x_0,r,R,\beta)$

\textbf{Input:} evaluation oracle for a proper lower-semicontinuous convex
potential $V$; the rounding data
$x_0\in\R^n$ and
$r,R>0$; a failure parameter
$\beta\in(0,1)$.

\textbf{Output:} a positive semidefinite $\widehat M$ and an invertible affine
map
$\mathcal T$.

\textbf{Assumptions:} $(V,x_0,r,R)$ is a grounded input in the sense of
\cref{def:grounded-input}.

\begin{enumerate}[leftmargin=*,itemsep=2pt]
\item Set $t_0=r^{-2}$,
  $\delta=\gamma n^{-1/4}$ with
  $\gamma$ from \cref{lem:renyi-trace}, and
  $t_j=(1-\delta)^jt_0$.
  Let $N$ be the first index with
  $t_NR^2\leq\eta_0$
  ($\eta_0$ from \cref{lem:exact-wr}), and set
  $p=\min(1/20,\beta/(N+1))$.
\item Set $F\leftarrow F_{\rm base}$ of \cref{lem:init}.
\item For $j=0,1,\dots,N-1$:
  \begin{enumerate}[itemsep=1pt]
\item Call \algname{FullDensityRound} (\cref{lem:fdwr}) on the potential
  $\Phi_{t_j}\circ F^{-1}$ with center
  $0$, constants
  $(\rho_*,H_*,B_*)$, accuracy
  $\eps_{\rm cool}$, and failure tolerance
  $p$; on failure, abort and return failure.
\item With the returned $(\widehat m_j,\widehat C_j)$, set
  $F\leftarrow\widehat C_j^{-1/2}(F(\cdot)-\widehat m_j)$.
  \end{enumerate}
\item Call \algname{FullDensityRound} on the exact target $V\circ F^{-1}$ with
  center
  $0$, constants
  $(\rho_*,H_*',2B_*)$, accuracy
  $\eps_{\rm cool}$, and failure tolerance
  $p$, obtaining
  $(\widehat m,\widehat C)$.
\item Write $F(x)=A(x-b)$.
  Return
  \[
\widehat M =A^{-1}\bigl(\widehat C+(\widehat m+Ab)(\widehat m+Ab)^\top\bigr) A^{-\top},
\qquad
\mathcal T(x)=\widehat C^{-1/2}(F(x)-\widehat m).
  \]
\end{enumerate}
\end{algbox}
\caption{The general-body scalar-cooling rounding algorithm.}
\label{fig:asymmetric-second-moment}
\end{figure}

The next lemma concerns the cooling induction.

\begin{lemma}
\label{lem:induction}
Conditional on all earlier calls succeeding, every
\algname{FullDensityRound} call made by
\algname{AsymmetricSecondMoment} (\cref{fig:asymmetric-second-moment})
satisfies the input promises of \cref{lem:fdwr}, and each evaluation of the
transformed potential costs one evaluation of $V$.
\end{lemma}

\begin{proof}
At $j=0$, the promises hold by \cref{lem:init}.
If the call at stage $j\geq0$ succeeds, then \cref{lem:next-temp} applied in the
current coordinates shows that
$\pi_{t_{j+1}}$ is
$(\rho_*,H_*,B_*)$-well-rounded around
$0$ in the updated coordinates, which are exactly the input promises for the
call at stage
$j+1$.
For the final call, \cref{lem:exact-wr} applies since $t_NR^2\leq\eta_0$, so the
exact target
$V\circ F^{-1}$'s law is
$(\rho_*,H_*',2B_*)$-well-rounded around
$0$.
Evaluation of $\Phi_t\circ F^{-1}(z)=V(F^{-1}z)+t\|F^{-1}z-x_0\|^2$ (or of
$V\circ F^{-1}$) is one evaluation of
$V$ plus arithmetic, and the conditional guarantee of \cref{lem:fdwr}
(\cref{def:adaptive-oracle-guarantee}) is exactly what is invoked at each stage.
\end{proof}

The next lemma concerns the radius comparison.

\begin{lemma}
\label{lem:Rr}
Every grounded input obeys $R\geq c_{h_0}r$ with
$c_{h_0}=\frac14e^{-h_0}$.
\end{lemma}

\begin{proof}
Suppose $r>4e^{h_0}R$.
Since $\E{\pi}{\|X-x_0\|^2}\leq R^2$, Markov gives
$\pi(B(x_0,2R))\geq3/4$, and
$B(x_0,2R)\subseteq B(x_0,r)$.
Writing $Z_V=\int e^{-V}$, on
$B(x_0,r)$ we have
$V\leq\inf V+h_0n$, so
\[
\frac34
\leq\pi(B(x_0,2R))
\leq\frac{\Vol(B(x_0,2R))\,e^{-\inf V}}{Z_V},
\qquad
\pi(B(x_0,r)) \geq\frac{\Vol(B(x_0,r))\,e^{-\inf V-h_0n}}{Z_V}.
\]
Dividing,
$\pi(B(x_0,r))\geq(r/2R)^ne^{-h_0n}\cdot\frac34 >(2e^{h_0})^ne^{-h_0n}\cdot\frac34 =\frac34\cdot2^n>1$,
a contradiction.
\end{proof}

The next lemma concerns the stage count and total cost.

\begin{lemma}
\label{lem:count}
The number of cooling stages satisfies
\[
N=O(n^{1/4}\log(2+R/r)).
\]
With
$p=\min(1/20,\beta/(N+1))$, all
$N+1$ calls of
\algname{AsymmetricSecondMoment} succeed simultaneously with
probability at least $1-\beta$, and the total number of evaluation queries is
\[
 O\left(n^{13/4}\log\left(2+\frac{nR}{r\beta}\right)^{O(1)}\right)
\]
on every execution.
\end{lemma}

\begin{proof}
By the definition of $N$ and $(1-\delta)^j\leq e^{-\delta j}$,
\[
N
\leq\max\left\{0,
\left\lceil\delta^{-1}\log\frac{t_0R^2}{\eta_0}\right\rceil\right\}
=O\bigl(n^{1/4}\log(2+R/r)\bigr),
\]
since $\delta=\gamma n^{-1/4}$ and $\eta_0$ is universal.
By the conditional guarantees
(\cref{lem:induction,def:adaptive-oracle-guarantee}), the failure probability of
the whole run is at most $(N+1)p\leq\beta$.
Each call costs at most $O(n^3\log(n/p)^{O(1)})$ deterministically, and
\[
\log(n/p)
=O\bigl(\log\bigl(n(N+1)/\beta\bigr)\bigr)
=O\Bigl(\log\Bigl(2+\frac{nR}{r\beta}\Bigr)\Bigr),
\]
so the total is
$O(n^{1/4}\log(2+R/r)\cdot n^3\log(2+nR/(r\beta))^{O(1)}) =O(n^{13/4}\log(2+nR/(r\beta))^{O(1)})$.
\end{proof}

The next theorem concerns the grounded-potential rounding.

\begin{theorem}
\label{thm:grounded-rounding}
Given a grounded input $(V,x_0,r,R)$ as in \cref{def:grounded-input} and
$\beta\in(0,1)$, the algorithm
\algname{AsymmetricSecondMoment} (\cref{fig:asymmetric-second-moment})
makes $O(n^{13/4}\log(2+nR/(r\beta))^{O(1)})$ evaluation queries on every
execution, and with probability at least
$1-\beta$ returns a positive semidefinite
$\widehat M$ with
\[
0.99\widehat M
\preceq M_\pi
\preceq1.01\widehat M,
\qquad
\pi\propto e^{-V},
\]
together with an affine map $\mathcal T$ such that
$C_{\mathcal T_\#\pi}\approx_{1.01}I$ and
$\|m_{\mathcal T_\#\pi}\|\leq1/250$.
\end{theorem}

\begin{proof}
On the event that all $N+1$ calls succeed, which has probability at least
$1-\beta$ by \cref{lem:count}, every call returned accurate estimates in the
sense of \cref{lem:fdwr}.
In particular the final call returned $(\widehat m,\widehat C)$ with
$(1-\eps_{\rm cool})C_{F_\#\pi}\preceq\widehat C \preceq(1+\eps_{\rm cool})C_{F_\#\pi}$
and
$\|C_{F_\#\pi}^{-1/2}(\widehat m-m_{F_\#\pi})\|\leq\eps_{\rm cool}$, where
$\eps_{\rm cool}\leq1/500$.
With $F(x)=A(x-b)$ and center
$c=-Ab$, \cref{def:affine-pushforward} gives
$M_\pi=A^{-1}(C_{F_\#\pi}+(m_{F_\#\pi}+Ab)(m_{F_\#\pi}+Ab)^\top )A^{-\top}$,
which is exactly
$A^{-1}M_{F_\#\pi}(c)A^{-\top}$, while the returned matrix is
$\widehat M=A^{-1}\widehat M_{F_\#\pi}(c)A^{-\top}$.
\cref{lem:raw} with $\eps=\zeta=\eps_{\rm cool}\leq1/500$ and Loewner congruence
by
$A^{-1}$ give
$0.99\widehat M\preceq M_\pi\preceq1.01\widehat M$;
$\widehat M$ is positive semidefinite since
$\widehat C\succ0$.
The map $\mathcal T$ pushes
$\pi$ to a law with covariance
$\widehat C^{-1/2}C_{F_\#\pi}\widehat C^{-1/2} \in[(1+\eps_{\rm cool})^{-1},(1-\eps_{\rm cool})^{-1}]I \subseteq[1/1.01,1.01]I$
and mean norm
$\|\widehat C^{-1/2}(m_{F_\#\pi}-\widehat m)\| \leq(1-\eps_{\rm cool})^{-1/2}\,\eps_{\rm cool}\leq1/250$.
The query bound is \cref{lem:count}.
\end{proof}

\begin{proof}[Proof of \cref{thm:general-upper-bound}]
Set $V=I_K$, the indicator potential of \cref{def:oracle-models}: evaluation of
$V$ is exactly membership querying of
$K$, and
$\pi=\pi_K$.
The input is grounded for any $h_0>0$, since
$B(x_0,r)\subseteq K=\{V\leq0\}=\{V\leq\inf V+h_0n\}$ and
$\E{\pi_K}{\|X-x_0\|^2}\leq R^2$ from
$K\subseteq B(x_0,R)$.
\cref{thm:grounded-rounding} then gives the query bound, the success
probability, and
$0.99\widehat M\preceq\E{x\sim K}{xx^\top}=M_K \preceq1.01\widehat M$, as
required.
On the same success event, \cref{thm:grounded-rounding} also returns
$\mathcal T$ with
$C_{\mathcal T_\#\pi_K}\in[1/1.01,1.01]I$ and
$\|m_{\mathcal T_\#\pi_K}\|\leq1/250$; its linear part is invertible because
this covariance is positive definite.
\cref{fact:kls-inner-ball} with $\lambda=1/1.01$ gives
$B(m_{\mathcal T_\#\pi_K},1.01^{-1/2})\subseteq\mathcal T(K)$, and
$1.01^{-1/2}-1/250\geq0.99$, so
$B(0,0.99)\subseteq\mathcal T(K)$.
Finally
\[
\E{x\sim\mathcal T(K)}{\|x\|^2}
=\Tr{C_{\mathcal T_\#\pi_K}}+\|m_{\mathcal T_\#\pi_K}\|^2
\leq1.01n+250^{-2}
\leq1.02n.
\]
\end{proof}

\begin{proof}[Proof of \cref{cor:volume}]
Run \algname{AsymmetricSecondMoment} with $\beta=1/8$.
On its success event, \cref{thm:general-upper-bound} gives an invertible
affine map
$\mathcal T(x)=Ax+a$ with
$B(0,0.99)\subseteq\mathcal T(K)$ and
$\E{x\sim\mathcal T(K)}{\|x\|^2}\leq1.02n$.
The body $K'=0.99^{-1}\mathcal T(K)$ contains the unit ball, satisfies
$\E{x\sim K'}{\|x\|^2}\leq1.02n/0.99^2\leq1.05n$, and
$\Vol(K)=0.99^n\Vol(K')/|\det A|$; a membership query to
$K'$ is one membership query to
$K$.
\cref{fact:cv18-volume} with $R^2=1.05n$ and
$p=1/8$ returns, with probability at least
$7/8$, a number within a factor
$1\pm\eps$ of
$\Vol(K')$ after
$O(n^3\eps^{-2}\log(n/\eps)^{O(1)})$ queries, and the output is
$0.99^n|\det A|^{-1}$ times that number.
A union bound gives success probability at least $3/4$, and the query count
is the sum of the two bounds, with
$\log(2+8nR/r)=O(\log(2+nR/r))$.

If $K$ is symmetric and $x_0=0$, run
\algname{Symmetric-Covariance} with
$\beta=1/8$ instead.
On its success event, $\pi_K$ has mean
$0$ and covariance
$C_K=\E{x\sim K}{xx^\top}$ with
$0.99\Sigma\preceq C_K\preceq1.01\Sigma$, so
$\mathcal T=\Sigma^{-1/2}$ makes the covariance of the uniform law on
$\mathcal T(K)$ lie in
$[0.99,1.01]I$.
\cref{fact:kls-inner-ball} gives $B(0,0.99^{1/2})\subseteq\mathcal T(K)$, and
$\E{x\sim\mathcal T(K)}{\|x\|^2}\leq1.01n$, so the argument above applies
to
$K'=0.99^{-1/2}\mathcal T(K)$, with the rounding cost of
\cref{thm:symmetric-upper-bound}.
\end{proof}

\subsection{Remarks on the Step Size}
\label{sec:step-size-remarks}

The exponent $13/4$ is the cost of one full-density rounding call per cooling
stage, times the number
$O(n^{1/4}\log(2+R/r))$ of stages; the stage count is forced by the step size
$\delta=\Theta(n^{-1/4})$, which in turn is the largest step the moment transfer
of \cref{lem:renyi-trace} tolerates: the one-step R\'enyi cost of
\cref{lem:renyi} is
$O(n\delta^2)$, hence
$O(n^{1/2})$ at this
$\delta$, and the Paouris tail can absorb exactly such
$e^{O(\sqrt n)}$ warmness.
Two conjectural simplifications would remove the $n^{1/4}$ overhead, and both
fail.
First, the coordinate-moment bound of
\cref{conj:gaussian-cooling-coordinate-moment} would make the per-coordinate
derivative of the cooled trace $O(1)$-bounded, permitting constant-size steps;
it is refuted by \cref{thm:coordinate-moment-counterexample}.
Second, even the trace transfer itself is sharp: \cref{thm:sharp-trace-lower}
exhibits restricted Gaussians whose trace after a warming step is
$\Omega(n+n^2\delta^4)$, which stays
$O(n)$ exactly when
$\delta=O(n^{-1/4})$.
In this sense the scalar moment bounds used above cannot support larger steps;
the matrix-valued cooling of \cref{sec:symmetric} is what removes the overhead
for symmetric bodies.
% ===== end 4_asymmetric_kimi.tex =====
% ===== begin 5_symmetric.tex =====

\section{Cubic Algorithm for the Symmetric Case}
\label{sec:symmetric}

We now describe the symmetric cubic algorithm.
The algorithm still uses Gaussian cooling, but it replaces the scalar
temperature $tI$ by a matrix precision
$H$.
This allows the cooling schedule to treat different directions differently:
directions that remain well-conditioned can keep cooling, while directions that
threaten the invariant receive extra regularization.
The rest of this section gives the full proof.

Assume $K=-K$ and
$rB_2^n\subseteq K\subseteq RB_2^n$.
For $H\succeq0$, let
\[
U_H(x)=I_K(x)+x^\top Hx,
\qquad
\mu_H(x)\propto e^{-x^\top Hx}1_K(x),
\qquad
C_H=\E{x\sim\mu_H}{xx^\top}.
\]
We treat the function $U_H$ as the potential, or negative log-density up to an
additive constant, of
$\mu_H$.
It records both the constraint $x\in K$ and the quadratic weight
$x^\top Hx$.
The algorithm changes $H$ to remove this weight gradually.
At every successful iteration $i$, it maintains the potential
$U_{H_i}$ together with a covariance estimate satisfying
\[
W_i\approx_{\kappa_{\rm app}}C_{H_i},
\qquad
\kappa_{\rm app}=1.01.
\]
Initially $H_0=r^{-2}I$.
Under the change of variables $x=rz$, the promised inclusion
$rB_2^n\subseteq K$ becomes
$B_2^n\subseteq r^{-1}K$, and the initial transformed potential is
\[
P_0(z)
=U_{H_0}(rz)
=I_K(rz)+\|z\|_2^2.
\]
If $\|z\|_2\leq1$, then
$rz\in K$, so
\[
P_0(z)
=\|z\|_2^2
\leq1.
\]
Moreover, $P_0\geq0$ and
$P_0(0)=0$.
Therefore
\[
 B_2^n\subseteq\{z:P_0(z)\leq\inf P_0+1\}.
\]
We refer to this ball as the sublevel-set ball for $P_0$.
Since $K=-K$, both
$I_K(rz)$ and
$\|z\|_2^2$ are even functions of
$z$.
Hence $P_0$ is even, the density proportional to
$e^{-P_0}$ is even, and the corresponding random vector
$Z$ has mean zero.
Applying \cref{lem:bl} with $G=I_{r^{-1}K}$ and
$t=1$ gives
\[
\operatorname{Cov}(Z)\preceq\frac12I,
\qquad
\E{P_0}{\|Z\|_2^2}\leq\frac n2.
\]
Thus the initial law has the required $O(n)$ second-moment bound.
The near-cubic covariance estimator in \cref{run01:thm:primitive} can therefore
construct $W_0$.

Given a pair $(H,W)$, the next precision is given by
\[
H^+=(1-\delta)H+ \eta_{\rm ridge}\frac{\Tr{W\,H}}{n}W^{-1},
\qquad
\delta=L^{-2},
\quad
\eta_{\rm ridge}=L^{-3},
\]
where $L=\Theta(1+\log(2+nR/r))$.
The vector $z=W^{-1/2}x$ is the whitened version of
$x$ obtained from the current covariance estimate
$W$.
Since $W\approx_{\kappa_{\rm app}}C_H$, its covariance satisfies
\[
 W^{-1/2}C_HW^{-1/2}\approx_{\kappa_{\rm app}}I,
\]
so all directions have comparable variance in the $z$-coordinates.
In these coordinates, the update becomes
\[
J=(1-\delta)B+ \eta_{\rm ridge}\frac{\Tr{B}}{n}I,
\qquad
B=W^{1/2}HW^{1/2}.
\]
Thus $B$ is the quadratic penalty in these well-conditioned coordinates, and
$J=W^{1/2}H^+W^{1/2}$.
If $\lambda$ is an eigenvalue of
$B$, the update maps it to
\[
 (1-\delta)\lambda+
 \eta_{\rm ridge}\frac{\Tr{B}}{n}.
\]
The additive term therefore has its largest relative effect when $\lambda$ is
small.
Moreover,
\[
\Tr{B}
=\Tr{W\,H}
\approx_{\kappa_{\rm app}}\Tr{C_H\,H}.
\]
The ridge adds total trace $\eta_{\rm ridge}\Tr{B}$, whereas the cooling term
removes
$\delta\Tr{B}$ from
$B$.

The invariant rests on two estimates.
First, even-sector Poincar\'e bounds along the interpolation from $B$ to
$J$ show that the next law has
$O(n)$ second moment in the
$z$-coordinates.
The same fixed-radius ball remains in an $O(n)$-height sublevel set of the next
potential.
These are the two well-roundedness hypotheses of the near-cubic covariance
estimator, so the next covariance estimate costs $\widetilde O(n^3)$ membership
queries.
Second, define
\[
D_{\rm sym}(H)=\log\frac{\Vol(K)}{\int_K e^{-x^\top Hx}\,dx},
\qquad
q_{\rm sym}(H)
=\Tr{C_H\,H}
=\E{X\sim\mu_H}{X^\top HX},
\]
so, for $X$ uniform on
$K$,
\[
 e^{-D_{\rm sym}(H)}=\E{X\sim K}{e^{-X^\top H X}}.
\]
We can understand $D_{\rm sym}(H)$ as measuring how much the quadratic weight
decreases the normalizing integral, while
$q_{\rm sym}(H)$ is the expected quadratic penalty under the current
distribution
$\mu_H$.
In particular, $D_{\rm sym}(H)\geq0$ and
$D_{\rm sym}(0)=0$.
The regularized update gives
\[
D_{\rm sym}(H^+)\leq D_{\rm sym}(H)
-\frac{\delta}{2}q_{\rm sym}(H).
\]
For every matrix reached by the algorithm, a separate comparison gives
\[
q_{\rm sym}(H)
\leq D_{\rm sym}(H)
\leq C_DLq_{\rm sym}(H),
\]
where $C_D$ is a constant.
Consequently,
\[
D_{\rm sym}(H^+)
\leq\left(1-\frac{\delta}{2C_DL}\right)D_{\rm sym}(H)
=\left(1-\Theta(L^{-3})\right)D_{\rm sym}(H).
\]
The initial value satisfies $\log(D_{\rm sym}(H_0)/c_*)=O(L)$ for every fixed
terminal threshold
$c_*>0$.
Thus $O(L^4)$ successful updates make
$D_{\rm sym}(H)$ at most
$c_*$.

At this point the uniform measure on $K$ is well-rounded in the current
coordinates.
The terminal comparison shows that $C_H$ and
$C_K$ are comparable in relative Loewner order, and removing the nonnegative
quadratic term preserves the fixed-radius sublevel-set ball.
One final covariance call on $K$ gives the matrix
$\Sigma$ in \cref{thm:symmetric-upper-bound}.
The algorithm makes $O(L^4)$ calls, each costing
$\widetilde O(n^3)$ membership queries, so its overall query complexity is cubic
up to logarithmic factors.

\subsection{Proof Setup for Even Log-Concave Laws}
\label{sec:symmetric-proof-setup}

We now prove the symmetric upper bound stated in
\cref{thm:symmetric-upper-bound}.
Unlike the preceding scalar schedule, the algorithm cools a matrix-valued
Gaussian precision and adds an adaptive regularizer at each step.
Symmetry keeps every distribution with density proportional to
$e^{-V(x)-x^{\mathsf T}Hx}$ centered.
The remaining subsections give the analytic inputs, invariant, progress
argument, and final algorithm.

\subsection{Quadratic Weights and Analytic Inputs}
\label{sec:symmetric-inputs}

The following covariance estimator is used at each cooling step.
Given an even log-concave distribution satisfying the two displayed
well-roundedness conditions, it returns a $1\%$ relative covariance estimate
using
$O(n^3\log(n/\xi)^{O(1)})$ evaluations.
Its success guarantee remains valid after conditioning on the preceding calls,
and its query count has a deterministic upper bound.
We apply the estimator to distributions with density proportional to
$e^{-V(x)-x^{\mathsf T}Hx}$ for
$H\succeq0$.
The next theorem concerns the near-cubic covariance estimation for even
distributions.

\begin{theorem}
\label{run01:thm:primitive}
Fix universal $b_0,c_1,c_2>0$ and
$\kappa_{\rm app}=1.01$.
Let $U:\mathbb R^n\to\mathbb R\cup\{+\infty\}$ be a proper,
lower-semicontinuous, convex, and even potential whose normalized law
$\nu_U$ satisfies
\[
 \E{\nu_U}{\|Z\|^2}\leq c_1n,
 \qquad
 B(0,b_0)\subseteq\{U\leq\inf U+c_2n\}.
\]
Given an evaluation oracle for $U$ and a failure tolerance
$\xi\in(0,1/2]$, there is an algorithm that returns either a failure symbol or a
positive-definite matrix
$\widehat C$, and with probability at least
$1-\xi$ the latter satisfies
\[
 \widehat C\approx_{\kappa_{\rm app}}
 \operatorname{Cov}_{\nu_U}(Z).
\]
It has a deterministic query limit of
\[
 O\!\left(n^3\log(n/\xi)^{O(1)}\right)
\]
evaluations of $U$, and its guarantee holds under adaptive conditioning in the
sense of \cref{def:adaptive-oracle-guarantee}.
\end{theorem}

\begin{proof}
Set
\[
p=\min\{\xi,1/20\},
\qquad
\varepsilon_{\rm prim}=1/500.
\]
\algname{FullDensityRound} (\cref{lem:fdwr}) takes a ball contained in a
height-$O(n)$ sublevel set of the potential together with an
$O(n)$ second-moment bound and returns relative estimates of the mean and
covariance, with a conditional failure guarantee and a deterministic query
limit.
By \cref{lem:fdwr}, \algname{FullDensityRound} implements this interface.
Apply it with center $0$ and
$(\rho,H,B)=(b_0,c_2,c_1)$, accuracy
$\varepsilon_{\rm prim}$, and failure tolerance
$p$.
If the routine returns a failure symbol, pass that symbol to the caller.
Otherwise discard its mean estimate and return its covariance estimate.
Evenness makes the true mean zero.
Since $p\leq\xi$ and
\[
1-\varepsilon_{\rm prim}\geq\kappa_{\rm app}^{-1},
\qquad
1+\varepsilon_{\rm prim}\leq\kappa_{\rm app},
\]
the routine supplies the stated covariance guarantee, conditional failure bound,
and deterministic cap.
Replacing $\xi$ by
$p$ inside the logarithm changes the cap only by a universal factor.
\end{proof}

The next definition concerns the well-rounded even potential.

\begin{definition}
\label{run01:def:local-well-rounded}
Let $U:\mathbb R^n\to\mathbb R\cup\{+\infty\}$ be proper, lower-semicontinuous,
even, and convex, and suppose
$0<\int_{\mathbb R^n}e^{-U(z)}\,dz<\infty$.
For positive constants $b_0,c_1,c_2$, call
$U$
$(b_0,c_1,c_2)$-well-rounded if its normalized law
$\nu_U$ satisfies
\[
 \E{\nu_U}{\|Z\|^2}\leq c_1n
 \quad\text{and}\quad
 B(0,b_0)\subseteq\{z:U(z)\leq\inf U+c_2n\}.
\]
For a base potential $P$ and $J\succeq0$ for which the denominator below
is positive and finite, we use the shorthand
\[
 \nu_{P,J}(dz)=
 \frac{e^{-P(z)-z^{\mathsf T} Jz}\,dz}
 {\int_{\mathbb R^n}e^{-P(y)-y^{\mathsf T} Jy}\,dy}.
\]
Here ``well-rounded'' is a property of the probability distribution defined by a
potential, not a matrix property.
We call $\nu_{P,J}$ well-rounded when the potential
$U(z)=P(z)+z^{\mathsf T}Jz$ satisfies the two preceding conditions: its
normalized distribution has second moment at most
$c_1n$, and
$B(0,b_0)\subseteq\{U\leq\inf U+c_2n\}$.
\end{definition}
\space The next lemma supplies the whole-space Poincar\'e inequality from which
the normalized and singular-precision inequalities, and hence the energy bounds
used in the cooling analysis, will follow.

The next lemma concerns the whole-space strongly convex even-sector inequality.

\begin{lemma}
\label{run01:lem:whole-space-even-sector}
Let $V\in C^\infty(\mathbb R^n)$ be even, suppose
$\nabla^2V(y)\succeq2I$ for every
$y$, and suppose
$Z_V=\int_{\mathbb R^n}e^{-V(y)}dy<\infty$.
For $d\mu_V=Z_V^{-1}e^{-V(y)}dy$ and every
$h\in C_c^\infty(\mathbb R^n)$,
\begin{equation}
 \V{\mu_V}{h}
 \leq\frac12\int_{\mathbb R^n}\lVert\nabla h\rVert_2^2d\mu_V.
 \label{run01:eq:whole-space-general}
\end{equation}
If $h$ is even, the factor
$1/2$ may be replaced by
$1/4$.
\end{lemma}

\begin{proof}
The Brascamp--Lieb inequality, \cref{lem:brascamp_lieb}, and
$\nabla^2V\succeq2I$ give
\[
\V{\mu_V}{h}
\leq \int \left\langle(\nabla^2V)^{-1}\nabla h,\nabla h\right\rangle d\mu_V
\leq\frac12\int\|\nabla h\|_2^2\,d\mu_V.
\]
This proves \cref{run01:eq:whole-space-general}, without using that $h$ is even.
It remains to improve the constant when $h$ is even.

Suppose now that $h$ is even, and let
$\gamma(dy)=\pi^{-n/2}e^{-\|y\|_2^2}\,dy$.
Caffarelli's contraction theorem, \cref{thm:caffarelli-contraction}, gives a
transport map $T$ from
$\gamma$ to
$\mu_V$ that does not increase distances:
\[
T_\#\gamma=\mu_V,
\qquad
\|DT(y)\|_{\mathrm{op}}\leq1
\quad
\text{for almost every $y$}.
\]
Here the hypothesis of Caffarelli's theorem holds because the source potential
$\|y\|_2^2$ has Hessian
$2I$, while
$\nabla^2V\succeq2I$.
Moreover, $T$ is odd.
Indeed, both $\gamma$ and
$\mu_V$ are unchanged by
$y\mapsto-y$.
Writing $T=\nabla\varphi$ as in Caffarelli's theorem, the reflected map
$y\mapsto-T(-y)$ equals
$\nabla(\varphi(-y))$, and
$\varphi(-y)$ is convex.
This reflected map also sends $\gamma$ to
$\mu_V$, so the uniqueness assertion in Caffarelli's theorem gives
$T(-y)=-T(y)$ almost everywhere.

We recall the elementary Gaussian fact that will be used.
Hermite polynomials play the same role for Gaussian measure that Fourier
characters play on the discrete cube: they give an orthogonal decomposition by
degree.
If $g_k$ denotes the part of
$g$ of total Hermite degree
$k$, then, for the normalization of
$\gamma$ above,
\[
\V{\gamma}{g}=\sum_{k\geq1}\|g_k\|_{L_2(\gamma)}^2,
\qquad
\int\|\nabla g\|_2^2\,d\gamma =\sum_{k\geq1}2k\|g_k\|_{L_2(\gamma)}^2.
\]
Under $y\mapsto-y$, the
degree-$k$ part changes by the sign
$(-1)^k$.
Thus, if $g$ is even, all of its odd-degree parts vanish.
Every nonconstant term in the sums above then has $k\geq2$, and therefore
\begin{equation}
 \V{\gamma}{g}
 \leq\frac14\int\|\nabla g\|_2^2\,d\gamma.
 \label{run01:eq:gaussian-even-sector}
\end{equation}

Apply \cref{run01:eq:gaussian-even-sector} to $g=h\circ T$.
Because $h$ is even and
$T$ is odd,
$g$ is even.
The chain rule and the contraction bound give
\[
 \|\nabla(h\circ T)(y)\|_2
 \leq\|\nabla h(T(y))\|_2
 \quad\text{for almost every $y$}.
\]
Finally, $T_\#\gamma=\mu_V$, so
\[
\V{\mu_V}{h}
=\V{\gamma}{h\circ T}
\leq\frac14\int\|\nabla(h\circ T)\|_2^2\,d\gamma
\leq\frac14\int\|\nabla h\|_2^2\,d\mu_V.
\]
\end{proof}

The next lemma concerns the normalized even-sector Poincar\'e inequality.

\begin{lemma}
\label{run01:lem:normalized-even-sector}
Let $P:\mathbb R^n\to\mathbb R\cup\{+\infty\}$ be proper, lower-semicontinuous,
even, and convex, and suppose
\[
0<Z_P:=\int_{\mathbb R^n}e^{-P(z)-\lVert z\rVert_2^2}\,dz<\infty.
\]
Define $d\gamma_P(y)=Z_P^{-1}e^{-P(y)-\lVert y\rVert_2^2}\,dy$.
For every locally Lipschitz $h$,
\begin{equation}
 \V{\gamma_P}{h}
 \leq\frac12\E{\gamma_P}{\lVert\nabla h\rVert_2^2}.
 \label{run01:eq:normalized-general}
\end{equation}
If $h$ is even almost everywhere, then the factor
$1/2$ in \cref{run01:eq:normalized-general} may be replaced by
$1/4$:
\begin{equation}
 \V{\gamma_P}{h}
 \leq\frac14\E{\gamma_P}{\lVert\nabla h\rVert_2^2}.
 \label{run01:eq:normalized-even}
\end{equation}
\end{lemma}

The next lemma concerns the Poincar\'e inequality for possibly singular
precision.

\begin{lemma}
\label{run01:lem:singular-poincare}
Let $P:\mathbb R^n\to\mathbb R\cup\{+\infty\}$ be proper, lower-semicontinuous,
even, and convex, let
$H\succeq0$, and let
\[
0
<\int_{\mathbb R^n}e^{-P(x)-x^{\mathsf T} Hx}\,dx
<\infty.
\]
Let
\[
 d\nu(x)=
 \frac{e^{-P(x)-x^{\mathsf T} Hx}\,dx}
 {\int_{\mathbb R^n}e^{-P(z)-z^{\mathsf T} Hz}\,dz}.
\]
If $f$ is locally Lipschitz and
$\nabla f(x)\in\operatorname{im}(H)$ almost everywhere, then
\begin{equation}
 \V{\nu}{f}
 \leq\frac12\E{\nu}
 {\nabla f(X)^{\mathsf T} H^\dagger\nabla f(X)}.
 \label{run01:eq:singular-general}
\end{equation}
If $f$ is even, the factor
$1/2$ may be replaced by
$1/4$.
\end{lemma}

\begin{proof}[Proof of \cref{run01:lem:normalized-even-sector}]
After adding a constant, assume $P(0)=\inf P=0$.
Approximate $P$ pointwise by finite, smooth, even, convex functions
$P_j\geq0$; thus
\[
 P_j(x)\longrightarrow P(x)
 \qquad\text{for every }x\in\mathbb R^n.
\]

Set $Z_j=\int e^{-P_j(y)-\|y\|^2}dy$ and let
$\gamma_j$ be the corresponding normalized measure.
Since
\[
0
\leq e^{-P_j(y)-\|y\|^2}
\leq e^{-\|y\|^2},
\]
dominated convergence gives $Z_j\to Z$ and, for every
$h\in C_c^\infty$,
\[
 \int h^k\,d\gamma_j\to\int h^k\,d\gamma_P\quad(k=1,2),
 \qquad
 \int\|\nabla h\|^2d\gamma_j
 \to\int\|\nabla h\|^2d\gamma_P.
\]
Apply \cref{run01:lem:whole-space-even-sector} to $V_j=P_j+\|\cdot\|^2$.
Passing to the limit proves both constants for $C_c^\infty$ functions.
Now let $h$ be compactly supported and Lipschitz.
Choose a smooth, even, nonnegative function $\rho$ with compact support and
$\int\rho=1$, define
\[
 \rho_\varepsilon(x)=\varepsilon^{-n}\rho(x/\varepsilon),
\]
and let
\[
 h_\varepsilon=h*\rho_\varepsilon.
\]
Then $h_\varepsilon\in C_c^\infty$, and as
$\varepsilon\downarrow0$,
\[
 h_\varepsilon\to h,
 \qquad
 \nabla h_\varepsilon\to\nabla h
 \quad\text{almost everywhere}.
\]
For $\varepsilon\leq1$, these functions are supported in a fixed compact set,
and
$\|\nabla h_\varepsilon\|$ is bounded by the Lipschitz constant of
$h$.
Dominated convergence therefore gives convergence of the variances and of the
gradient integrals.
If $h$ is even, then
$h_\varepsilon$ is even because the kernel
$\rho_\varepsilon$ is even.
Applying the smooth inequality to $h_\varepsilon$ and sending
$\varepsilon\downarrow0$ proves the claim for compactly supported Lipschitz
functions.

For a locally Lipschitz $h$, let
$T_M(t)=\max\{-M,\min\{t,M\}\}$, and choose a radial
$\chi_R\in C_c^\infty$ such that
$\chi_R=1$ on
$B_R$,
$\chi_R=0$ outside
$B_{2R}$, and
$\|\nabla\chi_R\|\leq C/R$.
Define
\[
h_{R,M}=\chi_R T_M(h),
\qquad
\nabla h_{R,M} =\chi_R\mathbf 1_{\{|h|<M\}}\nabla h+T_M(h)\nabla\chi_R
\quad
\text{a.e.}
\]
The radial choice makes $h_{R,M}$ even whenever
$h$ is even.
Moreover,
\[
 \|\nabla h_{R,M}\|
 \leq \mathbf 1_{\{|h|<M\}}\|\nabla h\|
 +\frac{CM}{R}\mathbf 1_{B_{2R}\setminus B_R}.
\]
Suppose first that
\[
 \E{\gamma_P}{\|\nabla h\|^2}<\infty.
\]
For fixed $M$, apply the compactly supported inequality to
$h_{R,M}$ and let
$R\to\infty$.
The preceding gradient bound and dominated convergence give
\[
 \V{\gamma_P}{T_M(h)}
 \leq c\,\E{\gamma_P}
 {\mathbf 1_{\{|h|<M\}}\|\nabla h\|^2},
\]
where $c=1/2$ in general and
$c=1/4$ for even
$h$.
If $X,X'$ are independent with law
$\gamma_P$, then Fatou's lemma gives
\[
\V{\gamma_P}{h}
=\frac12\E{X,X'\sim\gamma_P\otimes\gamma_P}{(h(X)-h(X'))^2}
\leq\liminf_{M\to\infty} \V{\gamma_P}{T_M(h)}.
\]
Monotone convergence on the right-hand side proves the claim.
If the gradient integral is infinite, the claimed inequality is immediate.
\end{proof}

\begin{proof}[Proof of \cref{run01:lem:singular-poincare}]
First suppose $H\succ0$.
Set
\[
y=H^{1/2}x,
\qquad
g(y)=f(H^{-1/2}y),
\qquad
\widetilde P(y)=P(H^{-1/2}y).
\]
The law of $Y=H^{1/2}X$ has density proportional to
$e^{-\widetilde P(y)-\|y\|^2}$, so \cref{run01:lem:normalized-even-sector}
applies.
Moreover,
\[
\nabla_y g(y)=H^{-1/2}\nabla_x f(x),
\qquad
\|\nabla_y g(y)\|^2 =\nabla_x f(x)^{\mathsf T}H^{-1}\nabla_x f(x).
\]
Therefore, with $c=1/2$ in general and
$c=1/4$ when
$f$ is even,
\[
\V{\nu}{f}
=\V{\gamma_{\widetilde P}}{g(Y)}
\leq c\,\E{\gamma_{\widetilde P}}{\|\nabla g(Y)\|^2}
=c\,\E{\nu} {\nabla f(X)^{\mathsf T}H^{-1}\nabla f(X)}.
\]
Since $H\succ0$ implies
$H^\dagger=H^{-1}$, this proves both constants for positive-definite
$H$.

Now suppose $H$ is singular and let
$P_0$ be the orthogonal projector onto
$\ker H$.
Define $H_\varepsilon=H+\varepsilon P_0$.
Then $H_\varepsilon\succ0$ and
\[
 H_\varepsilon^{-1}=H^\dagger+\varepsilon^{-1}P_0.
\]
The densities with precision $H_\varepsilon$ converge to the density of
$\nu$ in total variation, because
\[
0
\leq e^{-P(x)-x^{\mathsf T}H_\varepsilon x}
\leq e^{-P(x)-x^{\mathsf T}Hx}
\]
and the right-hand side is integrable.
Since $\nabla f\in\operatorname{im}(H)=\ker(P_0)$ almost everywhere,
\[
 \nabla f^{\mathsf T}H_\varepsilon^{-1}\nabla f
 =\nabla f^{\mathsf T}H^\dagger\nabla f.
\]
For bounded, compactly supported Lipschitz $f$, the positive-definite result and
total-variation convergence give
\[
 \V{\nu}{f}
 \leq c\,\E{\nu}
 {\nabla f^{\mathsf T}H^\dagger\nabla f}.
\]
For the stated locally Lipschitz class, write
$\mathbb R^n=\operatorname{im}(H)\oplus\ker(H)$, and let
$P_H$ be the orthogonal projector onto
$\operatorname{im}(H)$.
The gradient condition implies that $f$ is constant along
$\ker(H)$-directions.
With the truncation $T_M$ and radial cutoff
$\chi_R$ used above, define
\[
 f_{R,M}(x)=\chi_R(P_Hx)T_M(f(x)).
\]
Then $f_{R,M}$ and its gradient are bounded,
$\nabla f_{R,M}\in\operatorname{im}(H)$ almost everywhere, and the
positive-definite inequality passes to the limit
$\varepsilon\downarrow0$ by total-variation convergence.
Letting first $R\to\infty$ and then
$M\to\infty$, exactly as in the preceding proof, gives the stated inequality.
If $H=0$, the gradient condition makes
$f$ constant, so the conclusion is immediate.
\end{proof}
\space The next lemma concerns the energy domination and cooling monotonicity.

\begin{lemma}
\label{run01:lem:energy-conclusions}
Let $V:\mathbb R^n\to\mathbb R\cup\{+\infty\}$ be proper, lower-semicontinuous,
even, and convex, suppose
$0<\int_{\mathbb R^n}e^{-V(x)}\,dx<\infty$, let
$B,D\succeq0$, and let
$a>0$.
Define
\[
 A_{aB+D}=
 \frac{\int_{\mathbb R^n}xx^{\mathsf T}
 e^{-V(x)-x^{\mathsf T}(aB+D)x}\,dx}
 {\int_{\mathbb R^n}e^{-V(x)-x^{\mathsf T}(aB+D)x}\,dx}.
\]
Let $P_B$ be the orthogonal projector onto
$\operatorname{im}(B)$.
Then
\begin{equation}
aB^{1/2}A_{aB+D}B^{1/2}\preceq\frac12P_B,
\qquad
a\operatorname{Tr}(BA_{aB+D}) \leq\frac12\operatorname{rank}(B),
\label{run01:eq:energy-domination}
\end{equation}
and the function $a\mapsto a\operatorname{Tr}(BA_{aB+D})$ is nondecreasing on
$(0,\infty)$.
\end{lemma}

\begin{proof}
The distribution in the lemma is even, so its covariance is $A_{aB+D}$.
Set $H=aB+D$.
For $v\in\mathbb R^n$, let
$u=B^{1/2}v$.
The linear function $x\mapsto u^{\mathsf T}x$ is generally not even, so we apply
the ordinary factor
$1/2$ in \cref{run01:lem:singular-poincare}; also
$u\in\operatorname{im}(B)\subseteq\operatorname{im}(H)$.
This inclusion verifies the gradient condition in that lemma.
We obtain
\[
v^{\mathsf T}B^{1/2}A_{aB+D}B^{1/2}v
\leq \frac12v^{\mathsf T}B^{1/2}H^\dagger B^{1/2}v
\leq \frac1{2a}v^{\mathsf T}P_Bv.
\]
Indeed, $(H^\dagger)^{1/2}B(H^\dagger)^{1/2}\preceq a^{-1}I$; hence
$B^{1/2}H^\dagger B^{1/2}$ has operator norm at most
$a^{-1}$ and is supported on
$\operatorname{im}(B)$, which proves the last inequality.
Thus the first matrix inequality in \cref{run01:eq:energy-domination} follows,
and taking traces gives the second.
If $e(a)=\operatorname{Tr}(BA_{aB+D})$, differentiation gives
\[
 \frac{d}{da}\bigl(ae(a)\bigr)
 =e(a)-a\V{\nu_a}{X^{\mathsf T} BX}.
\]
The function $q(x)=x^{\mathsf T}Bx$ is even and
$\nabla q(x)=2Bx\in\operatorname{im}(H)$.
The even factor $1/4$ in \cref{run01:lem:singular-poincare} therefore gives
\[
\V{\nu_a}{q(X)}
\leq \frac14\E{\nu_a}{4X^{\mathsf T}BH^\dagger BX}
=\E{\nu_a}{X^{\mathsf T}BH^\dagger BX}
\leq \frac1a\E{\nu_a}{X^{\mathsf T}BX}
=\frac{e(a)}a.
\]
Substituting this bound into the derivative identity gives
\[
\frac{d}{da}(ae(a))
\geq e(a)-a\frac{e(a)}a
=0.
\]
It remains to justify differentiation when $V$ is nonsmooth or may take the
value
$+\infty$.
Let $q(x)=x^{\mathsf T}Bx$, let
$\nu_a$ denote the normalized distribution with density proportional to
$e^{-V(x)-x^{\mathsf T}Dx-aq(x)}$, and define
\[
Z(a)=\int e^{-V(x)-x^{\mathsf T}Dx-aq(x)}\,dx,
\qquad
N(a)=\int q(x)e^{-V(x)-x^{\mathsf T}Dx-aq(x)}\,dx.
\]
The integrable log-concave density proportional to $e^{-V(x)-x^{\mathsf T}Dx}$
has finite polynomial moments.
Hence $q$ and
$q^2$ are integrable against this density, and differentiation under the
integral gives
\[
Z'(a)=-N(a),
\qquad
N'(a)=-\int q(x)^2e^{-V(x)-x^{\mathsf T}Dx-aq(x)}\,dx.
\]
Since $e(a)=N(a)/Z(a)$, the quotient rule yields
\[
 e'(a)=-\V{\nu_a}{q(X)}.
\]
The Poincar\'e inequality used above already allows proper lower-semicontinuous,
extended-valued $V$, so the same calculation proves the result without a
smoothness assumption.
\end{proof}
\space The next lemma concerns the even log-concave marginals.

\begin{lemma}
\label{run01:lem:marginal-peak}
Let $f:\mathbb R^n\to[0,\infty)$ be integrable, even, and log-concave in the
extended-valued sense.
For a unit vector $u$, define
\[
 g(t)=\int_{u^\perp} f(tu+y)\,dy,
\]
where in dimension $n=1$ the integral over the zero-dimensional space is
evaluation at its unique point.
Then $g$ is even and log-concave.
In particular, $g(t)\leq g(0)$ for every
$t\in\mathbb R$.
\end{lemma}

\begin{proof}
Evenness follows from the change of variables $y\mapsto-y$ in
$u^\perp$, and log-concavity follows from \cref{lem:prekopa_leindler}.
An even log-concave function on $\mathbb R$ attains its maximum at zero.
\end{proof}

The next lemma concerns the grounded central-section comparison.

\begin{lemma}
\label{run01:lem:central-section}
Let $V$ be proper, lower-semicontinuous, even, and convex, and suppose
\[
0
<\int_{\mathbb R^n}e^{-V(x)}\,dx
<\infty,
\qquad
B(0,r)\subseteq\{V\leq V(0)+h_0n\}.
\]
Let $H\succeq0$, let
$m=\|H\|_{\mathrm{op}}$, fix a unit vector
$u$, and define
\[
 g(t)=\int_{u^\perp}
 e^{-V(tu+y)-(tu+y)^{\mathsf T} H(tu+y)}\,dy.
\]
For every $|t|<r$,
\begin{equation}
 g(t)\geq
 \left(1-\frac{|t|}{r}\right)^{n-1}
 \exp\!\left(-h_0n\frac{|t|}{r}-mt^2\right)g(0).
 \label{run01:eq:section-comparison}
\end{equation}
The factor with exponent zero is interpreted as $1$ when
$n=1$.
\end{lemma}

\begin{proof}
Evenness and convexity imply $V(0)=\inf V$.
Write $\alpha=|t|/r$ and
$\sigma=\operatorname{sgn}(t)$, with an arbitrary choice of
$\sigma$ when
$t=0$.
For $y\in u^\perp$, set
$x_\pm=\pm tu+(1-\alpha)y$.
Convexity, the points $\pm\sigma ru$ in the promised sublevel-set ball, and
$V(0)\leq V(y)$ give
\[
V(\pm\sigma ru)
\leq V(0)+h_0n
\leq V(y)+h_0n.
\]
Since
\[
x_+=(1-\alpha)y+\alpha(\sigma ru),
\qquad
x_-=(1-\alpha)y+\alpha(-\sigma ru),
\]
convexity yields
\[
\frac{V(x_+)+V(x_-)}2
\leq (1-\alpha)V(y) +\frac\alpha2\bigl(V(\sigma ru)+V(-\sigma ru)\bigr)
\leq V(y)+\alpha h_0n.
\]
For the quadratic terms,
\[
 \begin{aligned}
 x_+^{\mathsf T}Hx_+
 &=(1-\alpha)^2y^{\mathsf T}Hy
   +2t(1-\alpha)u^{\mathsf T}Hy+t^2u^{\mathsf T}Hu,\\
 x_-^{\mathsf T}Hx_-
 &=(1-\alpha)^2y^{\mathsf T}Hy
   -2t(1-\alpha)u^{\mathsf T}Hy+t^2u^{\mathsf T}Hu.
 \end{aligned}
\]
The cross terms cancel upon averaging.
Since $(1-\alpha)^2\leq1$ and
$u^{\mathsf T}Hu\leq\|H\|_{\rm op}=m$,
\[
\frac{x_+^{\mathsf T} Hx_++x_-^{\mathsf T} Hx_-}{2}
=(1-\alpha)^2y^{\mathsf T} Hy+t^2u^{\mathsf T} Hu
\leq y^{\mathsf T} Hy+mt^2.
\]
Define
\[
 F_\pm(y)=
 \exp\bigl(-V(x_\pm)-x_\pm^{\mathsf T}Hx_\pm\bigr).
\]
Changing transverse variables by $z=(1-\alpha)y$ gives
\[
 g(\pm t)=(1-\alpha)^{n-1}
 \int_{u^\perp}F_\pm(y)\,dy.
\]
Evenness gives $g(t)=g(-t)$, and the arithmetic--geometric mean inequality
therefore gives
\[
 \begin{aligned}
 g(t)
 &=\frac{g(t)+g(-t)}2\\
 &\geq(1-\alpha)^{n-1}
 \int_{u^\perp}(F_+(y)F_-(y))^{1/2}\,dy\\
 &\geq(1-\alpha)^{n-1}
 e^{-\alpha h_0n-mt^2}
 \int_{u^\perp}e^{-V(y)-y^{\mathsf T} Hy}\,dy.
 \end{aligned}
\]
The last integral is $g(0)$, proving the claim.
\end{proof}

The next lemma concerns the grounded directional second moment.

\begin{lemma}
\label{run01:lem:directional-lower-bound}
Under the assumptions of \cref{run01:lem:central-section}, let $C_H$ be the
second moment of the law proportional to
$e^{-V(x)-x^{\mathsf T} Hx}$.
There is a constant $c_{h_0}>0$, depending only on
$h_0$, such that for every unit vector
$u$,
\begin{equation}
 u^{\mathsf T} C_Hu
 \geq\frac{c_{h_0}}{\|H\|_{\mathrm{op}}+n^2/r^2}.
 \label{run01:eq:directional-bound}
\end{equation}
\end{lemma}

\begin{proof}
The marginal density in direction $u$ is proportional to the function
$g$ in \cref{run01:lem:central-section}.
Its maximum is at zero by \cref{run01:lem:marginal-peak}.
Set
\[
 a=c_{h_0}'\left(\|H\|_{\mathrm{op}}+n^2/r^2\right)^{-1/2},
\]
where $c_{h_0}'>0$ is sufficiently small.
Then $a\leq c_{h_0}'r/n$ and, when
$\|H\|_{\mathrm{op}}>0$,
$a\leq c_{h_0}'\|H\|_{\mathrm{op}}^{-1/2}$.
For $|t|\leq a$, write
$\alpha=|t|/r$.
Choosing $c_{h_0}'\leq1/2$ gives
\[
(1-\alpha)^{n-1}
\geq e^{-2n\alpha}
\geq e^{-2c_{h_0}'},
\qquad
h_0n\alpha\leq h_0c_{h_0}',
\qquad
\|H\|_{\mathrm{op}}t^2\leq(c_{h_0}')^2.
\]
Thus Equation~\eqref{run01:eq:section-comparison} gives $g(t)\geq c_{h_0}''g(0)$
on
$[-a,a]$.
If $p$ is the normalized marginal density, then
\[
1
\geq\int_{-a}^ap(t)\,dt
\geq2ac_{h_0}''p(0),
\qquad
p(0)\leq\frac{1}{2c_{h_0}''a}.
\]

For completeness, any probability density bounded by $M$ satisfies
$\int t^2p(t)\,dt\geq1/(12M^2)$.
Indeed, for $s\geq0$,
\[
\mathbb P(|T|>s)
\geq1-\int_{-s}^sp(t)\,dt
\geq(1-2Ms)_+.
\]
Consequently,
\[
\int t^2p(t)\,dt
=\int_0^\infty2s\,\mathbb P(|T|>s)\,ds
\geq\int_0^{1/(2M)}2s(1-2Ms)\,ds
=\frac{1}{12M^2}.
\]
Apply this with $M=p(0)$ to obtain
\[
u^{\mathsf T} C_Hu
=\int t^2p(t)\,dt
\geq c_{h_0}a^2
\geq\frac{c_{h_0}}{\|H\|_{\mathrm{op}}+n^2/r^2},
\]
as claimed.
\end{proof}

The next lemma concerns the inverse covariance versus precision.

\begin{lemma}
\label{run01:lem:inverse-covariance}
Under the grounded assumptions above, for every $H\succeq0$,
\[
 \|C_H^{-1}\|_{\mathrm{op}}
 \leq C_{h_0}\left(\|H\|_{\mathrm{op}}+\frac{n^2}{r^2}\right),
\]
where $C_{h_0}=c_{h_0}^{-1}$ depends only on
$h_0$.
\end{lemma}

\begin{proof}
\cref{run01:lem:directional-lower-bound} gives
$C_H\succeq c_{h_0}(\|H\|_{\mathrm{op}}+n^2/r^2)^{-1}I$.
Inverting this positive-definite matrix inequality proves the claim.
\end{proof}

The next lemma concerns the covariance-scale sublevel-set ball.

\begin{lemma}
\label{run01:lem:covariance-ground-ball}
Let $U$ be an even convex potential with normalized law of covariance
$C$.
If $W\approx_{\kappa_{\rm app}}C$, then for universal
$b_{\rm g},H_{\rm g}>0$,
\[
 B(0,b_{\rm g})
 \subseteq
 \{z:U(W^{1/2}z)\leq\inf U+H_{\rm g}n\}.
\]
\end{lemma}

\begin{proof}
The transformed law is centered and has covariance between
$\kappa_{\rm app}^{-1}I$ and
$\kappa_{\rm app}I$.
Apply item~4 of \cref{prop:log-concave-toolkit}, the covariance-to-level-set
implication, with a universal rescaling to absorb $\kappa_{\rm app}=1.01$.
\end{proof}
\space
\paragraph{Standing parameters.}
Fix $\kappa_{\rm app}=1.01$.
Choose $C_L$ sufficiently large, depending only on
$h_0$ and the universal constants below, and throughout the symmetric argument
set
\[
L_{\rm sym}=C_L\max\{1,\log(2+nR/r)\},
\qquad
\delta_{\rm sym}=L_{\rm sym}^{-2},
\qquad
\eta_{\rm ridge}=L_{\rm sym}^{-3}.
\]

The next definition concerns the quadratically weighted distributions.

\begin{definition}
\label{run01:def:quadratic-target}
Fix the even potential $V$ from \cref{thm:symmetric-upper-bound}.
For $H\succeq0$, define
\[
Z_V(H)=\int_{\mathbb R^n}e^{-V(x)-x^{\mathsf T} Hx}\,dx,
\qquad
\nu_H(dx)=\frac{e^{-V(x)-x^{\mathsf T} Hx}}{Z_V(H)}\,dx,
\qquad
C_H=\E{\nu_H}{XX^{\mathsf T}}.
\]
Also set
\[
D_{\rm sym}(H)=\log\frac{Z_V(0)}{Z_V(H)},
\qquad
q_{\rm sym}(H)=\operatorname{Tr}(C_H\,H).
\]
Equivalently,
\[
e^{-D_{\rm sym}(H)}
=\frac{Z_V(H)}{Z_V(0)}
=\E{\nu_0}{e^{-X^{\mathsf T}H X}},
\qquad
q_{\rm sym}(H)=\E{\nu_H}{X^{\mathsf T}H X}.
\]
Thus $D_{\rm sym}(H)$ records the reduction in the normalizing integral caused
by the quadratic weight.
We refer to $\nu_H$ as the quadratic tilt of
$\nu_0$ by
$H$.
Every $\nu_H$ is even, so
$C_H$ is both its covariance and its uncentered second moment.
\end{definition}

\paragraph{Standing analytic estimate.}
The singular-precision Poincar\'e inequality in
\cref{run01:lem:singular-poincare} controls the variance of even scalar
functions by their gradient energy in the image of $H$.
Thus, for every $H\succeq0$ and every even locally Lipschitz function
$g$ satisfying
$\nabla g(z)\in\operatorname{im}(H)$ almost everywhere, we use
\begin{equation}
 \V{\nu_H}{g}
 \leq\frac14\E{\nu_H}
 {\nabla g(Z)^{\mathsf T} H^\dagger\nabla g(Z)}.
 \label{run01:eq:standing-poincare}
\end{equation}

The next definition concerns the trace and energy along the interpolation.

\begin{definition}
\label{run01:def:interpolation}
Suppose $B\succeq0$ and
$b=\operatorname{Tr}(B)/n>0$.
Let $\gamma=\eta_{\rm ridge}b$, and for
$s\in[0,1]$ define
\[
a_s=1-\delta_{\rm sym}s,
\qquad
H_s=a_sB+s\gamma I,
\qquad
F(s)=\operatorname{Tr}(C_{H_s}),
\qquad
q(s)=\operatorname{Tr}(BC_{H_s}).
\]
Here the covariance is computed with any fixed even convex base potential in the
current coordinates.
Also set $f(s)=F(s)/n$ and
$u(s)=q(s)/(bn)$.
\end{definition}

The next lemma concerns the coupled differential control.

\begin{lemma}
\label{run01:lem:coupled-differential}
In the setting of \cref{run01:def:interpolation}, suppose $0<\delta_{\rm sym}<1$
and use the even-sector Poincar\'e estimate \eqref{run01:eq:standing-poincare}.
For every $s\in(0,1]$,
\begin{align}
 f'(s)
 &\leq \frac{\delta_{\rm sym}^2}{4a_s\eta_{\rm ridge}}\,u(s),
 \label{run01:eq:f-differential}\\
 u'(s)
 &\leq \frac{\delta_{\rm sym}}{a_s}u(s)
   +\left(\frac{\eta_{\rm ridge}u(s)f(s)}{a_s s}\right)^{1/2}.
 \label{run01:eq:u-differential}
\end{align}
Moreover, these inequalities may be integrated down to $s=0$.
\end{lemma}

\begin{proof}
For brevity write
\[
\delta=\delta_{\rm sym},
\qquad
\eta=\eta_{\rm ridge},
\qquad
C_s=C_{H_s},
\]
and let $G(x)=\|x\|^2$ and
$K(x)=x^{\mathsf T} Bx$.
Differentiating the expectations with respect to $s$ uses
\[
H_s'=-\delta B+\gamma I,
\qquad
\frac{d}{ds}\E{s}{\varphi(X)} =-\operatorname{Cov}_s \bigl(\varphi(X),X^{\mathsf T}H_s'X\bigr).
\]
Taking $\varphi=G$ and
$\varphi=K$ gives
\[
F'(s)=\delta\operatorname{Cov}(G,K)-\gamma\V{s}{G},
\qquad
q'(s)=\delta\V{s}{K}-\gamma\operatorname{Cov}(K,G).
\]
For $s>0$, applying the even-sector variance bound
\eqref{run01:eq:standing-poincare} to
$K$ and
$G$, and using
$H_s\succeq a_sB$ and
$H_s\succeq s\gamma I$, gives
\[
 \begin{aligned}
 \V{s}{K}
 &\leq\frac14\E{s}
 {(2BX)^{\mathsf T}H_s^{-1}(2BX)}\\
 &=\operatorname{Tr}(C_sBH_s^{-1}B)
 \leq \frac{q(s)}{a_s},\\
 \V{s}{G}
 &\leq\frac14\E{s}
 {(2X)^{\mathsf T}H_s^{-1}(2X)}\\
 &=\operatorname{Tr}(C_sH_s^{-1})
 \leq \frac{F(s)}{s\gamma}.
 \end{aligned}
\]
Indeed, restricted to $\operatorname{im}(B)$, the inequality
$H_s\succeq a_sB$ gives
$H_s^{-1}\preceq a_s^{-1}B^\dagger$.
Therefore
\[
BH_s^{-1}B
\preceq a_s^{-1}BB^\dagger B
=a_s^{-1}B,
\]
which is the bound used above.
The scalar Cauchy--Schwarz inequality
\[
 |\operatorname{Cov}(U,V)|
 \leq(\V{s}{U}\V{s}{V})^{1/2},
\]
and the inequality
\[
 \delta xy-\gamma x^2\leq\frac{\delta^2y^2}{4\gamma}
\]
imply
\[
 \begin{aligned}
 F'(s)
 &\leq
 \delta(\V{s}{G}\V{s}{K})^{1/2}
 -\gamma\V{s}{G}\\
 &\leq \frac{\delta^2}{4\gamma}\V{s}{K}
 \leq \frac{\delta^2}{4a_s\gamma}q(s).
 \end{aligned}
\]
Since $\gamma=\eta b$,
$F=nf$, and
$q=bnu$, division by
$n$ proves Equation~\eqref{run01:eq:f-differential}.
Similarly,
\[
 \begin{aligned}
 q'(s)
 &\leq\delta\V{s}{K}
 +\gamma(\V{s}{K}\V{s}{G})^{1/2}\\
 &\leq\frac{\delta}{a_s}q(s)
 +\left(\frac{\gamma q(s)F(s)}{a_s s}\right)^{1/2}.
 \end{aligned}
\]
Dividing by $bn$ proves Equation~\eqref{run01:eq:u-differential}.

Finally, let $P$ denote the fixed base potential in
\cref{run01:def:interpolation}.
Since $H_s\succeq0$,
\[
0
\leq (1+\|x\|^2)e^{-P(x)-x^{\mathsf T}H_sx}
\leq (1+\|x\|^2)e^{-P(x)}.
\]
The function on the right is integrable: after an affine normalization, this
follows from the log-concave tail bound in \cref{prelim:lem:lg_subtails}.
Dominated convergence, applied to the normalizing integral and the second-moment
integrals, shows that $f$ and
$u$ extend continuously to
$s=0$.
In particular, for some $C<\infty$ and all sufficiently small
$s$,
\[
u(s)\leq C,
\qquad
f(s)\leq C,
\qquad
a_s
\geq1-\delta_{\rm sym}
>0.
\]
The right-hand side of Equation~\eqref{run01:eq:f-differential} is then bounded,
and that of Equation~\eqref{run01:eq:u-differential} is $O(1+s^{-1/2})$, which
is integrable at zero.
Integrating over $[\varepsilon,t]$ and letting
$\varepsilon\downarrow0$ proves the final assertion.
\end{proof}

The next lemma concerns the trace bound after one update.

\begin{lemma}
\label{run01:lem:endpoint-trace}
Let $P$ be a fixed proper, lower-semicontinuous, even convex base potential with
$0<\int_{\mathbb R^n}e^{-P(z)}\,dz<\infty$, and let
$C_H$ denote covariance under the law proportional to
$e^{-P-z^{\mathsf T} Hz}$.
Fix $\kappa_{\rm app}=1.01$.
Suppose $B\succeq0$ and
\[
\kappa_{\rm app}^{-1}I
\preceq C_B
\preceq\kappa_{\rm app}I,
\qquad
\|B\|_{\mathrm{op}}\leq\kappa_{\rm app}/2,
\qquad
\operatorname{Tr}(B)\leq\kappa_{\rm app}n/2.
\]
Let $\delta_{\rm sym}=L_{\rm sym}^{-2}$ and
$\eta_{\rm ridge}=L_{\rm sym}^{-3}$, and define
\[
 J_{\rm sym}=(1-\delta_{\rm sym})B+
 \eta_{\rm ridge}\frac{\operatorname{Tr}(B)}nI.
\]
There is a universal $L_{{\rm sym},0}$ such that, for every
$L_{\rm sym}\geq L_{{\rm sym},0}$,
\[
 \operatorname{Tr}(C_{J_{\rm sym}})\leq2\kappa_{\rm app}n.
\]
\end{lemma}

\begin{proof}
Write
\[
\kappa=\kappa_{\rm app},
\qquad
L=L_{\rm sym},
\qquad
\delta=L^{-2},
\qquad
\eta=L^{-3}.
\]
If $\operatorname{Tr}(B)=0$, then
$B=J_{\rm sym}=0$, and the conclusion follows from
$C_B\preceq\kappa I$.
Otherwise set $b=\operatorname{Tr}(B)/n>0$ and use the interpolation from
\cref{run01:def:interpolation}, which parametrizes the path from
$B$ to
$J_{\rm sym}$.
Then $H_0=B$,
$H_1=J_{\rm sym}$, and
\[
f(0)
=\frac{\operatorname{Tr}(C_B)}n
\leq\kappa,
\qquad
u(0)
=\frac{\operatorname{Tr}(BC_B)}{\operatorname{Tr}(B)}
\leq\kappa.
\]
Define $f_*=\sup_{0\leq s\leq1}f(s)$ and
$v(s)=u(s)^{1/2}$.
Since $B\neq0$ and
$C_{H_s}\succ0$,
\[
q(s)
=\operatorname{Tr}(BC_{H_s})
>0,
\qquad
v(s)>0
\quad
(0\leq s\leq1).
\]
The second inequality in \cref{run01:lem:coupled-differential} therefore implies
\[
 v'(s)\leq\frac{\delta}{2a_s}v(s)
 +\frac12\left(\frac{\eta f_*}{a_s s}\right)^{1/2}.
\]
Since $a_s=1-\delta s$,
\[
 \bigl(a_s^{1/2}\,v(s)\bigr)'
 \leq\frac12\left(\frac{\eta f_*}{s}\right)^{1/2}.
\]
Consequently
\[
a_s^{1/2}\,v(s)
\leq v(0)+(\eta f_*s)^{1/2}
\leq\kappa^{1/2}+(\eta f_*)^{1/2}.
\]
For all sufficiently large $L$,
$a_s\geq1/2$, and hence
\[
 u(s)\leq4\kappa+4\eta f_*
 \qquad(0\leq s\leq1).
\]

The first differential inequality, which controls the trace growth along this
path, now gives
\[
f'(s)
\leq\frac{\delta^2}{4a_s\eta}u(s)
\leq\frac{1}{2L}u(s).
\]
Hence, for every $s\in[0,1]$,
\[
 \begin{aligned}
 f(s)
 &\leq f(0)+\frac{1}{2L}\int_0^su(t)\,dt\\
 &\leq\kappa+\frac{2\kappa}{L}+\frac{2}{L^4}f_*.
 \end{aligned}
\]
Taking the supremum over $s$ yields
\[
 f_*\leq\kappa+\frac{2\kappa}{L}+\frac{2}{L^4}f_*.
\]
For every sufficiently large universal $L$, this implies
$f_*\leq2\kappa$.
Since $H_1=J_{\rm sym}$,
$\operatorname{Tr}(C_{J_{\rm sym}})=nf(1)\leq2\kappa_{\rm app}n$.
\end{proof}
\space The next lemma concerns the normalized sublevel set and quadratic
penalty.

\begin{lemma}
\label{run01:lem:normalized-ground}
Let $H\succeq0$, let
$A=C_H$, and let
$W\approx_{\kappa_{\rm app}}A$.
Define
\[
B=W^{1/2}HW^{1/2},
\qquad
P_{H,W}(z)=V(W^{1/2}z)+z^{\mathsf T} Bz.
\]
Then $P_{H,W}$ is a well-rounded even potential with universal constants, and
\[
\|B\|_{\mathrm{op}}\leq\kappa_{\rm app}/2,
\qquad
\operatorname{Tr}(B)\leq\kappa_{\rm app}n/2.
\]
\end{lemma}

\begin{proof}
The covariance of the normalized distribution is $W^{-1/2}C_HW^{-1/2}$, which
lies between
$\kappa_{\rm app}^{-1}I$ and
$\kappa_{\rm app}I$.
Its second moment is therefore at most $\kappa_{\rm app}n$, and
\cref{run01:lem:covariance-ground-ball} supplies a universal ball in an
$O(n)$-height sublevel set of
$P_{H,W}$.
Apply \cref{run01:lem:energy-conclusions} with $B=H$,
$D=0$, and
$a=1$ to obtain
\[
 H^{1/2}C_HH^{1/2}\preceq\frac12P_{\operatorname{im}(H)}.
\]
Since $W\preceq\kappa_{\rm app}C_H$, a congruence transformation by
$H^{1/2}$ gives
\[
H^{1/2}WH^{1/2}
\preceq\kappa_{\rm app}H^{1/2}C_HH^{1/2}
\preceq\frac{\kappa_{\rm app}}2I.
\]
The nonzero eigenvalues of $W^{1/2}HW^{1/2}$ and
$H^{1/2}WH^{1/2}$ agree, so
$\|B\|_{\mathrm{op}}\leq\kappa_{\rm app}/2$.
Taking traces in the same inequality gives
\[
\operatorname{Tr}(B)
=\operatorname{Tr}(W\,H)
\leq\kappa_{\rm app}\operatorname{Tr}(C_H\,H)
\leq\frac{\kappa_{\rm app}n}{2}.
\]
\end{proof}

The next lemma concerns the initial quadratic tilt.

\begin{lemma}
\label{run01:lem:cold-local}
Write
\[
H_0=r^{-2}I,
\qquad
P_0(z)=V(rz)+\|z\|^2.
\]
Then $P_0$ is a well-rounded even potential with constants depending only on
$h_0$.
Consequently, a successful invocation of the covariance estimator in
\cref{run01:thm:primitive} on $P_0$ returns
$\widehat C_0\approx_{\kappa_{\rm app}}\operatorname{Cov}_{P_0}(Z)$, and
\[
W_0
=r^2\widehat C_0
\approx_{\kappa_{\rm app}}C_{H_0}.
\]
\end{lemma}

\begin{proof}
Evenness gives $V(0)=\inf V$.
For $\|z\|\leq1$, the grounded promise gives
\[
P_0(z)
\leq V(0)+h_0n+1
\leq\inf P_0+(h_0+1)n.
\]
Thus the unit ball lies in an $O(n)$-height sublevel set.
Applying the energy-domination estimate \cref{run01:lem:energy-conclusions} with
base potential $z\mapsto V(rz)$,
$B=I$ (the identity matrix),
$D=0$, and
$a=1$ gives
$\E{P_0}{\|Z\|^2}\leq n/2$.
Under $x=rz$, the covariance transforms as
\[
 \operatorname{Cov}_{P_0}(Z)=r^{-2}C_{H_0}.
\]
Thus the covariance estimate returned for $P_0$ satisfies the displayed relation
$W_0=r^2\widehat C_0\approx_{\kappa_{\rm app}}C_{H_0}$.
\end{proof}
\space

\subsection{Regularized Update and Invariant}
\label{sec:symmetric-invariant}

Using the distributions and analytic estimates from \cref{sec:symmetric-inputs},
we define the matrix-regularized update and prove that every accepted state
remains within the well-rounded regime required by the covariance estimator.
The next lemma concerns the adaptive-ridge well-roundedness.

\begin{lemma}
\label{run01:lem:adaptive-well-roundedness}
Let $H\succeq0$, let
$W\approx_{\kappa_{\rm app}}C_H$, and define
\[
\widehat q=\operatorname{Tr}(W\,H),
\qquad
H^+=(1-\delta_{\rm sym})H+ \eta_{\rm ridge}\frac{\widehat q}{n}W^{-1}.
\]
In the coordinates $x=W^{1/2}z$, define
\[
B=W^{1/2}HW^{1/2},
\qquad
J_{\rm sym}=W^{1/2}H^+W^{1/2},
\]
and define the next transformed potential
\[
 P^+(z)=V(W^{1/2}z)+z^{\mathsf T} J_{\rm sym}z.
\]
Then $P^+$ is
$(b_0,c_1,c_2)$-well-rounded for universal positive constants
$b_0,c_1,c_2$.
\end{lemma}

\begin{proof}
The covariance of the current distribution in $z$-coordinates is
\[
C=W^{-1/2}C_HW^{-1/2},
\qquad
\kappa_{\rm app}^{-1}I
\preceq C
\preceq\kappa_{\rm app}I.
\]
By \cref{run01:lem:normalized-ground}, which supplies its sublevel-set ball and
quadratic-penalty bounds, the current transformed potential
\[
 P(z)=V(W^{1/2}z)+z^{\mathsf T} Bz
\]
has a universal ball in an $O(n)$-height sublevel set, while
\[
\|B\|_{\mathrm{op}}\leq\kappa_{\rm app}/2,
\qquad
\operatorname{Tr}(B)\leq\kappa_{\rm app}n/2.
\]
Moreover $\widehat q=\operatorname{Tr}(B)$ and
\[
 J_{\rm sym}=(1-\delta_{\rm sym})B+
 \eta_{\rm ridge}\frac{\operatorname{Tr}(B)}nI.
\]
If $\operatorname{Tr}(B)=0$, then
$B=0$, the next potential equals the current one, and the conclusion already
follows from \cref{run01:lem:normalized-ground}.
Now assume that $\operatorname{Tr}(B)>0$.
The hypotheses of \cref{run01:lem:endpoint-trace}, the one-step trace-growth
bound for the adaptive ridge, hold with base potential $z\mapsto V(W^{1/2}z)$,
so the law of
$P^+$ has second moment at most
$2\kappa_{\rm app}n$.

Let $B(0,b_{\rm g})$ be the ball supplied for
$P$ by \cref{run01:lem:normalized-ground}, so that
\[
 P(z)-P(0)\leq H_{\rm g}n
 \qquad (\|z\|\leq b_{\rm g}).
\]
On this ball,
\[
P^+(z)-P(z)
=-\delta_{\rm sym}z^{\mathsf T} Bz+ \eta_{\rm ridge}\frac{\operatorname{Tr}(B)}n\|z\|^2
\leq\frac{\kappa_{\rm app}\eta_{\rm ridge}b_{\rm g}^2}{2}.
\]
Evenness and convexity give $V(0)=\inf V$, and both quadratic matrices are
positive semidefinite.
Hence $P(0)=\inf P$ and
$P^+(0)=\inf P^+$.
Since $P^+(0)-P(0)=0$, for
$\|z\|\leq b_{\rm g}$,
\[
P^+(z)-\inf P^+
\leq H_{\rm g}n+ \frac{\kappa_{\rm app}\eta_{\rm ridge}b_{\rm g}^2}{2}
\leq H'_{\rm g}n
\]
for a universal $H'_{\rm g}$.
Together with the second-moment bound, this proves both well-roundedness
clauses.
\end{proof}

The pullback of a successful transformed covariance estimate is exact under
congruence.

The next lemma concerns the invariant propagation.

\begin{lemma}
\label{run01:lem:propagation-local}
In the setting of \cref{run01:lem:adaptive-well-roundedness}, let
$C^+=W^{-1/2}C_{H^+}W^{-1/2}$.
If a covariance call returns $\widehat C^+\approx_{\kappa_{\rm app}}C^+$ and
$W^+=W^{1/2}\widehat C^+W^{1/2}$, then
$W^+\approx_{\kappa_{\rm app}}C_{H^+}$.
\end{lemma}

\begin{proof}
The covariance guarantee is
\[
\kappa_{\rm app}^{-1}C^+
\preceq \widehat C^+
\preceq \kappa_{\rm app}C^+.
\]
Multiplying all three matrices on the left and right by $W^{1/2}$ preserves
Loewner order.
Since
\[
W^{1/2}C^+W^{1/2}=C_{H^+},
\qquad
W^{1/2}\widehat C^+W^{1/2}=W^+,
\]
we obtain
\[
\kappa_{\rm app}^{-1}C_{H^+}
\preceq W^+
\preceq \kappa_{\rm app}C_{H^+}.
\]
\end{proof}
\space

\subsection{Contraction of \texorpdfstring{$D_{\rm sym}$}{D sym} and Comparison with the Original Distribution}
\label{sec:symmetric-progress}

We prove one-step progress and reachable-state control, then show that when
$D_{\rm sym}$ is sufficiently small, the original distribution
$\nu_0$ is well-rounded and comparable in relative Loewner order to the final
distribution
$\nu_H$.
These results prove the contraction and show that the final
covariance-estimation call in the next subsection satisfies its hypotheses.
The next lemma concerns the one-step decrease of $D_{\rm sym}$.

\begin{lemma}
\label{run01:lem:deficit-step-local}
Let $H\succeq0$, let
$W\approx_{\kappa_{\rm app}}C_H$, and let
$H^+$ be the update in \cref{run01:lem:adaptive-well-roundedness}.
If $C_L$ is sufficiently large, then
\[
 D_{\rm sym}(H^+)\leq D_{\rm sym}(H)
 -\frac{\delta_{\rm sym}}2q_{\rm sym}(H).
\]
\end{lemma}

\begin{proof}
Set $\Delta=H^+-H$. Since $H,H^+\succeq0$, both normalizing constants
are positive and finite. Jensen's inequality gives
\[
\begin{aligned}
D_{\rm sym}(H^+)-D_{\rm sym}(H)
&=-\log\E{X\sim\nu_H}{e^{-X^{\mathsf T}\Delta X}}\\
&\leq\E{X\sim\nu_H}{X^{\mathsf T}\Delta X}
=\operatorname{Tr}(C_H\Delta).
\end{aligned}
\]
The quadratic expectation is finite because $\nu_H$ has finite second moments.
Substituting the update formula gives
\[
 \begin{aligned}
 D_{\rm sym}(H^+)-D_{\rm sym}(H)
 &\leq\operatorname{Tr}(C_H(H^+-H))\\
 &=-\delta_{\rm sym}q_{\rm sym}(H)
 +\eta_{\rm ridge}\frac{\widehat q}{n}
  \operatorname{Tr}(C_HW^{-1}).
 \end{aligned}
\]
The covariance comparison gives
\[
\widehat q
=\operatorname{Tr}(W\,H)
\leq\kappa_{\rm app}q_{\rm sym}(H),
\qquad
\operatorname{Tr}(C_HW^{-1})\leq\kappa_{\rm app}n.
\]
Substitution gives
\[
 D_{\rm sym}(H^+)-D_{\rm sym}(H)
 \leq-(\delta_{\rm sym}
 -\eta_{\rm ridge}\kappa_{\rm app}^2)q_{\rm sym}(H).
\]
Increasing $C_L$ ensures
$\eta_{\rm ridge}\kappa_{\rm app}^2\leq\delta_{\rm sym}/2$.
\end{proof}

The next lemma combines energy domination, the grounded inverse-covariance
bound, and the global moment promise to control the reachable precision and
compare $D_{\rm sym}(H)$ with
$q_{\rm sym}(H)$.

The next lemma concerns the reachable cap, comparison, and contraction.

\begin{lemma}
\label{run01:lem:global-progress}
Start from $H_0=r^{-2}I$ and from a successful initial estimate
$W_0\approx_{\kappa_{\rm app}}C_{H_0}$ supplied by \cref{run01:lem:cold-local}.
Suppose that along the first $j$ updates every covariance call succeeds, so
$W_i\approx_{\kappa_{\rm app}}C_{H_i}$ for
$0\leq i\leq j$.
Then a constant $C_{\rm prec}$ depending only on
$h_0$ satisfies
\[
 \|H_i\|_{\mathrm{op}}\leq C_{\rm prec}\frac{n^2}{r^2}
 \qquad(0\leq i\leq j).
\]
Every such reachable $H$ satisfies, for a constant
$C_D$ depending only on
$h_0$,
\[
q_{\rm sym}(H)
\leq D_{\rm sym}(H)
\leq C_DL_{\rm sym}q_{\rm sym}(H).
\]
Consequently
\[
 D_{\rm sym}(H_{i+1})\leq
 \left(1-\frac{\delta_{\rm sym}}{2C_DL_{\rm sym}}\right)
 D_{\rm sym}(H_i).
\]
For every sufficiently small universal $c_*>0$, there is a sufficiently large
constant
$C_T$ such that, if
$T=\lceil C_TL_{\rm sym}^4\rceil$ and the first
$T$ update calls succeed, then
\[
 D_{\rm sym}(H_T)\leq c_*.
\]
\end{lemma}

\begin{proof}
Let $C_p=C_{h_0}$ be the constant from the grounded lower covariance estimate in
\cref{run01:lem:inverse-covariance}, and define
$s=n^2/r^2$ and
$u_i=\|H_i\|_{\mathrm{op}}$.
By the energy bound in \cref{run01:lem:energy-conclusions},
$q_{\rm sym}(H_i)\leq n/2$, and therefore
$\widehat q_i\leq\kappa_{\rm app}n/2$.
From $W_i\succeq\kappa_{\rm app}^{-1}C_{H_i}$ and the resulting
inverse-covariance bound in \cref{run01:lem:inverse-covariance},
\[
 \|W_i^{-1}\|_{\mathrm{op}}
 \leq\kappa_{\rm app}C_p(u_i+s).
\]
For $c_0=\kappa_{\rm app}^2C_p/2$, the update satisfies
\[
 \begin{aligned}
 u_{i+1}
 &\leq(1-\delta_{\rm sym})u_i
 +\eta_{\rm ridge}\frac{\widehat q_i}{n}
   \|W_i^{-1}\|_{\rm op}\\
 &\leq(1-\delta_{\rm sym}+c_0\eta_{\rm ridge})u_i
 +c_0\eta_{\rm ridge}s.
 \end{aligned}
\]
Because $\eta_{\rm ridge}/\delta_{\rm sym}=L_{\rm sym}^{-1}$, increasing
$C_L$ makes
$c_0\eta_{\rm ridge}\leq\delta_{\rm sym}/2$.
Since $u_0=r^{-2}\leq s$, we have
\[
 u_{i+1}\leq
 \left(1-\frac{\delta_{\rm sym}}2\right)u_i
 +c_0\eta_{\rm ridge}s.
\]
Let $C_{\rm prec}=\max\{1,2c_0\}$.
If $u_i\leq C_{\rm prec}s$, then
\[
u_{i+1}
\leq C_{\rm prec}s -\left(\frac{C_{\rm prec}\delta_{\rm sym}}2 -c_0\eta_{\rm ridge}\right)s
\leq C_{\rm prec}s,
\]
where the last inequality uses $\eta_{\rm ridge}\leq\delta_{\rm sym}$.
Induction proves $u_i\leq C_{\rm prec}s$.

Fix a reachable $H$ and define
\[
Q_H(x)=x^{\mathsf T}Hx,
\qquad
Z(t)=\int e^{-V(x)-tQ_H(x)}\,dx,
\qquad
N(t)=\int Q_H(x)e^{-V(x)-tQ_H(x)}\,dx.
\]
The moment promise and the fourth-moment bound in
\cref{prop:log-concave-toolkit}(iii) give
\[
\E{\pi}{\|X\|^4}
\leq n\sum_{j=1}^n\E{\pi}{X_j^4}
\leq C_4n\left(\sum_{j=1}^n\E{\pi}{X_j^2}\right)^2
\leq C_4nR^4.
\]
Since
\[
0
\leq Q_H(x)
\leq\|H\|_{\rm op}\|x\|^2,
\]
both $Q_H$ and
$Q_H^2$ are integrable under
$\pi$.
For $t\in[0,1]$,
$e^{-tQ_H}\leq1$, so these same integrable functions dominate the derivatives of
$Z(t)$ and
$N(t)$.
Differentiation under the integral gives
\[
Z'(t)=-N(t),
\qquad
N'(t)=-\int Q_H(x)^2e^{-V(x)-tQ_H(x)}\,dx.
\]
Moreover,
\[
e(t)
:=\operatorname{Tr}(HC_{tH})
=\frac{N(t)}{Z(t)}.
\]
It follows that
\[
D_{\rm sym}(H)=\int_0^1e(t)\,dt,
\qquad
e'(t)
=-\V{\nu_{tH}}{X^{\mathsf T} HX}
\leq0.
\]
It follows that $D_{\rm sym}(H)\geq e(1)=q_{\rm sym}(H)$.
For the reverse inequality set
\[
\tau
=\min\!\left\{1,\frac{r^2}{n^2R^2}\right\}
=\min\{1,(sR^2)^{-1}\}.
\]
The moment promise gives $C_0\preceq R^2I$ and hence
\[
\int_0^\tau e(t)\,dt
\leq\tau e(0)
\leq\tau R^2\operatorname{Tr}(H)
\leq\frac{\operatorname{Tr}(H)}s.
\]
The inverse-covariance bound and the precision cap imply
\[
C_H\succeq\frac{1}{C_p(C_{\rm prec}+1)s}I,
\qquad
q_{\rm sym}(H)\geq \frac{\operatorname{Tr}(H)}{C_p(C_{\rm prec}+1)s}.
\]
Therefore
\[
 \int_0^\tau e(t)\,dt
 \leq C_p(C_{\rm prec}+1)q_{\rm sym}(H).
\]
If $\tau<1$, energy monotonicity in \cref{run01:lem:energy-conclusions} says
that
$t e(t)$ is nondecreasing, so for
$t\in[\tau,1]$,
\[
 e(t)\leq q_{\rm sym}(H)/t.
\]
Consequently
\[
\int_\tau^1e(t)\,dt
\leq q_{\rm sym}(H)\log(1/\tau)
\leq 2L_{\rm sym}q_{\rm sym}(H),
\]
where $\log(1/\tau)=2\log(nR/r)\leq2L_{\rm sym}$ when
$\tau<1$.
Adding the two integrals proves the comparison between $D_{\rm sym}$ and
$q_{\rm sym}$ with
$C_D=C_p(C_{\rm prec}+1)+2$, including
$H=0$.

Combining this comparison with the one-step decrease in
\cref{run01:lem:deficit-step-local} gives the explicit contraction
\[
 D_{\rm sym}(H_{i+1})
 \leq\left(1-\frac{\delta_{\rm sym}}{2C_DL_{\rm sym}}\right)
 D_{\rm sym}(H_i).
\]
Finally, Jensen's inequality and $H_0=r^{-2}I$ give
\[
D_{\rm sym}(H_0)
=-\log\E{\pi}{e^{-\|X\|^2/r^2}}
\leq\frac{\E{\pi}{\|X\|^2}}{r^2}
\leq\frac{R^2}{r^2}.
\]
The logarithm of the ratio between this bound and $c_*$ is
$O(L_{\rm sym})$.
Since $\delta_{\rm sym}=L_{\rm sym}^{-2}$, iterating the contraction gives
\[
 D_{\rm sym}(H_T)
 \leq
 \exp\left(-\frac{T\delta_{\rm sym}}{2C_DL_{\rm sym}}\right)
 D_{\rm sym}(H_0).
\]
Since $T\delta_{\rm sym}/L_{\rm sym}=\Theta(C_TL_{\rm sym})$, choosing
$C_T$ sufficiently large makes the last expression at most
$c_*$.
\end{proof}

\paragraph{Comparison with the original distribution.}
\space The next lemma concerns the comparison with the original distribution.

\begin{lemma}
\label{run01:lem:terminal-exact}
Let $H$ be reachable in the sense of \cref{run01:lem:global-progress}, let
$W\approx_{\kappa_{\rm app}}C_H$, and suppose
\[
 D_{\rm sym}(H)\leq c_*
\]
for a sufficiently small universal $c_*>0$.
Then the exact transformed potential
\[
 P_*(z)=V(W^{1/2}z)
\]
is $(b_0,c_1,c_2)$-well-rounded for universal constants.
Moreover, if $C_0=\E{\pi}{XX^{\mathsf T}}$, then
\[
e^{-D_{\rm sym}(H)}C_H
\preceq C_0
\preceq \frac{e^{-D_{\rm sym}(H)}} {1-(C_4D_{\rm sym}(H))^{1/2}}C_H,
\]
where $C_4$ is the fourth-moment constant from part (iii) of
\cref{prop:log-concave-toolkit}.
\end{lemma}

\begin{proof}
Let $Q(X)=X^{\mathsf T} HX$.
Splitting the original-law second moment according to $1=e^{-Q}+(1-e^{-Q})$
gives
\[
C_0=e^{-D_{\rm sym}(H)}C_H+E_H,
\qquad
\text{where}
\qquad
E_H=\E{\pi}{(1-e^{-Q(X)})XX^{\mathsf T}}.
\]
For every vector $u$, Cauchy--Schwarz and item~3 of
\cref{prop:log-concave-toolkit}, the fourth-moment estimate for affine
functionals, give
\[
 \begin{aligned}
 u^{\mathsf T} E_Hu
 &\leq
 \bigl(\E{\pi}{(1-e^{-Q(X)})^2}\bigr)^{1/2}
 \bigl(\E{\pi}{(u^{\mathsf T} X)^4}\bigr)^{1/2}\\
 &\leq
 C_4^{1/2}(1-e^{-D_{\rm sym}(H)})^{1/2}\;u^{\mathsf T} C_0u\\
 &\leq(C_4D_{\rm sym}(H))^{1/2}\;u^{\mathsf T} C_0u.
 \end{aligned}
\]
Here
\[
0
\leq1-e^{-Q(X)}
\leq1,
\qquad
\E{\pi}{1-e^{-Q(X)}}=1-e^{-D_{\rm sym}(H)},
\]
so $\E{\pi}{(1-e^{-Q(X)})^2}\leq1-e^{-D_{\rm sym}(H)}$.
Thus $0\preceq E_H\preceq(C_4D_{\rm sym}(H))^{1/2}\,C_0$.
Let $\alpha=(C_4D_{\rm sym}(H))^{1/2}$.
Substitution into the exact identity gives
\[
e^{-D_{\rm sym}(H)}C_H\preceq C_0,
\qquad
(1-\alpha)C_0\preceq e^{-D_{\rm sym}(H)}C_H.
\]
Since $\alpha<1$, division by
$1-\alpha$ proves the two stated relative-Loewner bounds.

Choose $c_*$ so that
$(C_4c_*)^{1/2}\leq1/4$.
Since $W\approx_{\kappa_{\rm app}}C_H$, the covariance of the exact law in
$W$-coordinates satisfies
\[
cI
\preceq W^{-1/2}C_0W^{-1/2}
\preceq CI
\]
for universal $c,C>0$.
In particular, its second moment is at most $Cn$.
Item~4 of \cref{prop:log-concave-toolkit}, the covariance-to-level-set
implication, supplies a universal-radius ball in an $O(n)$-height sublevel set
of
$P_*$.
Hence $P_*$ satisfies both clauses of \cref{run01:def:local-well-rounded}.
\end{proof}
\space

\subsection{Algorithm and Complexity}
\label{sec:symmetric-algorithm}

We assemble the regularized update into the full symmetric algorithm and use the
preceding invariant and contraction results for its correctness, failure
accounting, deterministic evaluation-query limit, and well-rounded linear
output.
Indicator potentials recover the membership-oracle convex-body specialization.
The next definition concerns the adaptive state update.

\begin{definition}
\label{run01:def:adaptive-update}
For a positive-definite $W$ and
$H\succeq0$, define
\[
\widehat q_{\rm sym}(H,W)=\operatorname{Tr}(W\,H),
\qquad
\mathcal U(H,W)=(1-\delta_{\rm sym})H+ \eta_{\rm ridge}\frac{\widehat q_{\rm sym}(H,W)}nW^{-1}.
\]
For $H^+=\mathcal U(H,W)$ define
\[
J_{\rm sym}(H^+,W)=W^{1/2}H^+W^{1/2},
\qquad
P(H^+,W)(z)=V(W^{1/2}z)+z^{\mathsf T} J_{\rm sym}(H^+,W)z.
\]
\end{definition}

Fix once and for all a single triple of constants for the estimator in
\cref{run01:thm:primitive}: take its radius constant no larger than the smallest
radius certified by
\cref{run01:lem:cold-local,run01:lem:adaptive-well-roundedness,run01:lem:terminal-exact},
and take its moment and height constants no smaller than the largest of the
corresponding three certificates.
Since $h_0$ is fixed, this common triple is universal.
Every estimator call below uses this same choice.

\begin{figure}[H]
    \centering
\begin{algbox}
\textbf{\uline{\algname{Symmetric-Covariance}}}.
\\ \uline{Input}: An even grounded input $(V,0,r,R)$ in the sense of
\cref{def:grounded-input} (the ground point of an even potential is
$x_0=0$), and a failure probability
$\beta\in(0,1)$.
\\ \uline{Parameters}:
\[
L_{\rm sym}=C_L\max\{1,\log(2+nR/r)\},
\qquad
T=\lceil C_TL_{\rm sym}^4\rceil,
\qquad
\xi=\frac{\beta}{2(T+2)}.
\]
In particular, $\xi\in(0,1/2]$ for every
$n\geq1$.
The update map $\mathcal U$ of \cref{run01:def:adaptive-update} is used with
$\delta_{\rm sym}=L_{\rm sym}^{-2}$ and
$\eta_{\rm ridge}=L_{\rm sym}^{-3}$, and every call to the covariance
estimator of \cref{run01:thm:primitive} runs at its fixed accuracy
$\kappa_{\rm app}=1.01$.
\\ \uline{Output}: A positive-definite matrix $\Sigma$ and a well-rounding
linear map $T_*$, or a failure symbol.
\begin{enumerate}
\item Set $H_0=r^{-2}I$ and define the rescaled initial potential
 \[
  P_0(z)=V(rz)+\|z\|^2.
 \]
Invoke the covariance estimator of \cref{run01:thm:primitive} on $(P_0,\xi)$.
If the call fails, return the failure symbol.
Otherwise, denote its output by $\widehat C_0$.
If a required inverse or positive square root does not exist, return the failure
symbol; otherwise, set $W_0=r^2\widehat C_0$.
\item For $j=0,\ldots,T-1$, form
 \[
H_{j+1}=\mathcal U(H_j,W_j),
\qquad
P_{j+1}=P(H_{j+1},W_j).
 \]
Invoke the covariance estimator of \cref{run01:thm:primitive} on
$(P_{j+1},\xi)$.
If the call fails, return the failure symbol.
Otherwise, denote its output by $\widehat C_{j+1}$.
If a required inverse or positive square root does not exist, return the failure
symbol; otherwise, set
 \[
  W_{j+1}=W_j^{1/2}\widehat C_{j+1}W_j^{1/2}.
 \]
\item Define the final exact transformed potential
 \[
  P_*(z)=V(W_T^{1/2}z).
 \]
Invoke the covariance estimator of \cref{run01:thm:primitive} on $(P_*,\xi)$.
If the call fails or a required inverse or positive square root does not exist,
return the failure symbol.
Otherwise, denote its output by $\widehat C_*$ and return
 \[
\Sigma=W_T^{1/2}\widehat C_*W_T^{1/2},
\qquad
T_*=W_T^{-1/2}.
 \]
\end{enumerate}
\end{algbox}
    \caption{Near-cubic covariance estimation for even potentials.}
    \label{fig:symmetric-covariance}
\end{figure}

For any transformed distribution, evaluating the displayed potential at $z$ uses
one query at the single point
$W^{1/2}z$ (or
$rz$ in the initial call), plus evaluation of a known quadratic.
Thus each estimator evaluation uses exactly one evaluation of $V$.
When $V=I_K$, this is exactly one membership query to
$K$.

\begin{proof}[Proof of \cref{thm:symmetric-upper-bound}]
The algorithm is \algname{Symmetric-Covariance}
(\cref{fig:symmetric-covariance}).
By \cref{run01:lem:cold-local}, the initial quadratic tilt has the fixed-radius
sublevel-set ball and bounded second moment required by the covariance
estimator, so $P_0$ is well-rounded.
On success of the initial call, $W_0\approx_{\kappa_{\rm app}}C_{H_0}$.
Condition on the event that all calls before update $j$ succeeded.
Repeated use of \cref{run01:lem:propagation-local}, which transports a
successful covariance estimate back by congruence, gives
$W_j\approx_{\kappa_{\rm app}}C_{H_j}$, and
\cref{run01:lem:adaptive-well-roundedness}, which controls the next trace and
sublevel-set ball, verifies both clauses of \cref{run01:def:local-well-rounded}
for the explicit potential
$P_{j+1}$.
The covariance estimator is therefore applicable with failure tolerance $\xi$,
and on success the same congruence argument establishes the next invariant.

Conditioned on success through the first $T$ update calls,
\cref{run01:lem:global-progress}, whose contraction estimate gives
$D_{\rm sym}(H_T)\leq c_*$, shows that the terminal threshold is reached.
Thus \cref{run01:lem:terminal-exact}, the comparison between the moments under
$\nu_{H_T}$ and
$\nu_0$, verifies both well-roundedness clauses for the original transformed
potential
$P_*$.
A successful final call gives
\[
 \widehat C_*\approx_{\kappa_{\rm app}}
 W_T^{-1/2}C_0W_T^{-1/2},
\]
and congruence by $W_T^{1/2}$ yields
$\Sigma\approx_{\kappa_{\rm app}}C_0$.
Since $\kappa_{\rm app}^{-1}>0.99$, this is the required matrix inequality.
The law $\pi$ is full dimensional because its density has positive integral with
respect to
$n$-dimensional Lebesgue measure, so
$C_0$ and hence
$\Sigma$ are positive definite.
The well-rounded output clauses are exactly those proved for $P_*$ in
\cref{run01:lem:terminal-exact}.

The covariance-estimator guarantee applies conditionally on every preceding
successful history.
There are $T+2$ calls, so a union bound gives
\[
\mathbb P(\text{some call fails})
\leq(T+2)\xi
=\frac\beta2
\leq\beta.
\]
This argument does not condition the cost bound on success.
Each call is stopped at the estimator's deterministic query limit.
At $H=0$, the grounded lower covariance estimate in
\cref{run01:lem:inverse-covariance} gives
\[
 C_0\succeq c_{h_0}\frac{r^2}{n^2}I.
\]
Consequently the moment promise implies
\[
R^2
\geq\operatorname{Tr}(C_0)
\geq c_{h_0}\frac{r^2}{n},
\qquad
\log n=O_{h_0}(L_{\rm sym}).
\]
Since $T+2=O(L_{\rm sym}^4)$, it follows that
\[
 \log(n/\xi)
 =O_{h_0}\!\left(\log(2+nR/(r\beta))\right).
\]
Multiplying the $T+2$ calls by the estimator's per-call
$n^3\log(n/\xi)^{O(1)}$ query limit gives the deterministic total query limit
\[
 O\!\left(n^3\log(2+nR/(r\beta))^{O(1)}\right)
\]
evaluations.
The one-query transformed-oracle implementation completes the resource
accounting and its membership-oracle specialization.
\end{proof}
\space
% ===== end 5_symmetric.tex =====

% ===== bibliography (bibtex alpha output of ref.bib, inlined) =====
\newcommand{\etalchar}[1]{$^{#1}$}
\begin{thebibliography}{CDWY18}

\bibitem[BL76]{BL76:journal}
Herm~Jan Brascamp and Elliott~H. Lieb.
\newblock {On Extensions of the Brunn-Minkowski and Pr{\'e}kopa-Leindler
  Theorems, Including Inequalities for Log-Concave Functions, and with an
  Application to the Diffusion Equation}.
\newblock {\em Journal of Functional Analysis}, 22(4):366--389, 1976.

\bibitem[Caf00]{C00:journal}
Luis~A. Caffarelli.
\newblock Monotonicity properties of optimal transportation and the fkg and
  related inequalities.
\newblock {\em Communications in Mathematical Physics}, 214(3):547--563, 2000.

\bibitem[CDWY17]{CDWY17}
Yuansi Chen, Raaz Dwivedi, Martin~J. Wainwright, and Bin Yu.
\newblock Vaidya walk: A sampling algorithm based on the volumetric barrier.
\newblock In {\em 2017 55th Annual Allerton Conference on Communication,
  Control, and Computing}, pages 1220--1227, 2017.

\bibitem[CDWY18]{CDWY18:journal}
Yuansi Chen, Raaz Dwivedi, Martin~J. Wainwright, and Bin Yu.
\newblock Fast mcmc sampling algorithms on polytopes.
\newblock {\em Journal of Machine Learning Research}, 19(55):1--86, 2018.

\bibitem[Che21]{C21:journal}
Yuansi Chen.
\newblock An almost constant lower bound of the isoperimetric coefficient in
  the {KLS} conjecture.
\newblock {\em Geometric and Functional Analysis}, 31:34--61, 2021.

\bibitem[CV18]{CV18:journal}
Ben Cousins and Santosh Vempala.
\newblock Gaussian cooling and $o^*(n^3)$ algorithms for volume and gaussian
  volume.
\newblock {\em SIAM Journal on Computing}, 47(3):1237--1273, 2018.
\newblock Available at:~\url{https://arxiv.org/abs/1409.6011}.

\bibitem[DFK91]{DFK91:journal}
Martin~E. Dyer, Alan~M. Frieze, and Ravi Kannan.
\newblock A random polynomial time algorithm for approximating the volume of
  convex bodies.
\newblock {\em Journal of the ACM}, 38(1):1--17, 1991.

\bibitem[EL14]{EL14:book}
Ronen Eldan and Joseph Lehec.
\newblock Bounding the norm of a log-concave vector via thin-shell estimates.
\newblock In Bo'az Klartag and Emanuel Milman, editors, {\em Geometric Aspects
  of Functional Analysis}, volume 2116 of {\em Lecture Notes in Mathematics},
  pages 107--122. Springer, Cham, 2014.

\bibitem[Eld13]{E13:journal}
Ronen Eldan.
\newblock Thin shell implies spectral gap up to polylog via a stochastic
  localization scheme.
\newblock {\em Geometric and Functional Analysis}, 23(2):532--569, 2013.

\bibitem[Fle10]{F10:journal}
B.~Fleury.
\newblock Concentration in a thin {E}uclidean shell for log-concave measures.
\newblock {\em Journal of Functional Analysis}, 259(4):832--841, 2010.

\bibitem[Fra97]{F97:journal}
Matthieu Fradelizi.
\newblock Sections of convex bodies through their centroid.
\newblock {\em Archiv der Mathematik}, 69:515--522, 1997.

\bibitem[GC11]{GC11:journal}
Mark Girolami and Ben Calderhead.
\newblock Riemann manifold langevin and hamiltonian monte carlo methods.
\newblock {\em Journal of the Royal Statistical Society: Series B (Statistical
  Methodology)}, 73(2):123--214, 2011.

\bibitem[GKM{\etalchar{+}}24]{GKMSZ24}
Yuzhou Gu, Nikki~Lijing Kuang, Yi-An Ma, Zhao Song, and Lichen Zhang.
\newblock Log-concave sampling from a convex body with a barrier: a robust and
  unified dikin walk.
\newblock In {\em Advances in Neural Information Processing Systems},
  volume~37, 2024.

\bibitem[GKV24]{GKV24}
Khashayar Gatmiry, Jonathan Kelner, and Santosh~S. Vempala.
\newblock Sampling polytopes with riemannian hmc: Faster mixing via the lewis
  weights barrier.
\newblock In Shipra Agrawal and Aaron Roth, editors, {\em Proceedings of Thirty
  Seventh Conference on Learning Theory}, volume 247 of {\em Proceedings of
  Machine Learning Research}, pages 1796--1881. PMLR, 2024.

\bibitem[GM11]{GM11:journal}
Olivier Gu{\'e}don and Emanuel Milman.
\newblock Interpolating thin-shell and sharp large-deviation estimates for
  isotropic log-concave measures.
\newblock {\em Geometric and Functional Analysis}, 21(5):1043--1068, 2011.

\bibitem[Gua24]{G24:arxiv}
Qingyang Guan.
\newblock A note on bourgain's slicing problem, 2024.

\bibitem[JC25]{JC25}
Minhui Jiang and Yuansi Chen.
\newblock Regularized {Dikin} walks for sampling truncated log-concave
  measures, mixed isoperimetry and beyond worst-case analysis.
\newblock In {\em Proceedings of Thirty Eighth Conference on Learning Theory},
  volume 291 of {\em Proceedings of Machine Learning Research}, pages
  3017--3078. PMLR, 2025.

\bibitem[JLLV26]{JLLV26:journal}
He~Jia, Aditi Laddha, Yin~Tat Lee, and Santosh Vempala.
\newblock Reducing isotropy and volume to kls: Faster rounding and volume
  algorithms.
\newblock {\em J. ACM}, 73(2), April 2026.

\bibitem[JLV22]{JLV22:arxiv}
Arun Jambulapati, Yin~Tat Lee, and Santosh~S. Vempala.
\newblock A slightly improved bound for the {KLS} constant, 2022.

\bibitem[KL22]{KL22:journal}
Bo'az Klartag and Joseph Lehec.
\newblock Bourgain's slicing problem and {KLS} isoperimetry up to polylog.
\newblock {\em Geometric and Functional Analysis}, 32(5):1134--1159, 2022.

\bibitem[KL25a]{KL25:journal}
Bo'az Klartag and Joseph Lehec.
\newblock Affirmative resolution of bourgain's slicing problem using guan's
  bound.
\newblock {\em Geometric and Functional Analysis}, 35:1147--1168, 2025.

\bibitem[KL25b]{KL25:arxiv}
Boaz Klartag and Joseph Lehec.
\newblock Thin-shell bounds via parallel coupling.
\newblock {\em arXiv preprint arXiv:2507.15495}, 2025.

\bibitem[Kla06]{K06:journal}
B.~Klartag.
\newblock On convex perturbations with a bounded isotropic constant.
\newblock {\em Geometric and Functional Analysis}, 16:1274--1290, 2006.

\bibitem[Kla23]{K23:journal}
Bo'az Klartag.
\newblock Logarithmic bounds for isoperimetry and slices of convex sets.
\newblock {\em Ars Inveniendi Analytica}, pages Paper No.~4, 17 pp., 2023.

\bibitem[KLS95]{KLS95:journal}
R.~Kannan, L.~Lov\'asz, and M.~Simonovits.
\newblock Isoperimetric problems for convex bodies and a localization lemma.
\newblock {\em Discrete Comput. Geom.}, 13(3-4):541--559, 1995.

\bibitem[KLS97]{KLS97:journal}
Ravi Kannan, L{\'a}szl{\'o} Lov{\'a}sz, and Mikl{\'o}s Simonovits.
\newblock Random walks and an {$O^*(n^5)$} volume algorithm for convex bodies.
\newblock {\em Random Structures \& Algorithms}, 11(1):1--50, 1997.

\bibitem[KLSV22]{KLSV22}
Yunbum Kook, Yin-Tat Lee, Ruoqi Shen, and Santosh Vempala.
\newblock Sampling with riemannian hamiltonian monte carlo in a constrained
  space.
\newblock In S.~Koyejo, S.~Mohamed, A.~Agarwal, D.~Belgrave, K.~Cho, and A.~Oh,
  editors, {\em Advances in Neural Information Processing Systems}, volume~35,
  pages 31684--31696. Curran Associates, Inc., 2022.

\bibitem[KLSV23]{KLSV23}
Yunbum Kook, Yin~Tat Lee, Ruoqi Shen, and Santosh Vempala.
\newblock Condition-number-independent convergence rate of riemannian
  hamiltonian monte carlo with numerical integrators.
\newblock In Gergely Neu and Lorenzo Rosasco, editors, {\em Proceedings of
  Thirty Sixth Conference on Learning Theory}, volume 195 of {\em Proceedings
  of Machine Learning Research}, pages 4504--4569. PMLR, 2023.

\bibitem[KN12]{KN12:journal}
Ravindran Kannan and Hariharan Narayanan.
\newblock Random walks on polytopes and an affine interior point method for
  linear programming.
\newblock {\em Mathematics of Operations Research}, 37(1):1--20, 2012.

\bibitem[KV25]{KV25}
Yunbum Kook and Santosh~S. Vempala.
\newblock Sampling and integration of log-concave functions by algorithmic
  diffusion.
\newblock In Michal Kouck{\'{y}} and Nikhil Bansal, editors, {\em Proceedings
  of the 57th Annual {ACM} Symposium on Theory of Computing, {STOC} 2025,
  Prague, Czechia, June 23-27, 2025}, pages 924--932. {ACM}, 2025.
\newblock Full version (used for theorem numbering):
\url{https://arxiv.org/abs/2411.13462v1}.

\bibitem[KZ24]{KZ24:arxiv}
Yunbum Kook and Matthew~S. Zhang.
\newblock Covariance estimation using {M}arkov chain {M}onte {C}arlo, 2024.
\newblock \url{https://arxiv.org/abs/2410.17147v1}.

\bibitem[Let26]{L26:arxiv}
Brayden Letwin.
\newblock The {KLS} constant is {$O(\log^{1/4} n)$}, 2026.

\bibitem[LLV20]{LLV20}
Aditi Laddha, Yin~Tat Lee, and Santosh~S. Vempala.
\newblock Strong self-concordance and sampling.
\newblock In {\em Proceedings of the 52nd Annual ACM SIGACT Symposium on Theory
  of Computing}, pages 1212--1222, 2020.

\bibitem[LS93]{LS93:journal}
L{\'a}szl{\'o} Lov{\'a}sz and Mikl{\'o}s Simonovits.
\newblock Random walks in a convex body and an improved volume algorithm.
\newblock {\em Random Structures \& Algorithms}, 4(4):359--412, 1993.

\bibitem[LV06a]{LV06b:journal}
L{\'a}szl{\'o} Lov{\'a}sz and Santosh~S. Vempala.
\newblock Hit-and-run from a corner.
\newblock {\em SIAM Journal on Computing}, 35(4):985--1005, 2006.

\bibitem[LV06b]{LV06c:journal}
L{\'a}szl{\'o} Lov{\'a}sz and Santosh~S. Vempala.
\newblock Simulated annealing in convex bodies and an {$O^*(n^4)$} volume
  algorithm.
\newblock {\em Journal of Computer and System Sciences}, 72(2):392--417, 2006.

\bibitem[LV07]{LV07:journal}
L{\'a}szl{\'o} Lov{\'a}sz and Santosh~S. Vempala.
\newblock The geometry of log-concave functions and sampling algorithms.
\newblock {\em Random Structures \& Algorithms}, 30(3):307--358, 2007.

\bibitem[LV18]{LV18}
Yin~Tat Lee and Santosh~S. Vempala.
\newblock {Convergence Rate of Riemannian Hamiltonian Monte Carlo and Faster
  Polytope Volume Computation}.
\newblock In {\em Proceedings of the 50th Annual ACM SIGACT Symposium on Theory
  of Computing}, pages 1115--1121, 2018.

\bibitem[LV22]{LV22:journal}
Yin~Tat Lee and Santosh~S. Vempala.
\newblock Geodesic walks in polytopes.
\newblock {\em SIAM Journal on Computing}, 51(2):STOC17--400--STOC17--488,
  2022.
\newblock Conference version in STOC 2017.

\bibitem[LV24]{LV24:journal}
Yin~Tat Lee and Santosh~S. Vempala.
\newblock Eldan's stochastic localization and the {KLS} conjecture:
  Isoperimetry, concentration and mixing.
\newblock {\em Annals of Mathematics}, 199(3):1043--1092, 2024.

\bibitem[MV23]{MV23}
Oren Mangoubi and Nisheeth~K. Vishnoi.
\newblock Sampling from structured log-concave distributions via a
  soft-threshold dikin walk.
\newblock In Alice Oh, Tristan Naumann, Amir Globerson, Kate Saenko, Moritz
  Hardt, and Sergey Levine, editors, {\em Advances in Neural Information
  Processing Systems}, volume~36, pages 31908--31942. Curran Associates, Inc.,
  2023.

\bibitem[Nar16]{N16:journal}
Hariharan Narayanan.
\newblock Randomized interior point methods for sampling and optimization.
\newblock {\em The Annals of Applied Probability}, 26(1):597--641, 2016.

\bibitem[NR17]{NR17:journal}
Hariharan Narayanan and Alexander Rakhlin.
\newblock Efficient sampling from time-varying log-concave distributions.
\newblock {\em Journal of Machine Learning Research}, 18(112):1--29, 2017.

\bibitem[Pao06]{P06:journal}
G~Paouris.
\newblock Concentration of mass on convex bodies.
\newblock {\em GAFA Geom. Funct. Anal.}, 16(5):1021--1049, December 2006.

\bibitem[Pao12]{P12:journal}
Grigoris Paouris.
\newblock Small ball probability estimates for log-concave measures.
\newblock {\em Transactions of the American Mathematical Society},
  364(1):287--308, 2012.

\bibitem[Pr{\'e}73]{P73:journal}
Andr{\'a}s Pr{\'e}kopa.
\newblock On log-concave measures and functions.
\newblock {\em Acta Sci. Math. (Szeged)}, 34:335--343, 1973.

\bibitem[SV16]{SV16:journal}
Sushant Sachdeva and Nisheeth~K. Vishnoi.
\newblock The mixing time of the dikin walk in a polytope---a simple proof.
\newblock {\em Operations Research Letters}, 44(5):630--634, 2016.

\bibitem[Tro15]{T15:journal}
Joel~A. Tropp.
\newblock An introduction to matrix concentration inequalities.
\newblock {\em Foundations and Trends in Machine Learning}, 8(1--2):1--230,
  2015.

\end{thebibliography}
% ===== end bibliography =====

\appendix

% ===== begin a1_asym_sharp.tex =====
\section{Proof of the Sharp Trace Growth Theorem}
\label{sec:sharp}

This section proves \cref{thm:sharp-trace-lower}, which shows that the trace
growth bounds used in the $\O(n^{13/4})$ analysis of the asymmetric algorithm
in \cref{sec:asymmetric} are sharp up to constants.
We construct restricted Gaussian examples where anisotropic warming creates
trace growth of the same order as those bounds.

\subsection{The lower-bound construction}
\label{app:warming:sec:lower-bound}

We now construct the matching family.
The mechanism is a temperature-driven switch in a one-dimensional control
variable.
Here a fiber means the slice of the construction obtained by fixing this control
variable and looking at the remaining coordinates.
One collection of such slices is temperature-sensitive and expands with the
control variable; another is almost temperature-insensitive and contracts with
it.
The parameters are chosen so that these two contributions balance at the
original temperature but move in opposite directions after warming.

\subsubsection{A truncated-Gaussian response function}
\label{app:warming:sec:response-function}

For $t>0$ and
$r>0$, set
\begin{equation}
J_t(r)=\int_{-r}^{r}e^{-v^2/(2t)}\,\mathrm dv,
\qquad
D_t(r)=r\,\partial_r\log J_t(r).
\label{app:warming:eq:J-D}
\end{equation}
Writing $D=D_1$, scaling gives
\begin{equation}
  D_t(r)=D(r/t^{1/2}).
  \label{app:warming:eq:D-scaling}
\end{equation}

The next lemma concerns the temperature response.

\begin{lemma}
\label{app:warming:lem:temperature-response}
The function $D$ is strictly decreasing on
$(0,\infty)$ and
\begin{equation}
  D(r)=1-\frac{r^2}{3}+O(r^4)
  \qquad (r\to0).
  \label{app:warming:eq:D-expansion}
\end{equation}
Consequently, one may fix a sufficiently small universal $b>0$ and constants
$c_0,C_0>0$ such that, for every
$T=1+\delta\in[1,2]$,
\begin{equation}
c_0\delta
\leq D_T(1)D_1(b)-D_1(1)D_T(b)
\leq C_0\delta.
\label{app:warming:eq:temperature-response}
\end{equation}
\end{lemma}

\begin{proof}
Direct differentiation gives
\begin{equation}
  D'(r)=\frac{D(r)}{r}\bigl(1-r^2-D(r)\bigr).
  \label{app:warming:eq:D-derivative}
\end{equation}
We claim that $D(r)>1-r^2$.
This is immediate for $r\geq1$.
For $0<r<1$, let
\[
  H(r)=2re^{-r^2/2}-(1-r^2)J_1(r).
\]
Then $H(0)=0$ and
$H'(r)=2rJ_1(r)>0$.
Since $H(r)=J_1(r)[D(r)-(1-r^2)]$, the claim follows.
It follows from Equation~\eqref{app:warming:eq:D-derivative} that $D'(r)<0$.
Expanding the numerator and denominator in the definition of $D$ gives
Equation~\eqref{app:warming:eq:D-expansion}.

By Equation~\eqref{app:warming:eq:D-scaling},
\[
  \partial_tD_t(1)=-\frac{1}{2t^{3/2}}D'(t^{-1/2}),
\]
which has a strictly positive minimum for $1\leq t\leq2$.
Hence $D_T(1)-D_1(1)\asymp\delta$.
On the other hand, Equation~\eqref{app:warming:eq:D-expansion} gives
\[
D_t(b)=1-\frac{b^2}{3t}+O(b^4),
\qquad
D_T(b)-D_1(b)=O(b^2\delta),
\]
uniformly for $1\leq t,T\leq2$; here the remainder is smooth in
$t$ and its
$t$-derivative is
$O(b^4)$.
Taking $b$ sufficiently small proves
Equation~\eqref{app:warming:eq:temperature-response}.
\end{proof}

\subsubsection{The polytope and its phase switch}
\label{app:warming:sec:polytope}

Assume first that $n\geq8$.
Let
\begin{equation}
m=\lfloor n/4\rfloor,
\qquad
r=n-2m-2.
\label{app:warming:eq:m-r}
\end{equation}
Then $m\geq n/8$ and
$r\geq n/4$.
Use coordinates
\[
  (x,y_1,\ldots,y_m,z_1,\ldots,z_m,w,v_1,\ldots,v_r).
\]

Fix the constant $b$ from \cref{app:warming:lem:temperature-response}.
This lemma gives a uniform linear separation between the temperature responses
at $1$ and
$T=1+\delta$, which is used to choose
$\eta$ below.
Let $\kappa>0$ be a sufficiently small universal constant, define
\begin{equation}
\varepsilon=\kappa\delta,
\qquad
\eta=\varepsilon\frac{D_1(1)+D_T(1)}{D_1(b)+D_T(b)},
\qquad
T=1+\delta.
\label{app:warming:eq:epsilon-eta}
\end{equation}
Uniformly in $0\leq\delta\leq1$, we have
$\eta\asymp\varepsilon$; after decreasing
$\kappa$, we may assume
$0\leq\eta\leq1/2$.

For $s\in[0,1]$, write
$u=x+s$ and define the polytope
$K_s$ by
\begin{equation}
  \begin{aligned}
    0&\leq u\leq1,\\
    |y_i|&\leq1+\varepsilon u &&(1\leq i\leq m),\\
    |z_i|&\leq b(1-\eta u) &&(1\leq i\leq m),\\
    |w|&\leq u,\\
    |v_j|&\leq1 &&(1\leq j\leq r).
  \end{aligned}
  \label{app:warming:eq:polytope}
\end{equation}
All constraints are affine, so $K_s$ is a full-dimensional compact convex
polytope.

Before whitening, take the ambient Gaussian covariance to be $tI_n$.
The $u$-marginal at scale
$t$ is proportional on
$[0,1]$ to
\begin{equation}
  H_{t,s}(u)e^{F_t(u)},
  \label{app:warming:eq:u-marginal}
\end{equation}
where
\begin{equation}
  H_{t,s}(u)
  =e^{-(u-s)^2/(2t)}\int_{-u}^{u}e^{-w^2/(2t)}\,\mathrm dw,
  \label{app:warming:eq:H}
\end{equation}
and
\begin{equation}
  F_t(u)
  =m\log J_t(1+\varepsilon u)
   +m\log J_t\bigl(b(1-\eta u)\bigr).
  \label{app:warming:eq:F}
\end{equation}
For $s,u\in[0,1]$ and
$t\in[1,2]$,
\begin{equation}
c_1u
\leq H_{t,s}(u)
\leq C_1u.
\label{app:warming:eq:H-comparison}
\end{equation}

The next lemma concerns the uniform slope reversal.

\begin{lemma}
\label{app:warming:lem:slope-reversal}
There are universal constants $c,C>0$ such that, with
$\kappa$ chosen sufficiently small and
\begin{equation}
  \lambda=m\kappa\delta^2,
  \label{app:warming:eq:lambda}
\end{equation}
one has, for every $0\leq u\leq1$,
\begin{equation}
-C\lambda
\leq F_1'(u)
\leq-c\lambda,
\qquad
c\lambda
\leq F_T'(u)
\leq C\lambda.
\label{app:warming:eq:slope-reversal}
\end{equation}
\end{lemma}

\begin{proof}
Let $L_t(q)=\partial_q\log J_t(q)=D_t(q)/q$.
Differentiating Equation~\eqref{app:warming:eq:F} gives
\begin{equation}
  F_t'(u)
  =m\Bigl[
    \varepsilon L_t(1+\varepsilon u)
    -b\eta L_t\bigl(b(1-\eta u)\bigr)
  \Bigr].
  \label{app:warming:eq:F-derivative}
\end{equation}
At $u=0$,
\[
  F_t'(0)=m\bigl(\varepsilon D_t(1)-\eta D_t(b)\bigr).
\]
The choice of $\eta$ gives
\begin{equation}
F_T'(0)
=-F_1'(0)
=m\varepsilon \frac{D_T(1)D_1(b)-D_1(1)D_T(b)}{D_1(b)+D_T(b)}.
\label{app:warming:eq:F-zero}
\end{equation}
By \cref{app:warming:lem:temperature-response}, the right-hand side is between
$cm\kappa\delta^2$ and
$Cm\kappa\delta^2$.

The radii in Equation~\eqref{app:warming:eq:F-derivative} stay in fixed compact
subintervals of $(0,\infty)$.
Differentiating once more and using $\eta\asymp\varepsilon$ gives
\begin{equation}
|F_t'(u)-F_t'(0)|
\leq Cm(\varepsilon^2+\eta^2)
\leq Cm\kappa^2\delta^2.
\label{app:warming:eq:F-variation}
\end{equation}
Choosing $\kappa$ sufficiently small relative to the lower constant in
Equation~\eqref{app:warming:eq:F-zero} proves
Equation~\eqref{app:warming:eq:slope-reversal}.
\end{proof}

\subsubsection{Centering, endpoint concentration, and whitening}
\label{app:warming:sec:whitening}

We first choose $s$ so that the
scale-$1$ law is centered.
Let $\nu_s$ be the standard Gaussian restricted to
$K_s$, and define
\[
  \Phi(s)=\E{u\sim\nu_s}{u}-s.
\]
The function $\Phi$ is continuous.
Since $u>0$ almost surely when
$s=0$ and
$u<1$ almost surely when
$s=1$, we have
$\Phi(0)>0$ and
$\Phi(1)<0$.
Therefore some $s\in(0,1)$ satisfies
\begin{equation}
  s=\E{u\sim\nu_s}{u}.
  \label{app:warming:eq:centering-s}
\end{equation}
For this choice, $\E{x\sim\nu_s}{x}=\E{u\sim\nu_s}{u-s}=0$.
Every other coordinate has mean zero by sign symmetry.
The same conditional sign symmetries imply that all cross-covariances vanish.
Thus the scale-$1$ covariance matrix
$C$ is positive definite and diagonal.

For the chosen $s$, write
$\mathbb P_t$ for probability under the law
$\mu_{K_s,tI_n}$ obtained by restricting
$N(0,tI_n)$ to
$K_s$.

Suppose first that $\lambda\geq1$.
From Equation~\eqref{app:warming:eq:H-comparison} and the first pair of
inequalities in Equation~\eqref{app:warming:eq:slope-reversal},
\begin{equation}
  \begin{aligned}
    \E{u\sim\mu_{K_s,I_n}}{u^2}
    &=\frac{\displaystyle
       \int_0^1 u^2H_{1,s}(u)e^{F_1(u)}\,\mathrm du}
      {\displaystyle
       \int_0^1 H_{1,s}(u)e^{F_1(u)}\,\mathrm du}\\
    &\leq C
      \frac{\displaystyle\int_0^\infty u^3e^{-c\lambda u}\,\mathrm du}
           {\displaystyle\int_0^{1/\lambda}ue^{-C\lambda u}\,\mathrm du}
     \leq\frac{C}{\lambda^2}.
  \end{aligned}
  \label{app:warming:eq:u-small}
\end{equation}
At scale $T$, Equation~\eqref{app:warming:eq:slope-reversal} and
Equation~\eqref{app:warming:eq:H-comparison} show, by comparing
$[0,1/2]$ with
$[3/4,1]$, that
\begin{equation}
\mathbb P_T(u\geq1/2)\geq c,
\qquad
\E{u\sim\mu_{K_s,TI_n}}{u^2}\geq c.
\label{app:warming:eq:u-large}
\end{equation}

Conditionally on $u$, the
$w$-coordinate is the centered
$N(0,t)$ law restricted to
$[-u,u]$.
Uniformly for $0\leq u\leq1$ and
$1\leq t\leq2$,
\begin{equation}
cu^2
\leq\E{w\sim\mu_{K_s,tI_n}}{w^2\mid u}
\leq Cu^2.
\label{app:warming:eq:w-conditional}
\end{equation}
Since $\E{w\sim\mu_{K_s,tI_n}}{w\mid u}=0$, Equations
\eqref{app:warming:eq:u-small}--\eqref{app:warming:eq:w-conditional} imply
\begin{equation}
\frac{\V{w\sim\mu_{K_s,TI_n}}{w}}{\V{w\sim\mu_{K_s,I_n}}{w}}
\geq c\lambda^2
\geq cn^2\delta^4.
\label{app:warming:eq:w-variance-ratio}
\end{equation}
Here and below the universal constants absorb the fixed $\kappa$ and the fact
that
$m\geq n/8$.

The filler coordinates $v_j$ are independent of the gadget.
At scales $1\leq t\leq2$, their one-dimensional densities on
$[-1,1]$ are uniformly comparable.
Consequently,
\begin{equation}
  \frac{\V{v_j\sim\mu_{K_s,TI_n}}{v_j}}{\V{v_j\sim\mu_{K_s,I_n}}{v_j}}\geq c.
  \label{app:warming:eq:v-variance-ratio}
\end{equation}

Now apply the linear whitening map
\begin{equation}
L=C^{-1/2},
\qquad
\widetilde K=LK_s,
\qquad
\widetilde Q
=LI_nL^{\mathsf T}
=C^{-1}.
\label{app:warming:eq:whitening}
\end{equation}
If $\widetilde X=LX$, its
scale-$1$ law is exactly
$\mu_{\widetilde K,\widetilde Q}$, with mean zero and covariance
$I_n$.
At scale $T$, its law is exactly
$\mu_{\widetilde K,T\widetilde Q}$.
No translation has been applied, so the ambient Gaussian remains centered at the
origin.
Since $C$ and the warmed covariance are diagonal, Equations
\eqref{app:warming:eq:w-variance-ratio} and
\eqref{app:warming:eq:v-variance-ratio} give
\begin{equation}
\operatorname{Tr}\operatorname{Cov} (\mu_{\widetilde K,T\widetilde Q})
\geq cr+c\lambda^2
\geq c\bigl(n+n^2\delta^4\bigr)
\label{app:warming:eq:trace-lower-bound}
\end{equation}
whenever $\lambda\geq1$.

If $\lambda<1$, the filler coordinates still contribute at least
$cr\geq cn$.
Moreover, $m\geq n/8$ and Equation~\eqref{app:warming:eq:lambda} imply
\[
n^2\delta^4
=\frac{n^2\lambda^2}{m^2\kappa^2}
\leq\frac{64}{\kappa^2},
\]
which is a universal constant and can be absorbed into a constant multiple of
$n$.
Thus Equation~\eqref{app:warming:eq:trace-lower-bound} remains valid, after
decreasing its universal constant.

For $1\leq n<8$, take a product of one-dimensional symmetric truncated Gaussians
on
$[-1,1]$ and apply scalar whitening.
Warming by a factor in $[1,2]$ leaves each variance bounded below by a universal
constant.
Since $n+n^2\delta^4\leq8n$ in these dimensions, the lower bound follows after
one last adjustment of the constant.
This completes the proof of \cref{thm:sharp-trace-lower}.
% ===== end a1_asym_sharp.tex =====
% ===== begin a2_counter_coordinate_moment.tex =====
\section{Counterexamples to The Coordinate Moment Conjecture}
\label{sec:counter_coordinate_moment}

We prove \cref{thm:coordinate-moment-counterexample}.
Throughout, $x':=(x_2,x_3,\ldots,x_n)$ denotes all but the first coordinate of
$x\in\R^n$.
The counterexample density is a linear image of the restriction of a
multivariate Gaussian to a symmetric convex body, defined next.

Assume $n\ge7$.
Define
\[
g(x):=
\begin{cases}
(2\pi)^{-(n-1)/2}\exp(-\norm{x'}^2/2)&\text{if }|x_1|\le1\text{ and }\norm{x'}\le(n-2)^{1/2}-|x_1|,\\
0&\text{otherwise.}
\end{cases}
\]
Also, for $|x_1|\le1$, define
\[
g_1(x_1)
:=\int_{\R^{n-1}}g(x)\,dx'
=\int_0^{(n-2)^{1/2}-|x_1|} \frac{1}{2^{(n-1)/2-1}\Gamma((n-1)/2)}t^{n-2}e^{-t^2/2}\,dt
\]
\[
g_{\mathrm{sq}}(x_1)
:=\int_{\R^{n-1}}\norm{x'}^2g(x)\,dx'
=\int_0^{(n-2)^{1/2}-|x_1|} t^2\cdot\frac{1}{2^{(n-1)/2-1}\Gamma((n-1)/2)}t^{n-2}e^{-t^2/2}\,dt
\]
where the last equality uses polar coordinates in $x'$: if
$Z\sim N(0,I_{n-1})$ and
$T=\norm{Z}$, then
$T$ has density
\[
\frac{1}{2^{(n-1)/2-1}\Gamma((n-1)/2)}t^{n-2}e^{-t^2/2}.
\]
Observe that $g_{\mathrm{sq}}(x_1)/g_1(x_1)$ is a decreasing function of
$|x_1|$.

Defining
\[
G
:=\int_{\R^n}g(x)\,dx
=\int_{-1}^1g_1(x_1)\,dx_1
\]
we have that $g(x)/G$ is a probability density function.
It is even and log-concave, being the standard Gaussian density in
$x'$ restricted to the symmetric convex body
$\{|x_1|\le1,\ \norm{x'}\le(n-2)^{1/2}-|x_1|\}$.

The next lemma concerns the moments of $g/G$ and its isotropic linear image.

\begin{lemma}
\label{lem:coordinate-moment-almost-isotropic}
Let $X\sim g/G$.
Then $\E{X\sim g/G}{X_1^2}=\Theta(1)$ and
$\E{X\sim g/G}{\norm{X'}^2}=\Theta(n)$.
With
\[
c_1:=\E{X\sim g/G}{X_1^2}^{1/2},
\qquad
c':=\left(\frac{\E{X\sim g/G}{\norm{X'}^2}}{n-1}\right)^{1/2},
\]
both of order $\Theta(1)$, the density
\[
f^\star(x):=c_1(c')^{n-1}\bigl(g(c_1x_1,c'x')/G\bigr)
\]
of $(X_1/c_1,X'/c')$ is even, isotropic and log-concave, and
\[
\E{x\sim f^\star}{x_1^2\bigl(\norm{x'}^2-(n-1)\bigr)}
=\frac{1}{c_1^2(c')^2}
\E{X\sim g/G}{\Bigl(X_1^2-\E{X\sim g/G}{X_1^2}\Bigr)\norm{X'}^2}.
\]
\end{lemma}

The next lemma concerns the covariance of $X_1^2$ and
$\norm{X'}^2$ under
$g/G$.

\begin{lemma}
\label{lem:coordinate-moment-covariance}
Let $X\sim g/G$.
Then
\[
\E{X\sim g/G}{\Bigl(X_1^2-\E{X\sim g/G}{X_1^2}\Bigr)\norm{X'}^2}
=-\Theta(n^{1/2}).
\]
\end{lemma}

\begin{proof}[Proof of \cref{thm:coordinate-moment-counterexample}]
Let $f=f^\star$ from \cref{lem:coordinate-moment-almost-isotropic}; it is even,
isotropic and log-concave.
Since $\norm{x}^2-n=(\norm{x'}^2-(n-1))+(x_1^2-1)$ and
$\E{x\sim f}{x_1^2}=1$,
\[
\E{x\sim f}{x_1^2\bigl(\norm{x}^2-n\bigr)}
=\E{x\sim f}{x_1^2\bigl(\norm{x'}^2-(n-1)\bigr)}
+\E{x\sim f}{x_1^4}-1,
\]
and $\E{x\sim f}{x_1^4}\le C_4$ by \cref{prop:log-concave-toolkit}(iii).
The first term on the right is $-\Theta(n^{1/2})/(c_1^2(c')^2)=-\Theta(n^{1/2})$ by
\cref{lem:coordinate-moment-almost-isotropic,lem:coordinate-moment-covariance},
so the left side has absolute value at least
$cn^{1/2}$ for a universal
$c>0$ once
$n$ is sufficiently large.
\end{proof}

\begin{proof}[Proof of \cref{lem:coordinate-moment-almost-isotropic}]
We first show that $g_1(1)/g_1(0)>\Omega(1)$ and
$g_{\mathrm{sq}}(0)/g_1(0),g_{\mathrm{sq}}(1)/g_1(1)=\Theta(n)$.
\begin{enumerate}[label=(\roman*)]
\item To show $g_1(1)/g_1(0)>\Omega(1)$, it suffices to show
\[
\frac{\int_{(n-2)^{1/2}-2}^{(n-2)^{1/2}-1}t^{n-2}e^{-t^2/2}\,dt}
{\int_{(n-2)^{1/2}-1}^{(n-2)^{1/2}}t^{n-2}e^{-t^2/2}\,dt}>\Omega(1).
\]
But we compute
\begin{align*}
\frac{(t-1)^{n-2}e^{-(t-1)^2/2}}{t^{n-2}e^{-t^2/2}}
&=\left(1-\frac1t\right)^{n-2}e^{t-(1/2)}\\
&=\exp\left(-\frac{n-2}{t}-O\left(\frac{n-2}{t^2}\right)+t-\frac12\right)\\
&=\exp\left(-\frac{((n-2)^{1/2}-t)((n-2)^{1/2}+t)}{t}-O(1)\right)\\
&=\exp(-O(1))\\
&=\Omega(1)
\end{align*}
for $t\in[(n-2)^{1/2}-1,(n-2)^{1/2}]$.
Integrating yields
\[
\int_{(n-2)^{1/2}-2}^{(n-2)^{1/2}-1}t^{n-2}e^{-t^2/2}\,dt
>\Omega(1)\cdot\int_{(n-2)^{1/2}-1}^{(n-2)^{1/2}}t^{n-2}e^{-t^2/2}\,dt,
\]
as desired.

\item To show $g_{\mathrm{sq}}(0)/g_1(0)=O(n)$, we note the trivial bound
\[
\frac{\int_0^{(n-2)^{1/2}}t^2\cdot t^{n-2}e^{-t^2/2}\,dt} {\int_0^{(n-2)^{1/2}}t^{n-2}e^{-t^2/2}\,dt}
\le((n-2)^{1/2})^2
<n.
\]

\item To show $g_{\mathrm{sq}}(1)/g_1(1)=\Omega(n)$, we note that
\[
\frac{\int_0^{(n-2)^{1/2}-1}t^2\cdot t^{n-2}e^{-t^2/2}\,dt}
{\int_0^{(n-2)^{1/2}-1}t^{n-2}e^{-t^2/2}\,dt}
\ge
\left(
\frac{\int_0^{(n-2)^{1/2}-1}t\cdot t^{n-2}e^{-t^2/2}\,dt}
{\int_0^{(n-2)^{1/2}-1}t^{n-2}e^{-t^2/2}\,dt}
\right)^2
\]
so it suffices to show
\[
\frac{\int_0^{(n-2)^{1/2}-1}t\cdot t^{n-2}e^{-t^2/2}\,dt}
{\int_0^{(n-2)^{1/2}-1}t^{n-2}e^{-t^2/2}\,dt}=\Omega(n^{1/2}).
\]
Now, we calculate
\[
\log(t^{n-2}e^{-t^2/2})=(n-2)\log t-\frac{t^2}{2}
\]
is a concave function of $t$, and
\[
\frac{d}{dt}\log(t^{n-2}e^{-t^2/2})
=\frac{n-2}{t}-t
>\frac{n-2}{(n-2)^{1/2}}-t
\ge1
\]
for $t\in[0,(n-2)^{1/2}-1]$.
It follows (by a comparison argument) that
\[
\frac{\int_0^{(n-2)^{1/2}-1}t\cdot t^{n-2}e^{-t^2/2}\,dt} {\int_0^{(n-2)^{1/2}-1}t^{n-2}e^{-t^2/2}\,dt}
\ge \frac{\int_{-\infty}^{(n-2)^{1/2}-1}t\cdot e^t\,dt} {\int_{-\infty}^{(n-2)^{1/2}-1}e^t\,dt}
=(n-2)^{1/2}-2
=\Omega(n^{1/2})
\]
as desired.
\end{enumerate}

The density of $X_1$ is
$g_1/G$ on
$[-1,1]$, and
$g_1$ is nonincreasing in
$|x_1|$ with
$g_1(1)/g_1(0)>\Omega(1)$, so this density is within constant factors of the
uniform density on
$[-1,1]$ and
$\E{X\sim g/G}{X_1^2}=\Theta(1)$.
Conditionally on $X_1=x_1$, the mean of
$\norm{X'}^2$ is
$g_{\mathrm{sq}}(x_1)/g_1(x_1)$, which lies between
$g_{\mathrm{sq}}(1)/g_1(1)=\Omega(n)$ and
$g_{\mathrm{sq}}(0)/g_1(0)=O(n)$; hence
$\E{X\sim g/G}{\norm{X'}^2}=\Theta(n)$, and
$c_1,c'=\Theta(1)$.
The law $g/G$ is invariant under $x_1\mapsto-x_1$ and under every rotation of
$x'$, so its mean is
$0$ and its covariance is
$\operatorname{diag}(c_1^2,(c')^2I_{n-1})$; therefore
$f^\star$, the density of
$(X_1/c_1,X'/c')$, has mean
$0$ and covariance
$I_n$, and it is even and log-concave as a linear image of
$g/G$.
Finally,
\[
\E{x\sim f^\star}{x_1^2\bigl(\norm{x'}^2-(n-1)\bigr)}
=\frac{1}{c_1^2(c')^2}
\E{X\sim g/G}{X_1^2\bigl(\norm{X'}^2-(c')^2(n-1)\bigr)},
\]
and $(c')^2(n-1)=\E{X\sim g/G}{\norm{X'}^2}$, so subtracting the constant
$\E{X\sim g/G}{X_1^2}\,\E{X\sim g/G}{\norm{X'}^2}$ in both forms gives the
stated identity.
\end{proof}

The proof of \cref{lem:coordinate-moment-covariance} uses the following
derivative bound for truncated second moments of a monotone log-concave
density.

\begin{lemma}
\label{lem:aux-logconcave-ratio}
Let $L_0>0$, and let
$f:[0,+\infty)\to\R$ be a log-concave probability density function which is
continuous, such that
$f$ is nondecreasing on
$[0,L_0]$ and
$\int_0^{L_0}f(t)\,dt>0$.
Then
\[
\frac12L_0
\le \left.\frac{d}{dL}\frac{\int_0^L t^2\cdot f(t)\,dt}{\int_0^L f(t)\,dt}\right|_{L=L_0}
\le2L_0.
\]
\end{lemma}
\begin{proof}
We compute
\begin{align*}
\frac{d}{dL}\frac{\int_0^L t^2\cdot f(t)\,dt}{\int_0^L f(t)\,dt}
&=\frac{L^2\cdot f(L)}{\int_0^Lf(t)\,dt}
-\frac{f(L)\int_0^Lt^2\cdot f(t)\,dt}{\left(\int_0^Lf(t)\,dt\right)^2}\\
&=\frac{f(L)\int_0^L(L^2-t^2)\cdot f(t)\,dt}{\left(\int_0^Lf(t)\,dt\right)^2}
\end{align*}
and thus
\[
\left.\frac{d}{dL}\frac{\int_0^L t^2\cdot f(t)\,dt}{\int_0^L f(t)\,dt}\right|_{L=L_0}
=\frac{f(L_0)\int_0^{L_0}(L_0^2-t^2)\cdot f(t)\,dt}{\left(\int_0^{L_0}f(t)\,dt\right)^2}.
\]
For the lower bound, we have
\begin{align*}
\left(\int_0^{L_0}f(t)\,dt\right)^2
&=2\int_0^{L_0}\int_t^{L_0}f(t)f(s)\,ds\,dt\\
&\le2\int_0^{L_0}\int_t^{L_0}f(t)f(L_0)\,ds\,dt \qquad\text{($f$ is nondecreasing on $[0,L_0]$)}\\
&=2f(L_0)\int_0^{L_0}(L_0-t)\cdot f(t)\,dt\\
&\le2f(L_0)\int_0^{L_0}(L_0-t)\frac{L_0+t}{L_0}\cdot f(t)\,dt\\
&=\frac2{L_0}f(L_0)\int_0^{L_0}(L_0^2-t^2)\cdot f(t)\,dt
\end{align*}
which implies
\[
\left.\frac{d}{dL}\frac{\int_0^L t^2\cdot f(t)\,dt}{\int_0^L f(t)\,dt}\right|_{L=L_0}
=\frac{f(L_0)\int_0^{L_0}(L_0^2-t^2)\cdot f(t)\,dt}{\left(\int_0^{L_0}f(t)\,dt\right)^2}
\ge\frac12L_0,
\]
as desired.

For the upper bound, we have
\begin{align*}
\left(\int_0^{L_0}f(t)\,dt\right)^2
&=\int_0^{L_0}\int_0^{L_0}f(t)f(s)\,ds\,dt\\
&\ge\iint_{t,s\in[0,L_0],\ t+s\ge L_0}f(t)f(s)\,ds\,dt\\
&\ge\iint_{t,s\in[0,L_0],\ t+s\ge L_0}f(L_0)f(t+s-L_0)\,ds\,dt \qquad\text{(log-concavity of $f$)}\\
&=f(L_0)\int_0^{L_0}(L_0-t)\cdot f(t)\,dt\\
&\ge f(L_0)\int_0^{L_0}(L_0-t)\frac{L_0+t}{2L_0}\cdot f(t)\,dt\\
&=\frac1{2L_0}f(L_0)\int_0^{L_0}(L_0^2-t^2)\cdot f(t)\,dt
\end{align*}
which implies
\[
\left.\frac{d}{dL}\frac{\int_0^L t^2\cdot f(t)\,dt}{\int_0^L f(t)\,dt}\right|_{L=L_0}
=\frac{f(L_0)\int_0^{L_0}(L_0^2-t^2)\cdot f(t)\,dt}{\left(\int_0^{L_0}f(t)\,dt\right)^2}
\le2L_0,
\]
as desired.
\end{proof}

\begin{proof}[Proof of \cref{lem:coordinate-moment-covariance}]
Let us rewrite the left-hand side as
\begin{align*}
&\int_{-1}^1\left(x_1^2-\E{x\sim g(x)/G}{x_1^2}\right)\cdot\frac{g_{\mathrm{sq}}(x_1)}{g_1(x_1)}\cdot\frac{g_1(x_1)}{G}\,dx_1\\
&=2\int_0^1\left(x_1^2-\E{x\sim g(x)/G}{x_1^2}\right)\cdot\frac{g_{\mathrm{sq}}(x_1)}{g_1(x_1)}\cdot\frac{g_1(x_1)}{G}\,dx_1\\
&=-2\int_0^1\left(\frac{d}{dx_1}\frac{g_{\mathrm{sq}}(x_1)}{g_1(x_1)}\right)\cdot
\left(\int_0^{x_1}\left(s^2-\E{x\sim g(x)/G}{x_1^2}\right)\cdot\frac{g_1(s)}{G}\,ds\right)\,dx_1.
\end{align*}
Since $t^{n-2}e^{-t^2/2}$ is log-concave and nondecreasing on
$t\in[0,(n-2)^{1/2}]$, \cref{lem:aux-logconcave-ratio} gives
\[
\frac{d}{dx_1}\frac{g_{\mathrm{sq}}(x_1)}{g_1(x_1)}
=-\Theta((n-2)^{1/2}-x_1)
=-\Theta(n^{1/2})
\]
for $x_1\in[0,1]$.
Since
\[
\int_0^{x_1}\left(s^2-\E{x\sim g(x)/G}{x_1^2}\right)\cdot\frac{g_1(s)}{G}\,ds\le0
\]
for all $x_1\in[0,1]$, it suffices to show
\[
\int_0^1\int_0^{x_1}\left(s^2-\E{x\sim g(x)/G}{x_1^2}\right)\cdot\frac{g_1(s)}{G}\,ds\,dx_1=-\Theta(1).
\]
We prove the upper and lower bounds separately.
\begin{enumerate}[label=(\roman*)]
\item The bound
\[
\left|\int_0^1\int_0^{x_1}\left(s^2-\E{x\sim g(x)/G}{x_1^2}\right)\cdot\frac{g_1(s)}{G}\,ds\,dx_1\right|=O(1)
\]
follows from the fact that $\int_0^1g_1(s)/G\,ds=1/2=O(1)$.

\item To show
\[
\int_0^1\int_0^{x_1}\left(s^2-\E{x\sim g(x)/G}{x_1^2}\right)\cdot\frac{g_1(s)}{G}\,ds\,dx_1=-\Omega(1),
\]
it suffices to show that there exist $d_1,d_2=\Omega(1)$ such that
\begin{equation}
\int_0^{x_1}\left(s^2-\E{x\sim g(x)/G}{x_1^2}\right)\cdot\frac{g_1(s)}{G}\,ds=-\Omega(1)
\label{eq:coordinate-moment-window}
\end{equation}
whenever $x_1\in[d_1,d_2]$.

We know that $g_1(s)/G=\Omega(1)$ for each
$s\in[0,1]$, due to the bound
$g_1(1)/g_1(0)>\Omega(1)$ we proved earlier.
Hence
\[
\E{x\sim g(x)/G}{x_1^2}
=2 \cdot\int_0^1s^2\cdot\frac{g_1(s)}{G}\,ds
=\Omega(1),
\]
so if we choose
\[
d_1:=\frac12\E{x\sim g(x)/G}{x_1^2}^{1/2},
\qquad
d_2:=\E{x\sim g(x)/G}{x_1^2}^{1/2}
\]
we will have $d_1,d_2=\Omega(1)$.

It remains to show Equation~\eqref{eq:coordinate-moment-window}.
We only need to check the case when $x_1=d_1$, i.e. we only need to verify
\[
\int_0^{d_1}\left(s^2-\E{x\sim g(x)/G}{x_1^2}\right)\cdot\frac{g_1(s)}{G}\,ds=-\Omega(1),
\]
which follows from
\begin{align*}
\int_0^{d_1}\left(s^2-\E{x\sim g(x)/G}{x_1^2}\right)\cdot\frac{g_1(s)}{G}\,ds
&\le\int_0^{d_1}\left(d_1^2-\E{x\sim g(x)/G}{x_1^2}\right)\cdot\frac{g_1(s)}{G}\,ds\\
&=-\frac34\E{x\sim g(x)/G}{x_1^2}\int_0^{d_1}\frac{g_1(s)}{G}\,ds\\
&=-\Omega(1)
\end{align*}
using $g_1(s)/G=\Omega(1)$ and
$\E{x\sim g(x)/G}{x_1^2}=\Omega(1)$, proved above.
\end{enumerate}
The two bounds together prove the lemma.
\end{proof}
% ===== end a2_counter_coordinate_moment.tex =====
% ===== begin a3_counter_ellipsoid_growth_asym.tex =====
\section{Asymmetric Counterexample to the Ellipsoid-Growth Conjecture}
\label{sec:counter_ellipsoid_growth_asym}

Here we provide an asymmetric counterexample to the Ellipsoid-Growth
Conjecture, which is Conjecture~3.5 of Jia--Laddha--Lee--Vempala
\cite{JLLV26:journal}, and stated in~\cref{conj:jllv-ellipsoid-growth}.
Specifically, to prove~\cref{thm:overview-counterexample-conj35}, we give an
explicit family of convex bodies $K \subset \mathbb{R}^n$ and ellipsoids
$E \subset \mathbb{R}^n$ for which
$K \cap E$ has mean zero and covariance
$I_n$, but the covariance trace of
$K \cap (1+\delta)E$ is at least
$c n^2\delta^2$.
Here $(1+\delta)E=\{(1+\delta)z:z\in E\}$ is dilation about the origin.
Our family works for every $n^{-1/2}\leq\delta\leq1$, including
$n^{-1/4}$.

The proof keeps track of the affine centering step, which does not commute with
dilation, and gives explicit estimates.
For a convex body $A$, all means and covariance matrices below refer to the
uniform probability measure on
$A$.
We use the standard origin-dilation notation
\[
sE := \{sz:z\in E\},
\qquad
s\geq 1.
\]
The distinction between origin dilation and dilation about the center of $E$
matters because our final isotropizing map is affine.

Conjecture~3.5 of \cite{JLLV26:journal} asserts that if $K\cap E$ is
$2$-isotropic, then
$K\cap 2E$ is well-rounded.
The covariance part of well-roundedness is
$\operatorname{tr}\operatorname{Cov}=\widetilde{O}(n)$.
We prove the following stronger negation, using mean-zero, covariance-$I_n$
isotropy rather than merely
$2$-isotropy.

The next theorem concerns the quantitative counterexample.

\begin{theorem}
\label{app:counterexample:theorem:quantitative}
There are absolute constants $c,C>0$ such that, for every sufficiently large
$n$, one can construct a convex body
$K\subset\mathbb{R}^n$ and an ellipsoid
$E\subset\mathbb{R}^n$ with the following properties.
\begin{enumerate}
\item \label{itm:quantitative:1} The center of $E$ lies in
  $\operatorname{int}K$.

\item \label{itm:quantitative:2} $K\cap E$ has mean zero and covariance
  $I_n$.

\item \label{itm:quantitative:3} Simultaneously for all $0<\delta\leq 1$
  satisfying
  $(n-1)\delta\geq C$,
\[
  \operatorname{tr}\operatorname{Cov}\bigl(K\cap(1+\delta)E\bigr)
  \geq c(n-1)^2\delta^2.
\]
\end{enumerate}
One may take, for example, $C=32$ and
$c=1/98304$.
\end{theorem}

\subsection{The raw body and its thin direction}
\label{app:counterexample:raw-body}

Write $n=d+1$ and points of
$\mathbb{R}^n$ as
$(x,y)\in\mathbb{R}\times\mathbb{R}^d$.
Define
\[
  f(x):=
  \begin{cases}
    1+x/2, & -2\leq x\leq 0,\\
    1,     & 0\leq x\leq 4,
  \end{cases}
\]
and
\begin{equation}
\label{app:counterexample:eq:raw-definitions}
\begin{aligned}
  K_0 &:= \{(x,y):-2\leq x\leq 4,\ \lVert y\rVert_2\leq f(x)\},\\
  E_0 &:= \{(x,y):(x+\tfrac12)^2+\lVert y\rVert_2^2\leq\tfrac54\}.
\end{aligned}
\end{equation}
The function $f$ is nonnegative and concave: its slope drops from
$1/2$ to
$0$ at the origin.
Consequently $K_0$ is convex.
Indeed, the triangle inequality and concavity of $f$ show that the segment
between any two points of
$K_0$ remains in
$K_0$.
The center $(-1/2,0)$ of
$E_0$ is strictly inside
$K_0$, since
$f(-1/2)=3/4$.

Let $L_0:=K_0\cap E_0$.
Let $g(x)$ be the radius of its
$x$-section.
The axial support of $E_0$ is
\[
\alpha
:=-\frac{1+5^{1/2}}{2}
\leq x
\leq \beta
:=\frac{5^{1/2}-1}{2},
\]
and the section radius of $E_0$ is
\[
  e_0(x)=(1-x-x^2)^{1/2}.
\]
For $x\leq 0$, direct subtraction gives
\begin{equation}
\label{app:counterexample:eq:radius-difference}
  e_0(x)^2-(1+x/2)^2=-x\left(2+\frac54x\right).
\end{equation}
It follows that the section function is exactly
\begin{equation}
\label{app:counterexample:eq:section-function}
  g(x)=
  \begin{cases}
    (1-x-x^2)^{1/2}, & \alpha\leq x\leq-8/5,\\
    1+x/2,          & -8/5\leq x\leq 0,\\
    (1-x-x^2)^{1/2}, & 0\leq x\leq\beta,\\
    0,              & \text{otherwise}.
  \end{cases}
\end{equation}
Thus there is a small secondary cap near $x=\alpha$, but its section radius is
at most
$1/5$.
More importantly, the unique section of radius $1$ is at
$x=0$, and the section radius decreases linearly to first order on both sides.

Let $X$ denote the first coordinate of a uniform point of
$L_0$.
After the common factor $\operatorname{vol}(B_2^d)$ is removed, its density is
\begin{equation}
\label{app:counterexample:eq:density}
  \frac{g(x)^d}{Z_d}, \qquad
  Z_d:=\int_{\mathbb{R}}g(x)^d\,dx.
\end{equation}

The next lemma concerns the inner moments.

\begin{lemma}
\label{app:counterexample:lemma:inner-moments}
For $d\geq 3$, there are absolute constants
$c_i,C_i>0$ such that
\[
\frac{c_0}{d}
\leq Z_d
\leq\frac{C_0}{d},
\qquad
\lvert\E{X\sim g^d/Z_d}{X}\rvert\leq\frac{C_1}{d},
\qquad
\frac{c_2}{d^2}
\leq\V{X\sim g^d/Z_d}{X}
\leq\frac{C_2}{d^2}.
\]
Moreover, by rotational symmetry,
\[
\operatorname{Cov}(L_0)=\operatorname{diag}(\sigma_x^2,\tau^2 I_d),
\qquad
\sigma_x^2=\Theta(d^{-2}),
\qquad
\tau^2=\Theta(d^{-1}).
\]
The explicit bounds $\sigma_x^2\leq 64/d^2$ and
$\sigma_x^2\geq 9/(2048d^2)$ are valid.
\end{lemma}

\begin{proof}
On the negative half-line, $g(x)\leq 1+x/2$, while for
$x\geq 0$,
\[
g(x)^d
=(1-x-x^2)^{d/2}
\leq e^{-dx/2}.
\]
Hence
\[
Z_d
\leq\int_{-2}^{0}(1+x/2)^d\,dx +\int_0^\infty e^{-dx/2}\,dx
\leq\frac4d.
\]
On $[-1/d,0]$, one has
$g(x)=1+x/2$ and
$g(x)^d\geq 1/2$, so
$Z_d\geq 1/(2d)$.

For $k=1,2$, the negative-side moments satisfy the exact beta-integral bound
\begin{equation}
\label{app:counterexample:eq:negative-moments}
\begin{aligned}
  \int_{-2}^{0}|x|^k g(x)^d\,dx
  &\leq 2^{k+1}\int_0^1(1-u)^k u^d\,du\\
  &=\frac{2^{k+1}k!}{(d+1)(d+2)\cdots(d+k+1)}.
\end{aligned}
\end{equation}
The positive-side moments obey
\[
  \int_0^\infty x^k e^{-dx/2}\,dx=k!\left(\frac2d\right)^{k+1}.
\]
Dividing by the lower bound for $Z_d$ gives
$\E{X\sim g^d/Z_d}{|X|}=O(d^{-1})$ and
$\E{X\sim g^d/Z_d}{X^2}\leq 64/d^2$, hence the desired upper variance bound.

For the lower bound, take
\[
I_-=[-2/d,-1/d],
\qquad
I_+=[1/(2d),1/d].
\]
For $d\geq 3$, the unnormalized density
$g^d$ is bounded below by
$1/4$ on both intervals.
Since $Z_d\leq 4/d$,
\[
\mathbb{P}(X\in I_-)\geq\frac1{16},
\qquad
\mathbb{P}(X\in I_+)\geq\frac1{32}.
\]
If $X'$ is an independent copy of
$X$, then
$\V{X\sim g^d/Z_d}{X}=\tfrac12\E{X,X'\sim g^d/Z_d}{(X-X')^2}$.
The two intervals are separated by at least $3/(2d)$, and therefore
\[
\V{X\sim g^d/Z_d}{X}
\geq \frac1{16}\frac1{32}\left(\frac3{2d}\right)^2
=\frac9{2048d^2}.
\]
Finally, conditional on $X=x$, the variable
$y$ is uniform in a
$d$-ball of radius
$g(x)$, so
\[
  \E{y\sim\operatorname{Unif}(g(x)B_2^d)}{y_i^2\mid X=x}=\frac{g(x)^2}{d+2}.
\]
The upper bound $\tau^2\leq 1/(d+2)$ is immediate.
A central interval of probability bounded below has $g(x)^2$ bounded below,
giving
$\tau^2\geq c/(d+2)$.
Rotational symmetry eliminates all mixed terms.
\end{proof}

Let
\[
m:=(\mu,0),
\qquad
\mu:=\E{X\sim g^d/Z_d}{X},
\qquad
\Sigma:=\operatorname{Cov}(L_0),
\]
and define the affine map
\begin{equation}
F(z):=\Sigma^{-1/2}(z-m),
\qquad
K:=F(K_0),
\qquad
E:=F(E_0).
\label{app:counterexample:eq:affine-map}
\end{equation}
Since $F$ is bijective,
\[
  K\cap E=F(L_0).
\]
The uniform law is carried to the uniform law under an invertible affine map, so
$K\cap E$ has mean zero and covariance
$I_n$ exactly.
Also, the center of $E$ belongs to
$\operatorname{int}K$, because
$(-1/2,0)\in\operatorname{int}K_0$.

\subsection{The enlarged ellipsoid creates a plateau}
\label{app:counterexample:enlarged-ellipsoid}

Fix $s=1+\delta\in(1,2]$.
This is the step at which the translation in $F$ must be handled carefully.
For every $z\in\mathbb{R}^n$,
\[
F^{-1}\bigl(sF(z)\bigr)
=m+s(z-m)
=sz+(1-s)m.
\]
Consequently
\begin{equation}
F^{-1}(sE)
=sE_0+(1-s)m
=:E_s,
\label{app:counterexample:eq:pullback-ellipsoid}
\end{equation}
and hence
\begin{equation}
K\cap sE=F(D_s),
\qquad
D_s:=K_0\cap E_s.
\label{app:counterexample:eq:enlarged-intersection}
\end{equation}
The set $E_s$ is the Euclidean ball with center
$(-s/2-\delta\mu,0)$ and radius
$s5^{1/2}/2$.
Its section radius is
\begin{equation}
\label{app:counterexample:eq:enlarged-section-radius}
  e_s(x)^2=\frac{5s^2}{4}-(x+s/2+\delta\mu)^2.
\end{equation}

On the cylindrical part $x\geq 0$ of
$K_0$, the complete unit
$d$-ball is present precisely while
$e_s(x)\geq 1$.
The positive endpoint of this interval is
\begin{equation}
x_*=-\frac{s}{2}-\delta\mu+A_s,
\qquad
A_s:=\left(\frac{5s^2}{4}-1\right)^{1/2}.
\label{app:counterexample:eq:plateau-endpoint}
\end{equation}
The other root is negative: since $A_s\geq s/2$ and
$\mu\geq\alpha$, it is at most
$-s-\delta\mu\leq-1+(-1-\alpha)\delta<0$.
To estimate $x_*$, define
\[
  q(s):=A_s-s/2.
\]
A direct differentiation gives
\[
\frac34
\leq q'(s)
\leq 2,
\qquad
1
\leq s
\leq 2,
\]
and $q(1)=0$.
Since the support of $X$ is
$[\alpha,\beta]$, its mean also lies in that interval.
Therefore
\begin{equation}
\frac{\delta}{8}
\leq x_*
=q(1+\delta)-\delta\mu
\leq 4\delta,
\qquad
0
<\delta
\leq 1.
\label{app:counterexample:eq:endpoint-bounds}
\end{equation}
Our choice of the upper endpoint $4$ in the definition of
$K_0$ guarantees that this whole interval lies in the cylindrical part.
We have thus found an exact radius-one cylindrical plateau
\begin{equation}
\label{app:counterexample:eq:plateau-set}
  \{(x,y):0\leq x\leq x_*,\ \lVert y\rVert_2\leq 1\}\subset D_s.
\end{equation}

The next lemma concerns the outer axial variance.

\begin{lemma}
\label{app:counterexample:lemma:outer-axial-variance}
If $d\delta\geq 32$, then the first-coordinate variance of the uniform law on
$D_s$ satisfies
\[
  \V{D_s}{x}\geq\frac{\delta^2}{1536}.
\]
\end{lemma}

\begin{proof}
After the common factor $\operatorname{vol}(B_2^d)$ is removed, the plateau has
mass exactly
$x_*$.
All mass with $x<0$ is at most
\[
  \int_{-2}^{0}(1+x/2)^d\,dx=\frac{2}{d+1}.
\]
For the positive tail, write $x=x_*+v$,
$v\geq 0$.
From Equation~\eqref{app:counterexample:eq:enlarged-section-radius}--
\eqref{app:counterexample:eq:plateau-endpoint},
\[
  e_s(x)^2=1-2A_sv-v^2.
\]
On the support of this tail, the $D_s$-section radius is at most $e_s(x)$.
Since $A_s\geq1/2$, its section volume, divided by
$\operatorname{vol}(B_2^d)$, is at most $e^{-dA_sv}$.
Hence the positive-tail mass is at most
\[
\int_0^\infty e^{-dA_sv}\,dv\leq\frac2d.
\]
Thus the total off-plateau mass is at most $4/d$.
By the lower endpoint bound $x_*\geq\delta/8$ in
Equation~\eqref{app:counterexample:eq:endpoint-bounds},
$d\delta\geq 32$ implies
$x_*\geq 4/d$, so the plateau has probability at least
$1/2$ under the uniform law on
$D_s$.

Conditioned on the plateau, the first coordinate is exactly uniform on
$[0,x_*]$.
Applying the law of total variance to the indicator of the plateau gives
\[
\V{D_s}{x}
\geq\frac12\V{x\sim\operatorname{Unif}[0,x_*]}{x}
=\frac{x_*^2}{24}
\geq\frac{\delta^2}{1536}.
\]
\end{proof}

\begin{proof}[Proof of \cref{app:counterexample:theorem:quantitative}]
The first two assertions, namely mean zero and covariance $I_n$ for
$K\cap E$, follow from the affine normalization in
Equation~\eqref{app:counterexample:eq:affine-map}.
Rotational symmetry gives
\[
  \Sigma=\operatorname{diag}(\sigma_x^2,\tau^2 I_d).
\]
By Equation~\eqref{app:counterexample:eq:enlarged-intersection}, the first
coordinate after applying $F$ is
$(x-\mu)/\sigma_x$.
Hence its variance under the uniform law on $K\cap sE$ is exactly
\[
  \frac{\V{D_s}{x}}{\sigma_x^2}.
\]
Using \cref{app:counterexample:lemma:inner-moments} and
\cref{app:counterexample:lemma:outer-axial-variance},
\[
\operatorname{tr}\operatorname{Cov}(K\cap sE)
\geq\frac{\delta^2/1536}{64/d^2}
=\frac{d^2\delta^2}{98304}.
\]
The construction of $K$ and
$E$ depends on
$n$ but not on
$\delta$, so the estimate holds simultaneously throughout the stated range.
\end{proof}

The next corollary concerns the scales of failure.

\begin{corollary}
\label{app:counterexample:corollary:scales-of-failure}
The choice $\delta=1$ gives
$\operatorname{tr}\operatorname{Cov}(K\cap2E)=\Omega(n^2)$ and disproves
Conjecture~3.5.
Moreover,
\[
\delta=n^{-1/4}
\quad
\Longrightarrow
\quad
\operatorname{tr}\operatorname{Cov}\bigl(K\cap(1+\delta)E\bigr) =\Omega(n^{3/2}).
\]
For every fixed $0<\varepsilon<1/2$, the smaller choice
$\delta=n^{-1/2+\varepsilon}$ gives trace
$\Omega(n^{1+2\varepsilon})$.
Finally,
\[
\delta(n)
=n^{-1/2}\exp\!\bigl((\log n)^{1/2}\bigr)
=n^{-1/2+o(1)}
\]
gives
\[
  \operatorname{tr}\operatorname{Cov}\bigl(K\cap(1+\delta)E\bigr)
  \geq n\exp\!\bigl(2(\log n)^{1/2}\bigr)/C,
\]
which is not $\widetilde{O}(n)$.
\end{corollary}
% ===== end a3_counter_ellipsoid_growth_asym.tex =====

\end{document}
